% chapters/ch11.tex —— 第 11 章 电化学催化研究方法
% 装配层：只负责按顺序引入各分片。分片缺失时跳过，便于增量成书。
% 分片文件：chapters/parts/ch11-{a..e}.tex
%
% 【2026-09-29 调序】原为 a→b→c→d，有两处结构问题：
%   (1) 分片 d = §11.5.1 阻抗谱基本原理（依源文本 p.855—861 写成），却被排在
%       分片 c（§11.5 后半与 §11.6 原位谱学）之后，于是成书里 §11.6 排在
%       §11.5.1 之前，而分片 c 开篇写「阻抗的基本定义……写在上文 §11.5.1」——
%       在成书里这是假陈述；
%   (2) \subsection 按出现顺序自动编号，§11.5.1 会拿到比 §11.5 大的号。
%   现改为 a→b→d→c→e：d（原理）紧接 b（§11.4）之后、c（§11.5 应用与 §11.6）
%   之前，指代方向与前后顺序恢复正确。**注意：编号与原书不再一一对应**——成书为 §11.5 阻抗谱基本原理、§11.6 电化学阻抗谱（原书 §11.5，仅其应用部分）、§11.7 电化学原位谱学（原书 §11.6）。要完全对齐须把分片 d 的标题降为 \subsubsection 并入分片 c 的 §11.5。。分片 d 不含 \chapterstart（章计数器
%   由分片 a 唯一推进），故调序不会影响章号。

\IfFileExists{chapters/parts/ch11-a.tex}{% ================================================================
%  chapters/parts/ch11-a.tex
%  第 11 章  电化学催化研究方法（分片 a）
%  源：S2（下册）第 11 章 11.1—11.2.3，书内 p.807—832 / PDF p.296—321
%  范围：11.1 电化学基本原理；11.2 控制电势的研究方法（至 11.2.3）
%  不含：11.3 及以后（控制电流、整体电解、流体动力学、阻抗、现场谱学）
% ================================================================

\chapterstart{电化学催化研究方法}%
{本章回答：怎样用电势和电流这两个可观测量，把电极反应的本征活性从传质、双电层、
溶液电阻等非本征因素中剥离出来。导读主线只有一条——\textbf{电极反应的总速率由传质与
电荷转移两步串联决定}，11.1 建立描述这两步的方程，11.2 的每一种控制电势方法，本质上
都是在实验上设法单独固定其中一步、观测另一步。全章按「基本原理（11.1）→ 控制电势的
研究方法（11.2）→ 控制电流、整体电解与流体动力学方法（11.3 起）」展开。}
\label{ch:11}

\subsection{电化学基本原理}

\subsubsection{电化学与催化过程关系}

电化学以化学反应中的电子转移以及外加电场中的化学反应为研究对象\src{S2 §11.1.1，
书内 p.809 / PDF p.297}。它与催化化学的交集就是\textbf{电催化}：既是用催化剂提升反应
速率的电化学过程，又是以电能为能量形式的催化过程。二者其实可以统一——电极本身就是
催化剂，改变电极种类或电极表面修饰状态，就能改变反应的速率、选择性和方向性；只有
这类「通过改变电极材料或表面状态来调控反应」的电极，才被称为电催化剂。

两个学科的关注点不同，落到数据分析上会得到不同的结论：催化视角看的是电催化剂如何
降低活化能垒、如何催化能量转化过程；电化学视角看的是界面电场、电解质组成、以及参与
反应的分子或离子被活化后反应速率如何提升。同一个活性数字，两个视角的解释并不相同，
这一点在读到「本征活性」类结论时必须追问研究者用的是哪一套语言。

\begin{probout}
电化学催化实验的可观测量只有两个：\textbf{流过工作电极的电流或电量}，以及\textbf{工作
电极相对参比电极的电势}。所有动力学参数都必须从这两个量的关系中反推。反推的前提是：
电流里不混入非法拉第过程，电势里不混入未补偿的溶液电阻降，传质速率已知或已被排除，
并且电势的参照点（参比电极）和面积基准（几何面积还是真实面积）都被明确交代过。
这四条前提缺任何一条，得到的速率常数和 Tafel 斜率都不可比。
\end{probout}

\subsubsection{电化学涉及的基本概念}

这组概念的用处是给后面的公式划定符号与判据的边界，其中最容易混淆的是\textbf{电势}
类的三个词：平衡电势、过电位、开路电势。

\paragraph{导体、半导体与绝缘体} 按电阻率 $\rho$ 划分：$\rho<\qty{1}{\milli\ohm\centi\meter}$ 为
导体（金、银、铜、铝、铁、电解液、石墨），$\rho>\qty{1}{\giga\ohm\centi\meter}$ 为绝缘体（云母、
橡胶、石英、陶瓷），介于两者之间为半导体（硅、锗、\ce{TiO2}、\ce{GaN}）
\src{书内 p.810 / PDF p.298}。对催化实验的直接含义是：电解液必须导电，但催化剂载体
（如氧化物、碳）可以处于半导体区，其导电性会随电势变化，这会带来额外的、与催化
无关的电流--电势非线性。

\paragraph{电解质与电解液} 电解质是在水溶液中或熔融态能导电的化合物，其导电来自荷电
溶剂化离子在电场下的定向迁移；电解质本身不一定导电，按是否完全电离分为强、弱电解质，
而溶解度大小与强弱无关。电解质并非都是液体：聚合物锂离子电池和质子交换膜燃料电池
使用的就是固体有机电解质。

\paragraph{电极} 按材料分为金属电极、氧化物电极、固体难溶盐电极等；按功能分为工作
电极、对电极、参比电极、离子选择电极。电极的作用是连接电解液与外电路，并在其表面或
近表面界面发生电化学反应。催化研究关心的是工作电极，但\textbf{没有对电极和参比电极的
明确交代，工作电极上的数据无法解释}。

\paragraph{活度、离子氛与电迁移} 溶液中的离子与其周围离子形成相对规律的离子排布区，
称为离子氛。离子到达电极表面发生反应前需从离子氛中剥离，要耗费能量，因此溶液中离子的
自由能与反应活性都略低于自由离子。为描述实际溶液与理想溶液的偏离引入\textbf{活度}：
它是离子相对理想溶液浓度的有效值。离子氛随浓度增大而变大，故活度随浓度增大而减小；
无限稀溶液中离子可近似为自由离子。电解质本体中离子在电场下的定向移动称为\textbf{电迁移}；
电极表面液层中则因部分离子持续反应造成浓度差，驱动离子\textbf{扩散}。扩散和电迁移
共同完成电活性物质的传递，对应电流也分为扩散电流与迁移电流。电极近表面电活性物质的
流量决定反应速率，从而控制外电路电流。

\paragraph{电势的三种含义} 这三个词在催化文献里最常被混用，必须分清：

\begin{itemize}
\item \textbf{平衡电势} $E_{\mathrm{eq}}$：给定温度、压力和\textbf{全部组分活度}下，
  电极上氧化与还原两个方向速率相等、净电流为零时的电势。它是一个由热力学确定的量
  （Nernst 方程给出），与电极材料无关，只与电对和活度有关。
\item \textbf{过电位（超电势）} $\eta=E-E_{\mathrm{eq}}$：实际电极电势偏离平衡电势的
  量，是驱动净反应的「额外推动力」，符号与大小同时编码了动力学信息。$\eta$ 只对确定的
  $E_{\mathrm{eq}}$ 才有定义；讨论 $\eta$ 前必须先说明用什么电对、什么活度算出的
  $E_{\mathrm{eq}}$。
\item \textbf{开路电势} OCP：外电路断开（无净电流）时测得的工作电极电势。它\textbf{等于}
  平衡电势只在电极表面处于热力学平衡时成立；催化研究中电极表面常被吸附物、氧化物层
  或混合电位占据，此条件下 OCP 是一个\textbf{混合电位}，与任何单一电对的
  $E_{\mathrm{eq}}$ 都不相等，因而不能用来计算 $\eta$。
\end{itemize}

\begin{pitfall}
把 OCP 当平衡电势、直接用它算过电位，是电催化文献里最高频的错误之一。判据很简单：
若在该 OCP 处体系的净电流为零是因为\textbf{正逆两个方向的同一反应}速率相等，OCP 才是
平衡电势；若净电流为零是\textbf{两个不同反应}（如氧还原与铂氧化）各自贡献正负电流相消，
那它是混合电位。区分办法是在 OCP 附近做微扰，看响应是容性还是可通过极化曲线外推。
\end{pitfall}

\subsubsection{电势和电池热力学}

电势差推动电子在电路中定向移动；恒温恒压下，若系统的非膨胀功只有电功，则

\begin{equation}
(\Delta_{\mathrm{r}}G)_{T,p}=-nEF
\label{eq:gibbs-emf}
\end{equation}

式中 $n$ 为反应输出电荷的物质的量，$E$ 为可逆电池的电动势，$F$ 为法拉第常量
\src{式 11-3，书内 p.812 / PDF p.300}。这是联系热力学与电化学的主要桥梁：测量电势
（可观测）即可解决催化过程中的热力学问题。

\begin{keycond}
\equref{eq:gibbs-emf} 的成立条件是\textbf{可逆电池、恒温恒压、非膨胀功只有电功}。
量纲检查：$[-nEF]=1\times\mathrm{C.mol^{-1}}\times\mathrm{V}
=\mathrm{J.mol^{-1}}$，与 $\Delta_{\mathrm{r}}G$ 一致。
\end{keycond}

\paragraph{可逆电池} 化学上可逆指仅改变电流方向即可让反应反向，不产生新反应。热力学
可逆则进一步要求充放电电流无限小、电池在平衡态附近工作。实际过程总以一定速率进行，
不具有严格的热力学可逆性；但外接足够大的电阻使反应足够慢、且反应程度足够小（各组分
活度基本不变）时，在要求精度内热力学公式仍然适用，此时电池对外做最大电功 $|\Delta G|$。

由 $\Delta G$ 与温度的关系可得电池反应的熵变与温度系数：

\begin{equation}
\Delta S=-\left(\frac{\partial \Delta G}{\partial T}\right)_{p},
\qquad
\Delta S=nF\left(\frac{\partial E}{\partial T}\right)_{p}
\label{eq:entropy-emf}
\end{equation}

\src{式 11-5、11-6，书内 p.813 / PDF p.301}。这样无法直接观测的热力学量就转化为电势
等可观测量——其中一个实用推论是：$\partial E/\partial T$ 给出熵变，因此\textbf{参比
电极的温度系数不被补偿时会直接污染 $\Delta S$ 的数值}。

\paragraph{半反应、参比电极与标准电极电势} 单个电极的绝对电势无法精确测量，因此用
标准参比电极建立统一标度。按惯例采用常规氢电极（NHE，亦称标准氢电极 SHE）：

\begin{equation}
\ce{Pt | H2}(p^\ominus)\,|\,\ce{H+}(a=1)
\label{eq:she}
\end{equation}

其电势在所有温度下被定义为零。电极与 NHE 构成可逆电池，测出的电池电势即为该电极的
半电池电势。例如 $E(\ce{Cl-}|\ce{AgCl}|\ce{Ag})=+\qty{0.22}{V}$（对 NHE），
$E(\ce{Zn^{2+}}|\ce{Zn})=-\qty{0.76}{V}$，故 $\ce{Zn|AgCl}$ 电池的
$E=0.22-(-0.76)=\qty{0.98}{V}$ \src{式 11-7、11-8，书内 p.813 / PDF p.301}。
通常把半反应写成还原形式，所得电极电势即相对 NHE 的还原电势。

\begin{pitfall}
「电势是 $-0.2\ \mathrm{V}$」这句话在电催化里是不完整的信息。一个可用的电势数字至少
需要四项交代：\textbf{参比电极种类}（NHE/SHE、Ag/AgCl 及其 \ce{KCl} 浓度、饱和甘汞
SCE、\ce{Hg/HgO}、\ce{Hg/Hg2SO4} 等）、\textbf{温度}、\textbf{是否做过 iR 补偿及其
补偿比例}、\textbf{换算到 RHE 或 NHE 时用的 pH}。参比电极还会随时间漂移、被 \ce{Cl-}
污染、液接界电势随溶液组成改变；长期实验必须定期用标准液标定。缺任何一项的 Tafel
斜率都不能与文献数值直接比较。
\end{pitfall}

\begin{principle}
整章的方法都服务于同一个速率模型。电极反应的总速率（即电流密度 $j$）由两个\textbf{串联}
的环节决定：\textbf{（i）反应物从本体输运到电极表面}，速率由传质系数 $m$ 表征；
\textbf{（ii）反应物在表面完成电荷转移}，速率由动力学参数 $k^0$、$\alpha$（等价地
$j_0$、$\alpha$）表征。串联意味着

\begin{equation}
\frac{1}{j_{\text{总}}}=\frac{1}{j_{\text{传质}}}+\frac{1}{j_{\text{电荷转移}}}
\label{eq:series-rate}
\end{equation}

因此实测电流永远只报告\textbf{较慢一步}的信息：传质快于电荷转移时电流反映本征动力学
（此时才能读 Tafel 斜率与 $j_0$）；电荷转移快于传质时电流被钉在扩散极限上，反映的是
$D$、$\delta$ 与 $C^*$，与本征活性无关。\equref{eq:series-rate} 的成立条件是两步彼此
独立、无耦合的中间步骤，且电极表面不积累中间体；量纲检查：各项均为
$\mathrm{A.cm^{-2}}$，成立。本章后面每一种方法，本质上都是在实验上设法\textbf{固定其中
一步、观测另一步}：传质系数由扩散层厚度与扩散系数按方法设定，电荷转移速率由电势按
Butler--Volmer 关系设定。
\end{principle}

\paragraph{Nernst 方程} 当各组分不处于标准状态时，对总反应
$\upsilon_{\mathrm{O}}\mathrm{O}+\upsilon_{\mathrm{H_2}}\mathrm{H_2}
\rightleftharpoons \upsilon_{\mathrm{R}}\mathrm{R}+\upsilon_{\mathrm{H^+}}\mathrm{H^+}$，
由化学反应等温式与 \equref{eq:gibbs-emf} 得

\begin{equation}
E=E^\ominus-\frac{RT}{nF}\ln\frac{a_{\mathrm{R}}^{\upsilon_{\mathrm{R}}}
a_{\mathrm{H^+}}^{\upsilon_{\mathrm{H^+}}}}
{a_{\mathrm{O}}^{\upsilon_{\mathrm{O}}}a_{\mathrm{H_2}}^{\upsilon_{\mathrm{H_2}}}}
\label{eq:nernst}
\end{equation}

\src{式 11-11，书内 p.814 / PDF p.302}。式中 $E^\ominus$ 为所有参加反应的组分都处于
标准状态时的电动势；活度约定为：纯液体或固态纯物质 $a=1$；气态物质
$a_i=\gamma_i(p_i/p^\ominus)$；溶质 $a_i=\gamma_i(C_i/C^\ominus)$，$C^\ominus$ 通常取
\qty{1}{mol.L^{-1}}。书中给出的常数为 $R=\qty{8.31441}{J.mol^{-1}.K^{-1}}$、
$F=\qty{96486.7}{C.mol^{-1}}$ \src{书内 p.811—812 / PDF p.299—300}。

\begin{keycond}
\equref{eq:nernst} 要求\textbf{电化学平衡、可逆条件、各组分活度已知且界面处与本体一致}。
量纲检查：$[RT/(nF)]=\mathrm{J.mol^{-1}}/\mathrm{C.mol^{-1}}=\mathrm{V}$，与 $E$ 一致。
两个常见误用：把本体浓度直接代入活度（高离子强度下 $a$ 可比 $C$ 低很多）；在有大电流
通过时用电势反解表面活度（此时表面活度已被传质改写，需换用 \equref{eq:eta-conc}）。
\tocheck{未对照原页面图像：书中 $F$ 取 96486.7 C/mol，是较老取值（现用 96485.33 C/mol），
本分片按原书照录；若全书统一约定用新值需在此处改写。}
\end{keycond}

Nernst 方程也是离子选择性电极的本质思想：若保持 $\alpha$ 相中物质 $i$ 的活度不变，
测得膜电势 $E_{\mathrm{m}}=-\dfrac{RT}{nF}\ln\dfrac{a_i^\alpha}{a_i^\beta}$ 即可得到
$\beta$ 相中 $i$ 的活度，常用 pH 计的玻璃电极即属此类
\src{式 11-12，书内 p.814 / PDF p.302}。

\subsubsection{电极反应动力学}

可逆电极电势是在电流无限小时得到的。有电流通过时电极电势偏离可逆值，这一现象称为
\textbf{电极极化}。催化关心的正是偏离的大小与方式。

\paragraph{电极电流} 电极过程同时进行氧化与还原两个相反过程，二者速率相同时电极处于
平衡态、外电路无电流。对

\begin{equation}
\ce{O + $n$e- <=>[k_{\mathrm{f}}][k_{\mathrm{b}}] R}
\label{eq:redox}
\end{equation}

还原与氧化速率按 Arrhenius 形式写出，正比于各自反应物的表面浓度，比例常数为速率常数
\src{式 11-13—11-16，书内 p.815 / PDF p.303}：

\begin{align}
i_{\mathrm{f}}&=nFAk_{\mathrm{f}}C_{\mathrm{O}},\qquad
i_{\mathrm{b}}=nFAk_{\mathrm{b}}C_{\mathrm{R}}\notag\\[2pt]
i&=i_{\mathrm{f}}-i_{\mathrm{b}}
=nFA\left(k_{\mathrm{f}}C_{\mathrm{O}}-k_{\mathrm{b}}C_{\mathrm{R}}\right)
\label{eq:current-basic}
\end{align}

式中 $A$ 为电极面积，$C_{\mathrm{O}}$、$C_{\mathrm{R}}$ 为\textbf{电极表面}浓度。
量纲检查：$[nFAkC]=\mathrm{C.mol^{-1}}\times\mathrm{cm^2}\times\mathrm{cm.s^{-1}}
\times\mathrm{mol.cm^{-3}}=\mathrm{C.s^{-1}}=\mathrm{A}$，自洽。
\equref{eq:current-basic} 把电极电流与活化能联系起来——而由热力学部分已知电极电势
直接决定自由能，因此改变电势就改变了活化能。

\paragraph{Butler--Volmer 模型} 电势变化改变了电子的相对能量，因而改变氧化与还原的
活化自由能垒。\textbf{先把参考电势说清楚}：标准速率常数 $k^0$ 定义在标准（形式）电势
$E^{0'}$ 处，因此下面先用 $E-E^{0'}$ 写标准速率常数形式，再用 Nernst 关系换算成以过电位
$\eta=E-E_{\mathrm{eq}}$ 表示的形式。设电势正移到 $E$，电子能量下降
$F(E-E^{0'})$，则

\begin{equation}
\Delta G_{\mathrm{a}}^\ddagger=\Delta G_{\mathrm{a}}^{0\ddagger}-(1-\alpha)F(E-E^{0'}),
\qquad
\Delta G_{\mathrm{c}}^\ddagger=\Delta G_{\mathrm{c}}^{0\ddagger}+\alpha F(E-E^{0'})
\label{eq:dg-eta}
\end{equation}

\src{式 11-17，书内 p.816 / PDF p.304}。式中 $\alpha$ 为\textbf{传递系数}，取值与
活化能垒交叉区域形状有关，范围在 0～1 之间；$\eta=E-E_{\mathrm{eq}}$ 为过电位。
其物理含义是：电势变化对阳极和阴极活化能垒的削弱/增强\textbf{按 $\alpha$ 与
$1-\alpha$ 分配}，$\alpha$ 因此直接决定 Tafel 斜率。

把 \equref{eq:dg-eta} 代回 Arrhenius 式，并以标准电势处
（$E=E^{0'}$，该处 $C_{\mathrm{O}}=C_{\mathrm{R}}$、$k_{\mathrm{f}}=k_{\mathrm{b}}=k^0$，
$k^0$ 为标准速率常数）为基准归一化，得

\begin{equation}
k_{\mathrm{f}}=k^0\exp\!\left(-\frac{\alpha F(E-E^{0'})}{RT}\right),
\qquad
k_{\mathrm{b}}=k^0\exp\!\left(\frac{(1-\alpha)F(E-E^{0'})}{RT}\right)
\label{eq:k-eta}
\end{equation}

\src{式 11-20、11-21，书内 p.816 / PDF p.304}。代入 \equref{eq:current-basic} 得
电流--电势特征关系式

\begin{equation}
i=FAk^0\left[C_{\mathrm{O}}\exp\!\left(-\frac{\alpha F(E-E^{0'})}{RT}\right)
-C_{\mathrm{R}}\exp\!\left(\frac{(1-\alpha)F(E-E^{0'})}{RT}\right)\right]
\label{eq:bv-first}
\end{equation}

\src{式 11-22，书内 p.816 / PDF p.304}。这是 Butler--Volmer 动力学表达式的一种形式，
是处理电极电流与电势关系的基本手段。\textbf{参考点不能混用}：式~\eqref{eq:bv-first}
用 $E-E^{0'}$ 才与 $k^0$ 匹配；若把它写成 $E-E_{\mathrm{eq}}$，则在
$C_{\mathrm{O}}^*\neq C_{\mathrm{R}}^*$ 时 $\eta=0$（平衡态）会给出非零电流。
\tocheck{未对照原页面图像：原书式（11-17）把 $E-E^{0}$ 与 $\eta$ 写在同一组式子里
（正文只写「采用平衡电势作为电势参考点」，OCR 中该参考电势的符号不可辨），而式
（11-20）—（11-24）使用的 $k^0$ 定义在标准电势处，二者只在
$C_{\mathrm{O}}^*=C_{\mathrm{R}}^*$ 时相等；本稿按可自洽的 $E-E^{0'}$ 写出标准速率常数
形式，下文再用 Nernst 关系换算为 $\eta$ 与 $i_0$ 的形式。原书该参考电势的写法未对照
原页面图像。}

\paragraph{交换电流密度与传递系数} 平衡时净电流为零、表面浓度等于本体浓度
$C_{\mathrm{O}}^*$、$C_{\mathrm{R}}^*$。把 $E=E_{\mathrm{eq}}$ 代入
\equref{eq:bv-first}，括号为零即得本体浓度下的 Nernst 关系
$\exp\!\left[F(E_{\mathrm{eq}}-E^{0'})/RT\right]=C_{\mathrm{O}}^*/C_{\mathrm{R}}^*$
（单电子，即原书式 11-23）。此时正逆两个方向的电流\textbf{大小相等方向
相反}，其值定义为交换电流 $i_0$：

\begin{equation}
i_0=FAk^0C_{\mathrm{O}}^{*\,(1-\alpha)}C_{\mathrm{R}}^{*\,\alpha},
\qquad
j_0=\frac{i_0}{A}
\label{eq:i0}
\end{equation}

\src{式 11-24、11-25，书内 p.817 / PDF p.305}。$j_0$ 的单位是 $\mathrm{A.cm^{-2}}$，
量纲检查：$[FAk^0C]=(\mathrm{C.mol^{-1}})\mathrm{cm^2}(\mathrm{cm.s^{-1}})
(\mathrm{mol.cm^{-3}})=\mathrm{A}$，成立。$i_0$ 表征催化反应过程中活化能垒的大小：
$j_0$ 越大，同一过电位下净电流越大，\textbf{在同样的传质条件下}表现出的活性越高。
把上述 Nernst 关系写成
$E-E^{0'}=\eta+\dfrac{RT}{F}\ln\dfrac{C_{\mathrm{O}}^*}{C_{\mathrm{R}}^*}$ 后代入
\equref{eq:bv-first}，并用 \equref{eq:i0} 消去 $FAk^0$，即得一般浓度比下
\textbf{代数等价}的常用形式：

\begin{equation}
i=i_0\left[\frac{C_{\mathrm{O}}}{C_{\mathrm{O}}^*}
\exp\!\left(-\frac{\alpha F\eta}{RT}\right)
-\frac{C_{\mathrm{R}}}{C_{\mathrm{R}}^*}
\exp\!\left(\frac{(1-\alpha)F\eta}{RT}\right)\right]
\label{eq:bv}
\end{equation}

\src{式 11-26，书内 p.817 / PDF p.305}。这一步是\textbf{经 Nernst 关系代入化简}，不是
简单「把 $FAk^0$ 换成 $i_0$」：$i_0$ 中含 $C_{\mathrm{O}}^{*\,(1-\alpha)}
C_{\mathrm{R}}^{*\,\alpha}$，只有代入 Nernst 关系才能把参考电势 $E^{0'}$ 换成
$\eta=E-E_{\mathrm{eq}}$ 而保持等价；当 $C_{\mathrm{O}}^*=C_{\mathrm{R}}^*$ 时
$E_{\mathrm{eq}}=E^{0'}$，两式逐项相同。自检：$\eta=0$ 时 \equref{eq:bv} 的括号恒为
零，故任意本体浓度比下都有 $i=0$。

\begin{keycond}
\equref{eq:bv} 的成立条件有三条：\textbf{（i）单步骤、单电子的基元反应}——本段式
\equref{eq:dg-eta}—\equref{eq:bv-first} 按单电子写出（$F$ 不乘 $n$，与原书式
11-17—11-30 一致，式~\eqref{eq:current-basic} 中的 $n$ 在此取 $1$）；若推广到 $n$ 电子
一步基元反应，各式中的 $F$ 须\textbf{同步}换成 $nF$，包括 Tafel 斜率
$RT/(\alpha nF)$（后文传质部分的式~\eqref{eq:diff-layer} 起即按 $nF$ 写出）；
\textbf{（ii）$k^0$ 定义在 $E=E^{0'}$ 处（该处 $C_{\mathrm{O}}=C_{\mathrm{R}}$），
$i_0$ 定义在平衡态 $\eta=0$（任意本体浓度比，由 Nernst 关系确定 $E_{\mathrm{eq}}$）}；
\textbf{（iii）浓度项取电极表面浓度 $C_{\mathrm{O}}$、$C_{\mathrm{R}}$，不是本体浓度}。
量纲检查：$[i]=[i_0]=\mathrm{A}$（或两侧同除 $A$
写成电流密度）。第三条最容易被忽略：一旦把表面浓度换成本体浓度，方程就默认了
「传质无穷快」，此时得到的 $j_0$ 是表观值而非本征值。
\end{keycond}

\paragraph{Butler--Volmer 方程的两个极限} 不考虑传质影响时：
\textbf{（1）小过电位}，由泰勒展开 $e^x\approx1+x$，\equref{eq:bv} 化为

\begin{equation}
i=-i_0\frac{F\eta}{RT}
\qquad\text{或}\qquad
j=-\frac{F}{RT}j_0\,\eta
\label{eq:bv-linear}
\end{equation}

电流与过电位呈线性关系，斜率即\textbf{电荷转移电阻}的倒数。这是电化学阻抗谱测 $j_0$ 的
理论基础，也是「低过电位区不能读 Tafel 斜率」的原因。\textbf{（2）大过电位}，括号中
一项可忽略，以较大负过电位（阴极）为例：

\begin{equation}
i=i_0\exp\!\left(-\frac{\alpha F\eta}{RT}\right)
\label{eq:bv-cathodic}
\end{equation}

重排为

\begin{align}
\eta&=\frac{RT}{\alpha F}\ln i_0-\frac{RT}{\alpha F}\ln i
\label{eq:tafel}\\[2pt]
\eta&=a+b\ln i,\qquad
a=\frac{RT}{\alpha F}\ln i_0,\qquad b=-\frac{RT}{\alpha F}\notag
\end{align}

\src{式 11-27—11-30，书内 p.817—818 / PDF p.305—306}。\equref{eq:tafel} 即 Tafel 经验
公式 $\eta=a+b\ln i$ 的具体形式；若改用常用对数（文献多如此），则
$b=-2.303RT/(\alpha F)$，\qty{25}{\celsius} 下 $|b|=\qty{59}{\mV.dec^{-1}}/\alpha$。

\begin{keycond}
\equref{eq:tafel} 的成立条件是\textbf{高过电位极限}（通常 $|\eta|>\qty{50}{mV}$ 以上，
视体系而定）且\textbf{传质不构成限制}（$i\ll i_{\mathrm{t}}$，一般取 $i<0.1i_{\mathrm{t}}$）。
量纲检查：$[RT/(\alpha F)]=\mathrm{V}$，$[\ln i]$ 无量纲，故 $b$ 为 V（或 mV/dec）。
三条推论必须记住：$\alpha$ 只由斜率给出，$j_0$ 必须由\textbf{截距}给出，故
\textbf{斜率错一点、截距就错一个数量级}；$\alpha$ 与 $j_0$ 通过 \equref{eq:i0} 耦合，
只有当 $C_{\mathrm{O}}^*=C_{\mathrm{R}}^*$ 时 $j_0=FAk^0C^*$ 才是干净的本征量；
由阴极支和阳极支分别外推得到的截距应给出同一个 $j_0$，不重合说明有附加过程。
\tocheck{未对照原页面图像：书中式（11-30）写作 $b=-RT/(\alpha F)$（自然对数形式），
本分片按原书保留自然对数并补注了常用对数的换算；若本章他处统一采用 $\lg$，需改符号。}
\end{keycond}

\begin{figure}[tbp]
\centering
\begin{tikzpicture}
\begin{axis}[
  width=\dimexpr\twocolwidth-2mm\relax, height=58mm,
  xlabel={过电位 $\eta$ / V}, ylabel={$\ln(|j|\,/\,\mathrm{A.cm^{-2}})$},
  xmin=-0.30, xmax=0.30, ymin=-10, ymax=1,
  xtick={-0.3,-0.2,-0.1,0,0.1,0.2,0.3},
  ytick={-8,-6,-4,-2,0},
  legend style={at={(0.02,0.98)},anchor=north west,legend columns=1},
  clip=false
]
% alpha = 0.5, f = F/RT ≈ 38.92 V^-1 at 298 K ; log10 scale factor ln -> use natural log
% ln|j| = ln j0 - alpha*f*eta   (cathodic, eta<0 gives increase), anodic mirrored
\addplot[mainline,black,domain=-0.30:0,samples=80]{ln(1e-6) - 0.5*38.92*x};
\addplot[mainline,black,domain=0:0.30,samples=80]{ln(1e-6) + 0.5*38.92*x};
\addlegendentry{$\alpha=0.5$, $j_0=\qty{e-6}{\A\per\square\cm}$}
\addplot[mainline,black,dashed,domain=-0.30:0,samples=80]{ln(1e-8) - 0.5*38.92*x};
\addplot[mainline,black,dashed,domain=0:0.30,samples=80]{ln(1e-8) + 0.5*38.92*x};
\addlegendentry{$\alpha=0.5$, $j_0=\qty{e-8}{\A\per\square\cm}$}
\addplot[mainline,black,dashdotted,domain=-0.30:0,samples=80]{ln(1e-6) - 0.25*38.92*x};
\addplot[mainline,black,dashdotted,domain=0:0.30,samples=80]{ln(1e-6) + 0.75*38.92*x};
\addlegendentry{$\alpha=0.25$, $j_0=\qty{e-6}{\A\per\square\cm}$}
\addplot[thinline,black,dotted,domain=-0.30:0.30,samples=2]{ln(1e-6)};
\node[anchor=west,font={\fontsize{8}{9.5}\selectfont}] at (axis cs:0.02,-9.4)
  {实线/虚线：斜率相同、截距不同 $\Rightarrow$ 只差 $j_0$};
\end{axis}
\end{tikzpicture}
\caption{Butler--Volmer 方程在高过电位极限下的 Tafel 图（模型计算，$\qty{298}{K}$）。
读者应看两件事：\textbf{同一斜率、不同截距}只反映 $j_0$ 的差别（实线对虚线，差两个
数量级）；\textbf{斜率改变}反映传递系数 $\alpha$ 的差别（点划线，阴极支 $\alpha=0.25$，
$\alpha$ 越小斜率越陡）。值得注意的局限：真实体系的 Tafel 图常在中段出现折点，
说明速率控制步骤随过电位切换，此时用一条直线拟合全区间会同时给出错误的 $\alpha$ 与
$j_0$；低过电位区（$|\eta|<\qty{50}{mV}$）应改用 \equref{eq:bv-linear} 的线性关系。
\modelmark}
\label{fig:tafel}
\end{figure}

\subsubsection{静态电极过程中的物质传递影响}

前面所有含有电极表面浓度 $C_{\mathrm{O}}$ 的式子都留了一个未定的量。除平衡态等特殊
情形外，本体浓度 $C^*$ 与表面浓度 $C$ 并不相等，因此\textbf{表面浓度信息是理解浓差
极化行为的关键}。溶液中物质传递有迁移、扩散、对流三种形式；对\textbf{静态电极}，
对流可忽略；若溶液中有\textbf{过量支持电解质}，电迁移也可不考虑。此时电极周围传质由
扩散控制，符合 Fick 扩散定律，速率与浓度梯度成正比：

\begin{equation}
N_{\mathrm{O}}=-D_{\mathrm{O}}\left(\frac{\partial C_{\mathrm{O}}}{\partial x}\right)
\label{eq:fick1}
\end{equation}

\src{式 11-31，书内 p.818 / PDF p.306}。式中 $N_{\mathrm{O}}$ 为距电极表面 $x$ 处物质
$\mathrm{O}$ 的流量（$\mathrm{mol.s^{-1}.cm^{-2}}$），$D_{\mathrm{O}}$ 为扩散系数
（$\mathrm{cm^2.s^{-1}}$）。\textbf{符号约定（本节与式~\eqref{eq:current-basic} 统一）}：
坐标 $x$ 由电极表面指向溶液本体，$N_{\mathrm{O}}$ 为沿 $+x$ 方向的流量；电流取
\textbf{「还原为正」}，即 $i=i_{\mathrm{f}}-i_{\mathrm{b}}$，故
$i=-nFA\,N_{\mathrm{O}}|_{x\to0}$——还原时物质向 $-x$（朝电极）输运，$N_{\mathrm{O}}<0$，
电流为正。

\paragraph{Nernst 扩散层与极限电流} 假设电极表面存在厚度为 $\delta_{\mathrm{O}}$ 的
静止层，$x=\delta_{\mathrm{O}}$ 以外浓度均为本体浓度 $C_{\mathrm{O}}^*$，传递到电极表面
的物质被完全消耗。静止层内浓度线性上升，
$\partial C_{\mathrm{O}}/\partial x\approx
(C_{\mathrm{O}}^*-C_{\mathrm{O}})/\delta_{\mathrm{O}}>0$，由式~\eqref{eq:fick1} 与上述
$i=-nFA\,N_{\mathrm{O}}$（Faraday 定律）得

\begin{equation}
\frac{i}{nFA}=D_{\mathrm{O}}
\frac{C_{\mathrm{O}}^*-C_{\mathrm{O}}}{\delta_{\mathrm{O}}}
\label{eq:diff-layer}
\end{equation}

当表面浓度 $C_{\mathrm{O}}=0$ 时传递速率达到最大，电流达到平台值，称为\textbf{极限电流}
$i_{\mathrm{t}}$：

\begin{equation}
\frac{i_{\mathrm{t}}}{nFA}=D_{\mathrm{O}}
\frac{C_{\mathrm{O}}^*}{\delta_{\mathrm{O}}}
\label{eq:il}
\end{equation}

\src{式 11-32、11-33，书内 p.819 / PDF p.307}。还原时 $C_{\mathrm{O}}^*>C_{\mathrm{O}}$，
故 $i>0$、$i_{\mathrm{t}}>0$，与动力学式~\eqref{eq:current-basic}（$i_{\mathrm{f}}>0$ 为
还原）以及后面的 Cottrell 式~\eqref{eq:cottrell} 同号，同时也与
\equref{eq:tafel} 成立条件中的 $i\ll i_{\mathrm{t}}$ 自洽。若讨论的是反向的氧化过程，
把浓度差换成 $C_{\mathrm{R}}-C_{\mathrm{R}}^*$，电流随之变号。$i_{\mathrm{t}}$ 只表征
电极表面的传质情况，\textbf{不含任何动力学信息}——这一点在后面每一种方法里都会再出现。
量纲检查：$[nFAD C^*/\delta]=(\mathrm{C.mol^{-1}})\mathrm{cm^2}
(\mathrm{cm^2.s^{-1}})(\mathrm{mol.cm^{-3}})/\mathrm{cm}=\mathrm{C.s^{-1}}=\mathrm{A}$。
\tocheck{未对照原页面图像：原书式（11-32）（11-33）写作
$i/nFA=-D_{\mathrm{O}}(C_{\mathrm{O}}^*-C_{\mathrm{O}})/\delta_{\mathrm{O}}$
（OCR 中负号可辨），即扩散部分取「还原为负」的取向，与式（11-16）的「还原为正」
相反；本稿统一为还原为正，并写明 $x$ 的取向与 $i=-nFA N_{\mathrm{O}}$。其后式
（11-36）—（11-40）只含比值 $i_{\mathrm{t}}/(i_{\mathrm{t}}-i)$，在两套取向下的数值完全
相同，故未改动；原页面图像未对照。}

\paragraph{浓差过电位的分解} 表面浓度对电极电势的影响可拆成热力学与动力学两部分。
由 Nernst 方程，本体浓度与表面浓度对应的电势差为
$\eta_1=\dfrac{RT}{nF}\ln\dfrac{C_{\mathrm{O}}}{C_{\mathrm{O}}^*}$
\src{式 11-34、11-35，书内 p.819 / PDF p.307}；把 \equref{eq:diff-layer} 与
\equref{eq:il} 代入，得

\begin{equation}
\eta_1=-\frac{RT}{nF}\ln\!\left(\frac{i_{\mathrm{t}}}{i_{\mathrm{t}}-i}\right)
\label{eq:eta1}
\end{equation}

再按 Butler--Volmer 方程在高电流密度区略去第二项，得到
$i=i_0\dfrac{C_{\mathrm{O}}}{C_{\mathrm{O}}^*}
\exp\!\left(-\dfrac{\alpha nF\eta}{RT}\right)$ 及其重排
$\eta=-\dfrac{RT}{\alpha nF}\left(\ln\dfrac{i}{i_0}
+\ln\dfrac{C_{\mathrm{O}}^*}{C_{\mathrm{O}}}\right)$
\src{式 11-37、11-38，书内 p.819 / PDF p.307}。只保留浓度项即

\begin{equation}
\eta_2=-\frac{RT}{\alpha nF}\ln\frac{C_{\mathrm{O}}^*}{C_{\mathrm{O}}}
=-\frac{RT}{\alpha nF}\ln\!\left(\frac{i_{\mathrm{t}}}{i_{\mathrm{t}}-i}\right)
\label{eq:eta2}
\end{equation}

两部分之和给出电极表面浓度造成的总电势损失：

\begin{equation}
\eta_{\mathrm{conc}}=\eta_1+\eta_2
=-\frac{RT}{nF}\left(1+\frac{1}{\alpha}\right)
\ln\!\left(\frac{i_{\mathrm{t}}}{i_{\mathrm{t}}-i}\right)
\label{eq:eta-conc}
\end{equation}

\begin{keycond}
\equref{eq:eta1}—\equref{eq:eta-conc} 的成立条件是\textbf{静态电极 + 过量支持电解质 +
Nernst 扩散层模型 + 高过电位下略去 Butler--Volmer 的阳极项}。量纲检查：
$[RT/(\alpha nF)]=\mathrm{V}$，$[\ln(\cdot)]$ 无量纲，成立。两个极限行为要记牢：
$i\to 0$ 时 $\eta_{\mathrm{conc}}\to 0$（传质不构成限制）；$i\to i_{\mathrm{t}}$ 时
$\eta_{\mathrm{conc}}\to\infty$（浓差极化发散，电流被钉在平台值上）。
\tocheck{未对照原页面图像：书中式（11-40）OCR 为
$\eta_{\mathrm{conc}}=\eta_1+\eta_2=-\frac{RT}{nF}(1+\frac{1}{\alpha})(\frac{i_{\mathrm{t}}}{i_{\mathrm{t}}-i})$，
疑漏对数号（否则量纲为 V 而右端无量纲），本分片按 \equref{eq:eta1}、\equref{eq:eta2}
的标准形式补回 $\ln$，系数 $(1+1/\alpha)$ 沿用原书。}
\end{keycond}

图~\ref{fig:series} 把本节与上一节的关系画成一张图：电荷转移与传质是\textbf{串联}的
两步，二者共用同一个电流，因此\textbf{总速率由较慢的一步决定}。

\begin{figure}[htbp]
\centering
\begin{tikzpicture}[x=1mm,y=1mm,>=Stealth,scale=0.95,transform shape]% width=\onecolwidth
% --- electrode (left) ---
\fill[black!12] (-2,0) rectangle (0,30);
\draw[axisline] (0,0) -- (0,30);
\node[rotate=90,anchor=south,font={\fontsize{8}{9.5}\selectfont}] at (-1.0,15) {电极（电子导体）};
% --- x axis ---
\draw[axisline,->] (0,-3) -- (70,-3) node[anchor=west,font={\fontsize{8}{9.5}\selectfont}]{本体溶液 $x$};
% --- IHP/OHP ---
\draw[accentsoft,dashed,line width=0.6pt] (4,0) -- (4,30);
\draw[accentsoft,dashed,line width=0.6pt] (8,0) -- (8,30);
\node[anchor=south,font={\fontsize{8}{9.5}\selectfont}] at (6,30) {IHP/OHP};
% --- double layer charges ---
\foreach \y in {4,9,14,19,24}{\node[font={\fontsize{8}{9.5}\selectfont}] at (2,\y) {$-$};}
\foreach \y in {5,10,15,20,25}{\node[font={\fontsize{8}{9.5}\selectfont}] at (6,\y) {$+$};}
% --- diffusion layer span ---
\draw[<->,thinline] (0,33) -- (38,33)
  node[midway,above,font={\fontsize{8}{9.5}\selectfont}]{扩散层 $\delta_{\mathrm{O}}$（Nernst 静止层）};
\draw[accentsoft,dotted,line width=0.6pt] (38,0) -- (38,30);
% --- concentration profile ---
\draw[mainline] (0,1) .. controls (11,1) and (25,20) .. (38,24) -- (68,24);
\node[anchor=west,font={\fontsize{8}{9.5}\selectfont}] at (40,28) {$C_{\mathrm{O}}(x)$};
\node[anchor=west,font={\fontsize{8}{9.5}\selectfont}] at (40,21) {$C_{\mathrm{O}}^*$（本体浓度）};
% --- arrows ---
\draw[->,mainline] (58,10) -- (42,10);
\node[anchor=west,font={\fontsize{8}{9.5}\selectfont}] at (59,10) {扩散};
\draw[->,mainline] (13,17) -- (2,17);
\node[anchor=west,font={\fontsize{8}{9.5}\selectfont},text width=26mm] at (15,18)
  {电荷转移\\ $k_{\mathrm{f}},k_{\mathrm{b}}$};
\end{tikzpicture}
\caption{静态电极界面的双电层与扩散层示意，以及传质与电荷转移的串联位置（示意）。
读者应看三处：紧贴电极的 \textbf{IHP/OHP} 与电荷符号代表双电层，其充放电构成
\textbf{非法拉第电流}；向右延伸的 $\delta_{\mathrm{O}}$ 是\textbf{扩散层}，其中的浓度
梯度决定传质速率；电流必须依次通过这两步，故总速率由较慢者决定。局限：本图为
Nernst 静止层的简化图像，真实扩散层没有明确边界，仅在无对流且有支持电解质时近似
成立；双电层被简化为单层符号，未画出溶剂化与特性吸附。\schemamark\figsrc{据 S2
（下册）§11.1.5 与图 11-2 改绘；书内 p.817—819，PDF p.305—307。}}
\label{fig:series}
\end{figure}

\begin{crosscheck}
\textbf{怎么判断自己测的是传质还是电荷转移？} 三个互不依赖的检验：（i）
\emph{改变传质速率}——提高搅拌或旋转速率，若电流随之增大，说明传质仍在限制；
（ii）\emph{改变浓度}——若电流正比于 $C^*$（扩散控制）而不随催化剂本征性质变化，
说明读到的是极限电流；（iii）\emph{改变电极面积}——真实面积加倍而电流不成比例增加，
说明电流被溶液侧传质「钉住」。三条中只要有一条指向传质，就不能把该电流当作本征活性。
与流体动力学方法（11.3 起）的 Levich 关系互为回查：若 $j^{-1}$ 对 $\omega^{-1/2}$ 作图
得直线，其截距才对应本征动力学。
\end{crosscheck}

\subsection{控制电势的研究方法}

\subsubsection{实验装置与共同假定}

本节的方法都靠\textbf{已知的程序强制电极电势按预定函数变化}，同时测定电流随时间或
电势变化的特征曲线\src{S2 §11.2，书内 p.819 / PDF p.307}。装置包含函数发生器、恒电势仪
（potentiostat）与电解池：函数发生器设定电势程序，恒电势仪通过高阻抗反馈回路测量工作
电极与参比电极间的电势差，并\textbf{随时注入电流}以保证该电势差与程序一致
\src{图 11-3，书内 p.820 / PDF p.308}。对电极承担全部电流，参比电极因接在高阻抗回路上
而基本不取电流——这是「测量不扰动体系」的前提。

\begin{keycond}
本节所有方法都假定\textbf{面积体积比 $A/V$ 很小}，以保证测试过程中电活性物质的
\emph{本体}浓度不因电解而显著变化。实现方式是选用\textbf{微电极}（microelectrode）
并在足够大体积的溶液中做短时间实验。若违反此假定，$C^*$ 会随实验进程下降，
所有基于恒定 $C^*$ 的方程（Cottrell、Levich、Tafel 截距）都会随时间漂移，
表现为「同一电极重复扫描，峰高逐圈下降」。
\end{keycond}

\begin{measure}
\textbf{报告任何一条电势--电流曲线之前，这五项必须实测并记录。}
\textbf{（1）参比电极}：种类、内充液浓度、\qty{25}{\celsius} 下的标定值、实验前后漂移
（用饱和 \ce{KCl} 或标准缓冲液标定，漂移超过约 \qty{1}{mV} 的数据应作废）；
\textbf{（2）温度}：溶液温度及其稳定性，因为 $RT/F$ 直接进入所有斜率与截距；
\textbf{（3）未补偿电阻 $R_{\mathrm{u}}$}：由高频阻抗外推或电流中断法测出，记录
\textbf{补偿比例}（如 85\% 补偿）与残余 $iR_{\mathrm{u}}$ 降；\textbf{（4）双电层电容
$C_{\mathrm{d}}$ 与时间常数 $\tau=R_{\mathrm{u}}C_{\mathrm{d}}$}：决定哪些时间窗的数据
可用；\textbf{（5）面积基准}：几何面积 $A_{\mathrm{g}}$ 与真实面积 $A_{\mathrm{m}}$ 及其
定标方法（氢吸附电量或 $C_{\mathrm{d}}$）。这五项缺任何一项，Tafel 斜率与 $j_0$ 都
不可与文献直接比较。
\end{measure}

\subsubsection{电势阶跃法}

\paragraph{计时电流法} 给体系一个电势阶跃，测定电流随时间的变化
（chronoamperometry）\src{图 11-4，书内 p.820 / PDF p.308}。典型设置是：反应物
$\mathrm{O}$ 在初始电势 $E_1$ 不反应，阶跃到 $E_2$ 后以扩散极限速度被还原。
相关的两类变体用于获得更多信息：

\begin{itemize}
\item \textbf{取样电流伏安法}（sampled-current voltammetry）：做一系列电势阶跃，
  在初始条件相同的前提下分别测 $i$--$t$ 曲线，再用阶跃后某一固定时刻的电流对阶跃
  电势作图，得到「伏安图」\src{图 11-5，书内 p.821 / PDF p.309}。
\item \textbf{双电势阶跃计时电流法}（double potential step chronoamperometry）：
  先阶跃再跃回初始值，是反向技术，用于研究较复杂的电极反应，特别是中间产物的去向
  \src{图 11-6，书内 p.821 / PDF p.309}。
\end{itemize}

记录电流对时间的\textbf{积分}即通过的电量，相应方法统称库仑方法，最基本的是计时库仑法
与双电势阶跃计时库仑法，它们就是上述计时电流法的积分量
\src{图 11-7，书内 p.822 / PDF p.310}。

\paragraph{扩散控制的电势阶跃：Cottrell 方程} 指瞬间施加一个大电势阶跃，使电极表面
物质浓度为零；此后继续加大电势不再改变 $i$--$t$ 曲线，电流主要受溶液传质速率控制。
对平板电极，解线性扩散方程（Fick 第二定律）

\begin{equation}
\frac{\partial C_{\mathrm{O}}(x,t)}{\partial t}
=D_{\mathrm{O}}\frac{\partial^2 C_{\mathrm{O}}(x,t)}{\partial x^2}
\label{eq:fick2}
\end{equation}

边界条件为：初始 $C_{\mathrm{O}}(x,0)=C_{\mathrm{O}}^*$；
$\lim_{x\to\infty}C_{\mathrm{O}}(x,t)=C_{\mathrm{O}}^*$；阶跃后表面
$C_{\mathrm{O}}(0,t)=0\ (t>0)$ \src{式 11-41—11-44，书内 p.822 / PDF p.310}。解得
\textbf{Cottrell 方程}

\begin{equation}
i(t)=i_{\mathrm{d}}(t)=\frac{nFAD_{\mathrm{O}}^{1/2}C_{\mathrm{O}}^*}{(\pi t)^{1/2}}
\label{eq:cottrell}
\end{equation}

\src{式 11-45，书内 p.822 / PDF p.310}。表面附近电活性物质的贫化效应造成电流与
$t^{1/2}$ 呈\textbf{倒数}关系，这正是电解速度受扩散控制的标志。

\begin{keycond}
\eqref{eq:cottrell} 的成立条件：\textbf{（i）平板电极、半无限线性扩散}；
\textbf{（ii）阶跃足够大，使表面浓度瞬间降为零（极限扩散）}；\textbf{（iii）无对流、
有惰性支持电解质（电迁移被抑制）}；\textbf{（iv）$C^*$ 在实验期间不变（$A/V$ 小）}。
量纲检查：$[nFAD^{1/2}C^*/(\pi t)^{1/2}]
=(\mathrm{C.mol^{-1}})(\mathrm{cm^2})(\mathrm{cm.s^{-1/2}})(\mathrm{mol.cm^{-3}})
/(\mathrm{s^{1/2}})=\mathrm{C.s^{-1}}=\mathrm{A}$，成立。
\textbf{实操上必须扣除两个非理想项}：电流起始段受恒电势仪输出能力限制；阶跃瞬间还有
非法拉第的\textbf{双电层充电电流}，它按电解池时间常数 $\tau=R_{\mathrm{u}}C_{\mathrm{d}}$
（$R_{\mathrm{u}}$ 为未补偿电阻，$C_{\mathrm{d}}$ 为双电层电容）指数衰减，即使经过
$5\tau$ 后对总电流仍有可观贡献，从中很难精确分离法拉第电流。此外长时间实验中浓度梯度
与振动会造成对流扰动，表现为\textbf{电流大于} Cottrell 计算值。
\end{keycond}

相应的浓度分布为

\begin{equation}
C_{\mathrm{O}}(x,t)=C_{\mathrm{O}}^*\left[1-\mathrm{erfc}
\left(\frac{x}{2(D_{\mathrm{O}}t)^{1/2}}\right)\right]
=C_{\mathrm{O}}^*\,\mathrm{erf}\!\left(\frac{x}{2(D_{\mathrm{O}}t)^{1/2}}\right)
\label{eq:conc-profile}
\end{equation}

\src{式 11-46、11-47，书内 p.823 / PDF p.311}。不同时刻的浓度分布显示电极附近的氧化态
被耗尽，且\textbf{表面浓度梯度随时间变小}——这正是电流随时间衰减的几何原因。

\paragraph{球形扩散与超微电极} 对球形扩散（如悬汞电极）需用球坐标下的扩散方程

\begin{equation}
\frac{\partial C_{\mathrm{O}}(r,t)}{\partial t}
=D_{\mathrm{O}}\left[\frac{\partial^2 C_{\mathrm{O}}(r,t)}{\partial r^2}
+\frac{2}{r}\frac{\partial C_{\mathrm{O}}(r,t)}{\partial r}\right]
\label{eq:fick2-sph}
\end{equation}

$r$ 为距球心的径向距离，电极半径 $r_0$，边界条件为
$C_{\mathrm{O}}(r,0)=C_{\mathrm{O}}^*(r>r_0)$、
$\lim_{r\to\infty}C_{\mathrm{O}}(r,t)=C_{\mathrm{O}}^*$、
$C_{\mathrm{O}}(r_0,t)=0\,(t>0)$ \src{式 11-48—11-51，书内 p.824 / PDF p.312}。解得

\begin{equation}
i_{\mathrm{d}}(t)=nFAD_{\mathrm{O}}C_{\mathrm{O}}^*
\left[\frac{1}{(\pi D_{\mathrm{O}}t)^{1/2}}+\frac{1}{r_0}\right]
=i_{\mathrm{d}}(\text{线性})+\frac{nFAD_{\mathrm{O}}C_{\mathrm{O}}^*}{r_0}
\label{eq:spherical}
\end{equation}

\src{式 11-52、11-53，书内 p.824 / PDF p.312}。量纲检查：第二项
$[nFADC^*/r_0]=(\mathrm{C.mol^{-1}})(\mathrm{cm^2})(\mathrm{cm^2.s^{-1}})
(\mathrm{mol.cm^{-3}})/\mathrm{cm}=\mathrm{A}$，与第一项同量纲。
关键结论是：\textbf{球形扩散多出一个与时间无关的项}，这使电流在长时间区趋向稳态，
而不是像平板电极那样衰减到零。浓度分布的形式与线性情形非常相似，只差一个
$r_0/r$ 系数；若扩散层远薄于电极半径，球形电极的行为与平面电极无异。

\paragraph{粗糙度与超微电极} 实际电极表面总不像理想条件那样平滑洁净，故引入粗糙度
$\rho=A_{\mathrm{m}}/A_{\mathrm{g}}$：$A_{\mathrm{m}}$ 为微观（真实）面积，可由测双电层
电容或测电极表面形成单分子层所需电量来测定；$A_{\mathrm{g}}$ 为几何面积，即电极表面
正投影所得的面积 \src{书内 p.824 / PDF p.312}。\textbf{所有电流密度都必须说明用的是
哪个面积。} 超微电极因尺寸小而使实验结果与理论计算相近，故被广泛使用；本书界定其
临界尺寸小于 \qty{10}{nm}。短时间区扩散层厚度小于临界尺寸，可用 Cottrell 方程
（半无限线性扩散）计算电流；长时间区扩散层厚度较大，电流趋于稳态或准稳态

\begin{equation}
i_{\mathrm{ss}}=nFAm_0C_{\mathrm{O}}^*
\label{eq:iss-ch11}
\end{equation}

\src{式 11-55，书内 p.825 / PDF p.313}。\tocheck{未对照原页面图像：式中 $m_0$ 在原书
未在前文定义，按上下文应为传质系数（$\mathrm{cm.s^{-1}}$），其与 $D_{\mathrm{O}}$ 及
电极几何的关系（半球形 $m_0=2D/r_0$、球形 $4D/r_0$、嵌盘 $4D/(\pi r_0)$ 等）书中未给，
本分片不代为指定。}

\paragraph{取样电流伏安法的响应与信息边界} 对可逆电极反应，平板电极（满足半无限线性
扩散）对\textbf{任意}单电势阶跃的电流响应为

\begin{equation}
i(t)=\frac{nFAD_{\mathrm{O}}^{1/2}C_{\mathrm{O}}^*}
{\pi^{1/2}t^{1/2}\,(1+\theta)}
\label{eq:sampled-rev}
\end{equation}

它是可逆体系对电势阶跃的一般响应式，极限扩散（$\theta\to0$）时即还原为 Cottrell 方程
\src{式 11-56，书内 p.825 / PDF p.313}。此外，\textbf{电极形状并不是关键}，只要满足
「扩散层厚度小于电极曲率半径」即可。超微球形电极上述模型的稳态区响应为
$i=\dfrac{nFAD_{\mathrm{O}}C_{\mathrm{O}}^*}{(1+\xi_2\theta)r_0}$
\src{式 11-57，书内 p.825 / PDF p.313}。
\tocheck{未对照原页面图像：式（11-56）（11-57）中的符号 $\theta$、$\xi_2$ 在原书未在
本分片可见范围内定义，OCR 亦可能把 $\theta$/$\vartheta$ 或下标记错。本分片保留原符号
并只使用其定性结论（$\theta\to0$ 退化为 Cottrell），未据此做数值计算。}

可逆体系下取样电流伏安法能给出与不能给出的信息很不相同：

\begin{itemize}
\item \textbf{能}：由波高测浓度，用于校正或测量标准添加物的浓度；可逆波的形状与位置
  反映电极反应的能量依赖性，由此可得\textbf{标准电势、自由能、平衡常数}；半波电势接近
  形式电势，故可用于估算尚未定性的化学体系的电势。
\item \textbf{不能}：\textbf{不能得到动力学信息}。可逆波只报告热力学，因为传质已经完全
  接管了电流的大小。
\end{itemize}

对准可逆与不可逆反应，平板电极的响应形如
$i=FA\left(k_{\mathrm{f}}C_{\mathrm{O}}^*-k_{\mathrm{b}}C_{\mathrm{R}}^*\right)
\exp(H^2t)\,\mathrm{erfc}(Ht^{1/2})$，球形电极的稳态响应另含
$D_{\mathrm{O}}$、$r_0$ 与速率常数之比 \src{式 11-58、11-59，书内 p.825 / PDF p.313}。
\tocheck{未对照原页面图像：式（11-58）中 $H$ 的定义、式（11-59）中括号内的
$(\delta+1)$ 结构均无法从 OCR 复原，本分片只用其定性形式与下述结论，未代入数值。}
其信息边界为：准可逆与不可逆波的\textbf{极限电流平台仍完全受扩散控制}，故仍可用于求
$C^*$、$n$、$A$、$D$、$r_0$ 等任何对 $i_{\mathrm{d}}$ 有贡献的参数；准可逆波有时可
给出形式电势的近似值，而不可逆波因半波电势偏离形式电势\textbf{无法求得热力学信息}；
但完全不可逆波的波形与位置可以提供\textbf{动力学}信息。

\paragraph{计时电量法} 计时电量法记录电量与时间的关系 $Q(t)$，常用于代替计时电流法，
有三个优势：信号随时间增长，暂态后期受阶跃瞬间非理想电势变化的影响较轻微，信噪比好；
暂态电流中的随机噪声因积分而平滑；\textbf{可以区分双电层充电、吸附物质的电极反应与
扩散反应物法拉第反应对电量的贡献}，因而便于研究表面过程。对 Cottrell 公式积分得

\begin{equation}
Q_{\mathrm{d}}=\frac{2nFAD_{\mathrm{O}}^{1/2}C_{\mathrm{O}}^*t^{1/2}}{\pi^{1/2}}
\label{eq:chronocoul}
\end{equation}

但实际电量还含双电层充电与还原吸附的某种氧化态所贡献的部分，$Q$ 对 $t^{1/2}$ 的直线
一般不通过原点。这些电量只在瞬间出现、不随时间累积，故作为两个与时间无关的附加项写出：

\begin{equation}
Q=\frac{2nFAD_{\mathrm{O}}^{1/2}C_{\mathrm{O}}^*t^{1/2}}{\pi^{1/2}}
+Q_{\mathrm{dl}}+nFA\Gamma_{\mathrm{O}}
\label{eq:chronocoul-full}
\end{equation}

式中 $Q_{\mathrm{dl}}$ 为电容电量，$nFA\Gamma_{\mathrm{O}}$ 为表面吸附的
$\mathrm{O}$ 还原的法拉第分量 \src{式 11-60、11-61，书内 p.826 / PDF p.314}。
量纲检查：$[nFAD^{1/2}C^*t^{1/2}/\pi^{1/2}]=(\mathrm{C.mol^{-1}})(\mathrm{cm^2})
(\mathrm{cm.s^{-1/2}})(\mathrm{mol.cm^{-3}})(\mathrm{s^{1/2}})=\mathrm{C}$，成立；
$[nFA\Gamma]=\mathrm{C.mol^{-1}}\cdot\mathrm{cm^2}\cdot\mathrm{mol.cm^{-2}}=\mathrm{C}$，
同量纲。\textbf{这条式子的实践价值在于把「非法拉第背景」变成了可读出的截距}：
用不同电势阶跃做 $Q$--$t^{1/2}$ 图，截距对电势作图为直线，斜率即 $C_{\mathrm{d}}$，
由此可测真实面积；剩余截距来自吸附物种。这一方法常用于测定电活性物质的表面余量，
但需要再做反向实验。

\begin{dataflow}
\textbf{从一条 $i$--$t$ 曲线到结论，只能走这五步。} （1）确认阶跃电势足够大，
表面浓度确实为零（否则 \equref{eq:cottrell} 不适用）。（2）判断时间窗：
$t<5\tau$（$\tau=R_{\mathrm{u}}C_{\mathrm{d}}$）的数据必须丢弃，双电层充电尚未衰减；
长时间区要检查 $i\,t^{1/2}$ 是否仍为常数。（3）作 $i$--$t^{-1/2}$ 图，直线斜率给出
$nFAD^{1/2}C^*/\pi^{1/2}$；斜率过原点才说明纯扩散控制。（4）由斜率反解 $D$ 或 $C^*$
时，必须固定另一个量；用\emph{同一个}面积定义（$A_{\mathrm{m}}$ 或 $A_{\mathrm{g}}$）。
（5）只有在确认 $i$--$t^{-1/2}$ 直线的截距为零之后，才能说「本体系在此时标下是纯扩散
控制」；截距非零通常意味着还有电荷转移阻力或吸附贡献。需要本征动力学时，改做取样电流
伏安法并处理 \equref{eq:sampled-rev} 形式的响应。
\end{dataflow}

\subsubsection{电势扫描法}

电势扫描法随时间扫描电势并记录电流--电势曲线，以获得体系完整的电化学行为。主要分为
\textbf{线性扫描伏安法}（linear sweep voltammetry，LSV）与\textbf{循环伏安法}
（cyclic voltammetry，CV）\src{S2 §11.2.2、图 11-9，书内 p.826—827 / PDF p.314—315}。

\paragraph{从一条 CV 读出什么：\ce{Pt} 单晶在 \ce{H2SO4} 中的例子} 以 \ce{Pt} 在
\qty{0.5}{mol.L^{-1}} \ce{H2SO4} 中的 CV（\qty{50}{mV.s^{-1}}）为例
\src{图 11-10，书内 p.827 / PDF p.315}。测量前先把单晶表面在氢氧火焰中退火：一是完全
除去表面吸附的杂质，二是使原子完全按规则排列。退火过程中 \ce{Pt} 表面会吸附氧形成氧
吸附层。扫描从\textbf{开路电势 OCP} 出发往低电势方向进行。

读图顺序与每个峰的归属：先出现的峰来自表面已吸附的氧随电势降低被电还原；第二个过程是
溶液中 \ce{H+} 得电子还原为氢原子并吸附在 \ce{Pt} 表面，而正向扫描时对应的是吸附氢的
电氧化，即该反应的逆反应。其他单晶上这些氧峰与氢峰同样存在，但\textbf{峰的形状和位置
不同}，因为表面 \ce{Pt} 原子的排列方式不同——因此峰形特征可用来辨别晶面是 (100)、
(110) 还是 (111)。

\paragraph{氢吸附电量与真实面积} 对氢吸附峰积分，所得面积即为氢吸附产生的电量
（电流$\times$时间 = 电量）。由电量可算出吸附氢原子的数目；不同晶面 $Q_{\mathrm{H}}$
不同。再由晶面上原子的排列方式与原子直径可算出单位面积的 \ce{Pt} 原子数，两者比较
发现数值很接近，说明\textbf{一个 \ce{Pt} 原子吸附一个 H}。由于电活性物质（H 和 O）反复
氧化还原吸脱附，表面会逐渐粗糙，单晶表面最终变成多晶表面；多晶表面由若干小单晶面
组成，反映的是这些单晶的平均行为 \src{图 11-11，书内 p.828 / PDF p.316}。

多晶 \ce{Pt}（\ce{Pt} 丝、纳米 \ce{Pt} 黑）的 CV 分三个区域：左侧峰对应氢的吸附与脱附，
不同峰来自氢在不同晶面上的吸附；右侧峰是氧的吸附与脱附，即 \ce{Pt} 表面的氧化与还原，
不同峰代表不同氧化态（可能有 \ce{Pt-OH}、\ce{PtO_x} 等，更详细的归属仍未知）；中间区域
是\textbf{双电层区}，该区电流由双电层充放电引起、没有氧化还原反应。对该区积分可求得
双电层电量。由于一个 \ce{Pt} 原子只吸附一个氢原子，\textbf{$Q_{\mathrm{H}}/Q$ 即 \ce{Pt}
的真实表面积}——这就是用 H 的吸脱附测定 \ce{Pt} 真实表面积的方法
\src{书内 p.829 / PDF p.317}。

\begin{dataflow}
\textbf{为什么必须用真实面积。} 实验能测的是整个电极面积上产生的电流，面积越大电流
越大。因此\textbf{忽略电极面积、直接用电流大小表示反应速率的快慢，既不科学也不准确}。
规范做法是用\textbf{电流密度}（单位电极面积上的电流）表示反应速率，而这个面积必须是
\textbf{真实（电化学活性）面积}，即由氢吸附电量 $Q_{\mathrm{H}}$ 或双电层电容定标得到的
面积，而不是电极的几何尺寸。报告数据时必须同时给出：面积基准、定标方法、以及
$\rho=A_{\mathrm{m}}/A_{\mathrm{g}}$ 的数值或范围。
\end{dataflow}

\paragraph{可逆性判据} 由 CV 的形状可以判断电极反应属于可逆、准可逆还是不可逆，其判据
与「能得到什么信息」直接对应：

\begin{keycond}
\textbf{可逆（Nernst 型）}：峰电流 $i_{\mathrm{p}}$ 正比于 $v^{1/2}$；峰电位差
$\Delta E_{\mathrm{p}}=|E_{\mathrm{pa}}-E_{\mathrm{pc}}|\approx 2.303RT/(nF)$，即
\qty{25}{\celsius} 下约 $59/n$ mV，且\textbf{不随扫描速率变化}；$i_{\mathrm{pa}}/i_{\mathrm{pc}}
\approx1$；$E_{\mathrm{p}}$ 与 $v$ 无关。\textbf{准可逆}：$\Delta E_{\mathrm{p}}$ 随 $v$
增大而增大，峰仍然成对。\textbf{不可逆}：只出现单向峰，$E_{\mathrm{p}}$ 随 $v$ 增大而
向更负（阴极）或更正（阳极）方向移动，约每十倍 $v$ 移动 $30/(\alpha n)$ mV。
\textbf{判据的成立条件是反应物与产物都溶于溶液且无后续化学反应}；若产物发生吸附或
偶联化学反应，$i_{\mathrm{pa}}/i_{\mathrm{pc}}$ 会显著偏离 1，此时不能直接判为不可逆。
\tocheck{未对照原页面图像：分片 a 的源文本（图 11-9、11-10 说明）未展开 CV 的定量判据
公式，本节判据按电化学标准形式重建；若本书 11.2.2 后有专门的判据小节，应由后一分片
或统一核验时对齐符号（是否用 $D^{1/2}$、是否含 $n^{3/2}$、是否讨论 $\Delta E_{\mathrm{p}}$
的严格表值）。}
\end{keycond}

\begin{pitfall}
\textbf{控制电势实验的五类系统性误判，按出现频率排序。}
\textbf{（1）未做 iR 补偿导致 Tafel 斜率失真。} 实测电势含 $iR_{\mathrm{u}}$ 降：
$E_{\text{实测}}=E_{\text{真}}+iR_{\mathrm{u}}$，故按 $\eta$--$\lg i$ 作图得到的表观斜率
为 $b_{\text{表观}}=b+2.303R_{\mathrm{u}}A$（以电流作图时）或 $b+2.303R_{\mathrm{u}}j$
（以电流密度作图时）。这意味着\textbf{高电流段被系统性抬高}，Tafel 图向上弯折、斜率变大；
电流越大失真越重。识别办法：改变电极与参比电极的距离或溶液电导率，若斜率随之变化，
则斜率是 $R_{\mathrm{u}}$ 的产物而非 $\alpha$ 的产物。
\textbf{（2）传质未排除时把「扩散极限」当本征活性。} 若 Tafel 区取在
$i\gtrsim0.1i_{\mathrm{t}}$ 的区间，测到的实际上是浓差极化（\equref{eq:eta-conc}），
其斜率依赖 $i_{\mathrm{t}}$ 而与 $\alpha$ 无关。判据：$i$--$t^{-1/2}$ 直线不通过原点，
或提高搅拌后电流增大。
\textbf{（3）参比漂移。} 漂移直接叠加到所有 $\eta$ 上，表现为整条曲线平移——斜率不动而
$j_0$ 整体错一个数量级。
\textbf{（4）双电层充电电流被计入法拉第电流。} $\tau=R_{\mathrm{u}}C_{\mathrm{d}}$ 内的
电流以非法拉第成分为主；把它当作法拉第电流会使低过电位区（恰是 $j_0$ 的取值区）的
电流全部偏高，从而\textbf{系统性高估 $j_0$}。排除办法：同一电势程序下做空白电解液
对照，或在 $5\tau$ 之后取样（计时电量法用 $Q$--$t^{1/2}$ 截距分离）。
\textbf{（5）几何面积与电化学活性面积混用。} $j$ 的分母用 $A_{\mathrm{g}}$ 还是
$A_{\mathrm{m}}$ 可使数值相差数倍到数十倍；纳米催化剂的 $A_{\mathrm{m}}$ 尤其依赖定标
方法。跨文献比较 $j_0$ 或质量活性时，必须先核对面积基准是否一致。
\end{pitfall}

\paragraph{线性扫描与循环扫描的分工} LSV 只扫一个方向，得到的是一条单向的 $i$--$E$
曲线，适合测起始电势、半波电势，也是极化曲线（Tafel 图）的常规来源；CV 扫过去再扫回来，
额外给出\textbf{逆反应的响应}，因而能识别中间产物、判断可逆性、测表面吸脱附电量。
二者的共同弱点是：\textbf{扫描过程中表面状态与扩散层都在变}，所以从 CV 直接读出的
「活性」是扫描速率、传质与动力学三者耦合的结果。

\subsubsection{极谱法和脉冲伏安法}

\paragraph{极谱法与滴汞电极} 极谱法由海洛夫斯基使用滴汞电极（dropping mercury
electrode，DME）创建。使用时 \ce{Hg} 从毛细管流出，在管口逐渐长成直径约 \qty{1}{mm}
的汞滴，重力超过表面张力后落下 \src{S2 §11.2.3、图 11-12，书内 p.829 / PDF p.317}。
汞滴长大期间若有电解反应，就有随时间变化的电流通过，$i$--$t$ 曲线同时反映电极长大与
电解贫化两种效应。汞滴落下时能在一定程度上混合溶液，\textbf{极大减弱贫化效应}；因为
每一滴都在新鲜溶液中产生，若在汞滴寿命（2～6 s）期间维持电极电势不变，则每个汞滴的
生长过程都可以当成一次电势阶跃实验。

经典 DME 有两个主要缺点：\textbf{（i）汞滴面积不断变化}，使扩散处理复杂，双电层充电
还产生不断变化的背景电流；\textbf{（ii）时间尺度受汞滴寿命限制}，一般在 0.5～10 s 之间，
在此范围外工作很难。静态滴汞电极（static mercury drop electrode，SMDE）通过电控阀门
控制汞滴生长时间以保证每次落下的汞滴大小相同，既保留 DME 的优点，面积又不变，简化了
扩散处理并降低了背景噪声 \src{图 11-13，书内 p.830 / PDF p.318}。

\paragraph{脉冲伏安法的五类} 脉冲伏安法起源于经典极谱技术，最初目的是抑制 DME 中汞滴
连续扩张引起的充电电流；自 SMDE 发明后逐渐向静态电极靠近。脉冲伏安法主要分五种：
断续极谱法/阶梯伏安法、常规脉冲伏安法、反向脉冲伏安法、示差脉冲伏安法、方波伏安法。

\textbf{（1）断续极谱法/阶梯伏安法}：在汞滴寿命期间，法拉第电流不断上升而充电电流逐渐
下降 \src{图 11-14，书内 p.831 / PDF p.319}。为得到最佳的法拉第电流与充电电流之比以
提高灵敏度，应\textbf{刚好在汞滴下落前采样} \src{图 11-15，书内 p.831 / PDF p.319}。
断续极谱法每一滴汞对应一个循环，适用于各种可周期更新的电极，但对静止电极（如悬汞电极）
不适用，此时需用阶梯伏安法。阶梯伏安法与断续极谱法类似，只是不需要电极更新，与电势
扫描方法类似，可替代线性扫描并能抑制充电电流，但其\textbf{低信噪比以及速度与分辨率
不能同时满足}的缺陷限制了使用。

\textbf{（2）常规脉冲伏安法}：断续极谱法只在很短时间内采集电流，其余时间的电流没有
用途，而采样前的电流会使待测物质贫化、降低灵敏度。为禁止这种无用的电解，此法仅在
\textbf{采样前的一小段时间}才施加脉冲电势 \src{图 11-16，书内 p.832 / PDF p.320}，
从而大大减小等待期间的电解对测量精度的影响。

\textbf{（3）反向脉冲伏安法}：与常规脉冲伏安法类似，但脉冲电势逐渐变至电化学物质不电解
的电势，属反向技术，常用于测定前一步初始电解产物的行为。

\textbf{（4）示差脉冲伏安法}：与常规脉冲伏安法类似，但 $E_{\mathrm{b}}$ 在每个循环中都有
所提升，实验记录结果是\textbf{两个采样电流之差}。其法拉第电流不会超过常规脉冲伏安法
测得的电流，因此\textbf{高灵敏度不是来自增强法拉第响应，而是来自降低背景噪声}
\src{书内 p.832 / PDF p.320}。

\begin{pitfall}
脉冲类方法最容易被误用的一点：把灵敏度提高当成活性提高。示差脉冲伏安法给出的电流是
\textbf{差分信号}，其数值与常规脉冲伏安法的电流不可直接比较，更不能跨方法比较
$j_0$。跨方法比较电流或电流密度之前，必须确认采样时刻、脉冲宽度/振幅、扫描速率与
面积基准都一致。
\end{pitfall}

\paragraph{四类方法的测量对象与信息边界} 表~\ref{tab:methods} 把本节四种控制电势方法
按「测量什么、能得到什么、不能得到什么、关键条件」并列，便于在拿到一条曲线时先确定
它属于哪一格。

\begin{table}[htbp]
\centering
\caption{四类控制电势方法的测量对象与信息边界}
\label{tab:methods}
{\srcsize
\begin{tabularx}{\columnwidth}{@{}lXX@{}}
\toprule
方法 & 测量对象 & 能得到 / 不能得到 \\
\midrule
电势阶跃（计时电流 / 计时电量） &
$i$--$t$ 或 $Q$--$t$，表面浓度被强制为零 &
\textbf{能}：$D$、$C^*$、$n$、面积；由 $Q$--$t^{1/2}$ 截距分离 $C_{\mathrm{d}}$ 与吸附量。
\textbf{不能}：单次阶跃给不出 $j_0$ 与 $\alpha$（电势未作自变量） \\
\addlinespace
取样电流伏安 / 双电势阶跃 &
系列阶跃后固定时刻的电流；反向阶跃的电流 &
\textbf{能}：可逆波给 $E^{\ominus\prime}$、$\Delta G$、平衡常数；准可逆/不可逆波给
$k^0$、$\alpha$ 等动力学信息，极限平台仍给 $C^*$、$n$、$A$、$D$、$r_0$。
\textbf{不能}：可逆波不给动力学；不可逆波不给热力学 \\
\addlinespace
线性扫描 / 循环伏安 &
$i$--$E$ 曲线（CV 含回扫） &
\textbf{能}：峰电位定性与可逆性判据；峰电流随 $v^{1/2}$ 的依赖；氢/氧吸脱附电量定标
真实面积；极化曲线给 Tafel 斜率与 $j_0$。
\textbf{不能}：扫描中表面与扩散层同时在变，直接读出的活性是三者耦合结果；双电层充电
电流混在总电流中 \\
\addlinespace
脉冲伏安（NPV/DPV/SWV 等） &
脉冲前后电流之差或采样电流 &
\textbf{能}：抑制充电电流、提高信噪比；DPV 以差分降背景；反向脉冲测初始电解产物。
\textbf{不能}：差分电流与常规电流不可直接比较；速度与分辨率不能同时满足；断续极谱法
不适用于静止电极 \\
\bottomrule
\end{tabularx}}
\end{table}

\begin{crosscheck}
\textbf{本节方法的互补关系。} 电势阶跃确定传质参数与面积基准，取样电流伏安把电势作为
自变量分离出动力学，电势扫描给出可逆性与表面吸脱附信息，脉冲方法在低浓度/弱信号
条件下压低背景。因此一个完整的电催化动力学结论通常需要\textbf{至少两类方法互证}：
例如 CV 给出的 Tafel 斜率，必须用某一种「传质已知」的方法（旋转圆盘或取样电流伏安）
确认斜率确实出现在 $i\ll i_{\mathrm{t}}$ 的区间内，否则该斜率可能只是浓差极化的表观
结果。反过来，若某方法给出的 $j_0$ 与另一方法相差超过一个数量级，先查\textbf{面积基准
与 iR 补偿}，再查参比电极标定，最后才怀疑模型。
\end{crosscheck}
}{}
\IfFileExists{chapters/parts/ch11-b.tex}{\subsection{控制电流方法与整体电解}
\subsection{流体动力学方法}
}{}
\IfFileExists{chapters/parts/ch11-d.tex}{% ================================================================
%  chapters/parts/ch11-d.tex
%  第 11 章  电化学催化研究方法（分片 d）
%  节次：§ 11.5.1  阻抗谱基本原理
%  源：S2（下册）第 11 章 §11.5.1，书内 p.855—861 / PDF p.343—349
%      （实测：该节源文本在 sources/parts/ch11-b.txt 第 864—1094 行，
%        起于「11.5.1 阻抗谱基本原理」标题，止于「11.5.2 阻抗技术在电催化
%        领域中的应用」标题之前，含原书图 11-32 至图 11-41、表 11-1，
%        公式编号（11-115）至（11-120）；源文本未给（11-121）以后的编号。）
%
%  【顺序问题——已于 2026-09-29 由装配层解决】
%  装配器 chapters/ch11.tex 原先按 a→b→c→d 的 \IfFileExists 顺序引入，
%  本分片实际排在 ch11-c（§11.5 后半与 §11.6）之后，而不是 §11.5.1 应处的
%  §11.4 之后，由此产生三处节序错乱：
%    (1) 正文顺序变成 §11.4 → §11.5.1 → §11.5.2 → §11.6 → §11.5（ch11-c）：
%        「阻抗谱基本原理」与其应用之间被 §11.5.2 与整节 §11.6 隔开；
%    (2) ch11-c 开篇「阻抗的基本定义、等效电路推导与 Nyquist／Bode 图的定义
%        写在上文 §11.5.1」这句话，在装配后指向的是**下文**，指代方向与实际相反；
%    (3) 编号 11.5.1 会出现在 11.5 之后（LaTeX 的 \subsection 计数器只按出现
%        顺序编号，本分片标题沿用分片 b 的层级，故书内编号将与源书目录不符）。
%  处置办法（三选一，均须改本分片之外的文件）：
%    (A) 在 chapters/ch11.tex 里把 ch11-d 的 \input 提到 ch11-b 之后、
%        ch11-c 之前（最小改动，编号与指代同时恢复正确）；
%    (B) 把本分片内容并入 ch11-b 末尾（需改 ch11-b，与并发作业冲突）；
%    (C) 接受现状，把本分片标题改为「阻抗谱基本原理（补）」并把 ch11-c 的
%        指代句改为「见前文」。
%  【已执行 (A)】chapters/ch11.tex 现为 a→b→d→c→e，实测渲染出的节标题序列为
%  11.1 → 11.2 → 11.3 → 11.4 → §11.5.1（本分片，编号 11.5）→ §11.5（分片 c，
%  编号 11.6）→ §11.6（分片 c，编号 11.7）。注意：本分片标题仍是 \subsection，
%  故**书内**自动编号为 11.5，与源书目录的「11.5.1」不同号。这是源书节号与
%  本项目层级不一致的通例——源书「11.5.1」以下的内容在本项目里被拆成两个
%  \subsection（原理与应用），源书节号一律用 \src{}／正文「§11.5.1」标注出处，
%  不映射到 book 计数器；分片 c 的指代句与 \srcnote 已同步改为
%  「上一节（§11.5.1）」并写明该编号指源书节号。
%  本分片按任务约定**不写 \chapterstart**，只用 \subsection，以免与装配层的
%  章计数器冲突（章计数器由 ch11-a 唯一推进，故调序不影响章号）。
%
%  【诚实纪律】本分片看不到原书页面图像。公式按「源文本 OCR 上下文 + 学科
%  标准形式」重建；凡是源文本未给而由学科标准形式补出的内容，都在正文或
%  图题里显式说明；凡数值、条件与符号约定存疑处，一律标 \tocheck{} 写明存疑点。
%  本分片**没有**核对过任何页面图像，正文中不作与此相反的表述。
% ================================================================

\subsection{阻抗谱基本原理}

电化学阻抗谱（\term{EIS}{electrochemical impedance spectroscopy}{电化学阻抗谱}）给电化学系统
施加一个\textbf{频率不同、振幅很小的交流正弦电势波}，测量交流电势与电流信号的比值（即系统的
阻抗）随正弦波频率的变化，或者是阻抗的相位角 $\varphi$ 随频率的变化。由于扰动电信号是交流信号，
电化学阻抗谱也称\textbf{交流阻抗谱}（alternating current impedance spectroscopy）。从装置上看得
更直观：用波形发生器产生一个小幅正弦电势信号，经恒电位仪施加到电化学系统上，把输出的电流或电势
信号经转换后送入锁相放大器或频谱分析仪，输出阻抗及其模量或相位角；改变正弦波的频率，就得到一
系列不同频率下的阻抗、阻抗模量与相位角\src{S2 §11.5.1，书内 p.855—856 / PDF p.343—344}
（原书图 11-32——电化学阻抗法实验装置示意图）。

\begin{probout}
阻抗实验回答的是「这套电极过程一共有几个可分辨的时间尺度、各是多大」。已知的是一条稳态极化
曲线只能给出所有步骤过电位的\textbf{总和}；要问的是能否把总响应按时间尺度拆开，分别读出溶液
电阻、双电层电容、电荷转移电阻与扩散阻抗。可观测量是每个频率下的两个实数——阻抗的实部与虚部
（或等价的模与相位角），而待定的等效电路元件数通常多于两个，因此本方法的核心风险不是「测不到」，
而是\textbf{解不唯一}：能拟合数据的电路不止一个，而只有一个是真实机理。（超）微电极把未补偿
电阻与时间常数压到极小，因而特别适合在低电导率介质中研究电化学过程，并能在很小的
$R_{\mathrm{u}}C_{\mathrm{d}}$ 时间常数下测量快速反应\src{S2 §11.4.3，书内 p.855 / PDF p.343}。
\end{probout}

\subsubsection{交流电、相量与简单元件的响应}

一个纯正弦电压可表示为

\begin{equation}
e=E\sin\omega t
\label{eq:eis-d-sine}
\end{equation}

式中 $\omega$ 为角频率，它是 $2\pi$ 乘以用 Hz 表示的常规频率值；$E$ 为幅值
\src{式（11-115），书内 p.856 / PDF p.344}。可以把电压用旋转矢量（或\textbf{相量}）表示：
观察到的电压 $e$ 是任意时刻投影在某一特定轴（通常取 $0$°）上相量的分量（见原书图 11-33）。
两个正弦信号——电流 $i$ 与电压 $e$——可以表示成以同样频率旋转的独立相量 $\dot I$ 和 $\dot E$，
它们通常不同相，相差一个相角 $\varphi$，以 $\dot E$ 为参考信号，$\varphi$ 即相对它测出。图 11-34
中电流滞后于电压，此时 $\varphi<0$，可表示为

\begin{equation}
i=I\sin(\omega t+\varphi)
\label{eq:eis-d-phase}
\end{equation}

\src{式（11-116），书内 p.857 / PDF p.345；原书图 11-34 给出频率为 $\omega$ 的交流电压与电流
相互关系的相量图}。\tocheck{未对照原页面图像：源文本此式即印作
$i=I\sin(\omega t+\varphi)$，本稿照印、不再改写成 $\sin(\omega t-\varphi)$——按本节的测量相角
约定 $\varphi\equiv-\arg Z$（见下文式~\eqref{eq:eis-d-zmod} 后的约定行），该符号是自洽的：电流
滞后即 $\varphi<0$；纯电容时式（11-118）的 $i=(E/X_C)\sin(\omega t+\pi/2)$ 给出
$\varphi=+\pi/2$。该式的符号未能逐字与页面核对。}

简单元件对电压的响应是本节的积木。纯电阻 $R$ 上施加 $e=E\sin\omega t$ 时，由欧姆定律电流为
$(E/R)\sin\omega t$，相角为零（原书图 11-35）。把电阻换成纯电容 $C$，所需的基本关系式是
$q=it$、$q=Ce$ 与 $i=C(\mathrm{d}e/\mathrm{d}t)$，因此

\begin{equation}
i=\omega CE\cos(\omega t),\qquad
i=\frac{E}{X_C}\sin\!\left(\omega t+\frac{\pi}{2}\right)
\label{eq:eis-d-cap}
\end{equation}

式中 $X_C$ 为\textbf{容抗}，其值为 $1/\omega C$\src{式（11-117）、（11-118），书内 p.857—858 /
PDF p.345—346}。相关电流的公式与电阻相似，只是以 $X_C$ 代替 $R$，且相角由 $0$ 变为 $\pi/2$：
\textbf{电流导前于电压}（原书图 11-36）。由于相角改变后矢量图扩展为一个平面，用复数符号表示
相量就成为必要：规定横坐标分量为实部、纵坐标分量为虚部，并乘以 $j=\sqrt{-1}$。这里引入复数符号，
是想使矢量的各分量都成为直线——实部与虚部这两种形式在相角可测量的意义上讲都是真实的。

\begin{equation}
\dot E=-jX_C\dot I
\label{eq:eis-d-capohm}
\end{equation}

\src{式（11-119），书内 p.858 / PDF p.346}。比较 $\dot E=\dot I R$ 与式~\eqref{eq:eis-d-capohm}
可知，$X_C$ 具有电阻的量纲，但与 $R$ 不同——\textbf{其值随频率的增加而下降}。这一点是后面所有
读图规则的根源：电容元件在频率轴上有一个确定的位置，改变频率等于改变它在电路中的「权重」。

\paragraph{串联叠加与阻抗的定义} 在电阻 $R$ 与电容 $C$ 的串联上施加电压 $\dot E$，其值无论何时
都等于电阻与电容上电压降之和：

\begin{equation}
\dot E=\dot E_R+\dot E_C=\dot I\,(R-jX_C)=\dot I Z
\label{eq:eis-d-series}
\end{equation}

电压与电流通过一个称为\textbf{阻抗}的矢量 $Z=R-jX_C$ 联系在一起
\src{式（11-120）至（11-122），书内 p.858 / PDF p.346}。更一般地，阻抗写成实部与虚部之和：

\begin{equation}
Z(\omega)=Z_{\mathrm{Re}}+jZ_{\mathrm{Im}},\qquad
Z_{\mathrm{Re}}=R,\qquad Z_{\mathrm{Im}}=-X_C=-\frac{1}{\omega C}
\label{eq:eis-d-zdef}
\end{equation}

$Z$ 的幅值与相角为

\begin{equation}
\lvert Z\rvert^{2}=R^{2}+X_C^{2}
=\left(Z_{\mathrm{Re}}\right)^{2}+\left(Z_{\mathrm{Im}}\right)^{2},
\qquad
\tan\varphi=-\frac{Z_{\mathrm{Im}}}{Z_{\mathrm{Re}}}=\frac{X_C}{R}
=\frac{1}{\omega RC}
\label{eq:eis-d-zmod}
\end{equation}

\src{式（11-123）至（11-125），书内 p.858—859 / PDF p.346—347}。式~\eqref{eq:eis-d-zmod} 中的
$\varphi$ 是\textbf{测量相角}，本节统一取 $\varphi\equiv-\arg Z$（$\arg Z$ 为 $Z$ 的辐角），
即\textbf{容性为正}：纯电阻 $\varphi=0$、纯电容 $\varphi=+\pi/2$、纯电感 $\varphi=-\pi/2$；
对混合体则观察到介于相应两者之间的相角。\textbf{Nyquist 图纵轴取 $-Z''$（向上为正）与 Bode 图的
相角曲线同属这一约定}——同一条谱在两种图上必须给出同一个 $\varphi$，换图不换号。
\textbf{更为复杂的电路，可以按与电阻相同的规则通过合并阻抗来分析——串联阻抗的总阻抗是各阻抗值
之和；并联则用导纳 $Y\equiv1/Z$ 表述：总导纳是各支路导纳之和，即}
\begin{equation}
Y_{\text{并}}=\frac{1}{Z_{\text{并}}}=\sum_i\frac{1}{Z_i},
\qquad
Z_{\text{并}}=\Bigl(\sum_i\frac{1}{Z_i}\Bigr)^{-1}
\label{eq:eis-d-parallel}
\end{equation}
\src{书内 p.859 / PDF p.347。原文作「并联阻抗的总阻抗是单个阻抗值倒数之和」，本式即按「总阻抗的
倒数是各阻抗倒数之和」理解}。核对两个数：两个 $10$ \unit{\ohm} 的电阻并联，
$Y_{\text{并}}=0.1+0.1=0.2$ S，故 $Z_{\text{并}}=1/0.2=5$ \unit{\ohm}（若把「倒数之和」
$0.2$ S 直接当成阻抗值，会得到 $0.2$ \unit{\ohm}，与 $5$ \unit{\ohm} 差 $25$ 倍）；$R$ 与
$C$ 并联时 $Y=1/R+j\omega C$，低频极限 $Z\to R$（$\varphi\to0$）、高频极限
$Z\to-j/(\omega C)$（$\varphi\to+\pi/2$）。这两条合并规则加上式~\eqref{eq:eis-d-series}，就是
全部等效电路的数学基础；后文所有电路的阻抗表达式都只是这两条规则的反复应用。
\tocheck{未对照原页面图像：源文本式（11-123）以下的 $Z_{\mathrm{Im}}$ 印作 $+X_C$、式（11-125）
印作 $\tan\varphi=X_C/R$，二者与式（11-121）$Z=R-jX_C$ 相差一个符号；本稿取
$Z_{\mathrm{Im}}=-X_C$（即 $j$ 的系数）并令 $\varphi\equiv-\arg Z$，两式即同时成立，也与源文本
p.859「纯电容 $\varphi=\pi/2$」一致。源文本式（11-125）OCR 只剩
$\tan\varphi=1/\omega RC$ 的骨架，符号与绝对值未能逐字核对。}

\begin{principle}
\textbf{把电化学系统看成等效电路。} 利用电化学阻抗谱研究一个电化学系统时，它的基本思路是把
电化学系统看作一个\textbf{等效电路}——由电阻（$R$）、电容（$C$）和电感（$L$）等基本元件按串联
或并联等不同方式组合而成；通过 EIS 测定等效电路的构成及各元件的大小，再利用这些元件的电化学
含义分析电化学系统的结构和电极过程的性质\src{书内 p.859 / PDF p.347}（原书图 11-37——电阻、
电容及电感等元件的串联和并联）。EIS 研究是在\textbf{可逆}电极反应过程的基础上发展起来的，
用线性元件作为等效元件，构成能给出与实测 EIS 一样谱图的等效电路\src{书内 p.860 / PDF p.348}。
「可逆」与「线性元件」这两个前提是整套方法的边界：它要求扰动足够小，使电流–电势关系可以线性化，
从而能够用传递函数（而不是非线性微分方程）描述系统。
\end{principle}

\subsubsection{两种常用画法与其定义}

电化学阻抗技术通过测定不同频率下扰动信号 $X$ 与响应信号 $Y$ 的比值，得到各频率下阻抗的实部、
虚部、模值和相位角，再把这些量绘成各种形式的曲线，就得到电化学阻抗谱
\src{书内 p.859 / PDF p.347}。常用的阻抗谱有两种。

\paragraph{Nyquist 图（能斯特图）} 以阻抗的实部 $Z_{\mathrm{Re}}$ 为横轴、虚部 $Z_{\mathrm{Im}}$
为纵轴，\textbf{图中的每个点代表不同的频率，左侧的频率高，成为高频区，右侧的频率低，成为低频
区}\src{书内 p.859 / PDF p.347}。本书从本节起统一按电化学惯例以 $-Z''$ 向上为正作图
（即纵轴画 $-Z_{\mathrm{Im}}$），以便常见的容性数据落在第一象限；这一点在读书中的数据图时必须
先确认，否则整个图会上下翻转。

\paragraph{Bode 图（波特图）} 包括\textbf{两条}曲线，它们的横坐标都是频率的对数 $\lg f$，纵坐标
一个是阻抗模值的对数 $\lg\lvert Z\rvert$，另一个是阻抗的相位角 $\varphi$
\src{书内 p.859 / PDF p.347}。两图的分工是本节的第一个实用结论：\textbf{Nyquist 图把频率信息
抹掉了}——图上每个点对应哪个频率必须另查数据表；\textbf{Bode 图把频率显式保留}，因而判断「哪个
频段被哪个过程支配」更可靠。因此\textbf{同一组数据必须两图并看}：用 Nyquist 图读几何量（截距、
直径、斜率），用 Bode 图定频段与数时间常数。

\paragraph{三类极限体系的特征谱} 源文本给出三种基本情况，可以当作读图的「字典」
\src{书内 p.859—860 / PDF p.347—348}：

\begin{itemize}
\item \textbf{纯电阻}：Nyquist 图上表现为 $Z'$ 轴（横轴）上的一个点，该点到原点的距离为电阻值的
大小；Bode 图中 $\lg\lvert Z\rvert$–$\lg f$ 为一条水平直线，相角 $\varphi$ 为 $0$°且不随测量
频率变化。
\item \textbf{纯电容体系}：Nyquist 图上表现为与纵轴（$-Z''$ 轴，按本节 $\varphi\equiv-\arg Z$ 的
约定向上为正）重合的一条直线；Bode 图中 $\lg\lvert Z\rvert$–$\lg f$ 是斜率为 $-1$ 的直线，
$\varphi$ 为 $+90$°。
\item \textbf{Warburg 阻抗} $Z_{\mathrm{W}}$：Nyquist 图上表现为相角为 $45$°的直线；Bode 图上
表现为斜率为 $-1/2$、$\varphi$ 为 $+45$°的直线。
\end{itemize}

这幅「字典」的价值在于它给出了\textbf{三条可直接反查的判据}：$|Z|$ 对频率的幂律指数 $0$、$-1$、
$-1/2$ 分别对应电阻、电容与半无限扩散。测得一段数据的局部斜率后，先问「它接近 $0$、$-1/2$ 还是
$-1$」，再去谈具体机理——这个顺序可以挡住大部分把 Warburg 当半圆、把 CPE 当电容的误判。

\src{式（11-115）至（11-125）依次见书内 p.856—859 / PDF p.344—347，原书图 11-33 至图 11-37
依次见书内 p.857—859 / PDF p.345—347。}

\subsubsection{等效元件与四类电极的特征谱}

等效电路中的常用等效元件如表~\ref{tab:eis-d-elements} 所示，系源文本表 11-1 的内容
\src{原书表 11-1，书内 p.860 / PDF p.348}。

\begin{table}[htbp]
\centering
\caption{等效电路中的常用等效元件，以及读图时各自的成立条件。第四列是本分片
按学科标准形式补出的：\textbf{它回答的是「凭什么把这个图形特征归到这个元件上」}，条件不成立时
读出的数不是「误差大一点」，而是物理归属错了。\src{据 S2 表 11-1，书内 p.860 / PDF p.348
整理。}}
\label{tab:eis-d-elements}
\begin{tabularx}{\onecolwidth}{@{}p{17mm}p{8mm}p{40mm}X@{}}
\toprule
名称 & 符号 & 源文本给出的意义 & 读图用法与成立条件 \\
\midrule
溶液电阻 & $R_{\Omega}$ & 工作电极与参比电极之间的电阻 &
$R_{\Omega}$ 与接触电阻、引线电阻在谱上不可分辨——EIS 本身分不开它们。故
「高频截距 $=R_{\Omega}$」是\textbf{约定}而非测量结果：须另用改变鲁金毛细管位置（截距随之变）或
电流中断法标定，才可把截距当溶液电阻用。 \\
双电层电容 & $C_{\mathrm{dl}}$ & 工作电极与电解质之间的电容 &
与 $R_{\mathrm{ct}}$ 并联，构成一个时间常数 $\tau=R_{\mathrm{ct}}C_{\mathrm{dl}}$。仅在电位区间内
无法拉第电流时，它才单独决定响应（此时半圆不出现，谱退化为垂直线）。 \\
极化电阻 & $R_{\mathrm{p}}$ & 当电位远离开路电位时，导致电极表面产生电流，电流受到反应动力学和
反应物扩散的控制 &
$R_{\mathrm{p}}$ 是\textbf{动力学与扩散的总和}，不是纯动力学量；只在扩散可忽略时
$R_{\mathrm{p}}\approx R_{\mathrm{ct}}$。用它算 $i_0$ 前必须先确认低频没有扩散尾。 \\
电荷转移电阻 & $R_{\mathrm{ct}}$ & 电化学反应动力学控制 &
对应半圆直径，但须「弧闭合、只有一个主时间常数、扩散不参与」；三者缺一，直径里就混着别的贡献。 \\
扩散电阻（Warburg 阻抗） & $Z_{\mathrm{W}}$ & 反应物从溶液本体扩散到电极反应界面的电阻 &
对应低频 $45$°直线；成立条件是\textbf{半无限}线性扩散——扩散层厚度远小于体系最小边界尺度，且
低频点没有因体系漂移而失真。 \\
界面电容 & $C_{\mathrm{c}}$ & 通常每一个界面之间都会存在的一个电容 &
与 $C_{\mathrm{dl}}$ 并联出现（涂层、晶界、催化剂层／膜界面）。二者时间常数接近时只能读出两者的
组合量，无法按单个界面分开。 \\
电感 & $L$ & 影响电流的变化 &
图上表现为第四象限的弧。仪器电感应与频率无关、且随接线方式改变——能靠改接线把它改掉的，就不是
化学信息。 \\
\bottomrule
\end{tabularx}

\noindent{\srcsize 表中「意义」一列逐条取自源文本表 11-1（书内 p.860 / PDF p.348）；第四列为
按学科标准形式补出的操作判据，并已把源文本未给出的分界条件写明——须与原书 p.860—861 的四类
电极讨论对照。}
\end{table}

\paragraph{四类电极的 Nyquist 图} 理想极化电极、溶液电阻可以忽略的电化学极化电极、溶液电阻不能
忽略的电化学极化电极、以及电化学极化和浓差极化同时存在的电极，其阻抗谱均可用 Nyquist 图与 Bode 图
表示；源文本以 Nyquist 图为代表逐一介绍\src{书内 p.860 / PDF p.348}。

\begin{enumerate}
\item \textbf{理想极化电极}是不发生电极反应的电极，其 Nyquist 图为一条距 $Z''$ 轴为 $R_{\Omega}$
且垂直于实轴 $Z'$ 的直线；由直线在 $Z'$ 轴上的交点到原点的距离可求得 $R_{\Omega}$（原书图 11-38）
\src{书内 p.860 / PDF p.348}。
\item \textbf{溶液电阻可以忽略的电化学极化电极}对应一半径为 $R_{\mathrm{ct}}/2$ 的半圆，
由圆的直径可求得 $R_{\mathrm{ct}}$；\textbf{圆的顶点对应的 $Z''$ 最大，由对应的角频率 $\omega$
及得到的反应电阻，可求得双电层电容}
\begin{equation}
C_{\mathrm{dl}}=\frac{1}{\omega_{\max}R_{\mathrm{ct}}},
\qquad
\omega_{\max}=2\pi f_{\max}=\frac{1}{R_{\mathrm{ct}}C_{\mathrm{dl}}}
\label{eq:eis-d-cdl}
\end{equation}
（原书图 11-39）\src{书内 p.860—861 / PDF p.348—349}。\tocheck{未对照原页面图像：源文本此句
OCR 呈「由对应的角频率 $\omega$ 及得到的反应电阻，可以求得双电层电容 $C_{\mathrm{dl}}=1/(\omega
R_{\mathrm{ct}})$」，未写明是取半圆顶点的 $\omega$；本分片按顶点条件补出 $\omega_{\max}$ 的
定义，并保留原式形式。（若误把 $f_{\max}$ 当作 $\omega_{\max}$ 代入，结果会差 $2\pi$ 倍——
这是本节最容易犯的算术错误。）}
\item \textbf{溶液电阻不能忽略的电化学极化电极}：设扩散过程引起的阻抗可以忽略、电极过程由电荷
传递过程（电化学反应步骤）控制，所得 Nyquist 图同样是一半径为 $R_{\mathrm{ct}}/2$ 的半圆，
由其可求得溶液电阻 $R_{\Omega}$ 和电荷转移电阻 $R_{\mathrm{ct}}$，由半圆顶点的 $\omega$ 可求得
$C_{\mathrm{dl}}$（原书图 11-40）\src{书内 p.861 / PDF p.349}。
\item \textbf{电化学极化和浓差极化同时存在的电极}：电极过程由电荷传递过程和扩散过程共同控制，
其 Nyquist 图由\textbf{高频区一个半径为 $R_{\mathrm{ct}}/2$ 的半圆}和\textbf{低频区一条 $45$°的
直线}构成——半圆表示电荷传递过程为控制步骤，直线表示由电极反应的反应物或产物的扩散控制
（原书图 11-41）\src{书内 p.861 / PDF p.349}。
\end{enumerate}

这四类讨论给出的正是三条最常用的读数规则。但源文本是按「理想情形」逐类陈述的，读者若把每一类
的结论直接搬到实测数据上，就会丢掉前提。下面把三条规则各自的成立条件写全，它们分别对应
「高频截距」「半圆直径」「低频斜线」三个图形特征。

\paragraph{规则一：高频截距 $=R_{\Omega}$} 成立需要同时满足两条：
（1）\textbf{频率足够高}，高到 $C_{\mathrm{dl}}$ 的阻抗小到可以忽略、$R_{\mathrm{ct}}$ 被完全旁路，
即 $\omega R_{\mathrm{ct}}C_{\mathrm{dl}}\gg1$；
（2）\textbf{高频极限处的接触电阻与引线电感可以独立排除}——接触电阻在物理上与溶液电阻串联，
EIS 无法从谱上把它与 $R_{\Omega}$ 分开，只能靠改变电极压紧力或参比位置、看截距是否跟着变来判定；
引线电感则使实部先减小后增大，把截距抬偏甚至读不出。两条都不满足时，截距里含 $R_{\mathrm{ct}}$
成分，读出的「溶液电阻」偏大。

\paragraph{规则二：半圆直径 $=R_{\mathrm{ct}}$} 成立需要同时满足三条：
（1）\textbf{半圆闭合}，即频率范围覆盖了从 $\omega\tau\gg1$ 到 $\omega\tau\ll1$ 的全过程；
（2）\textbf{这段弧只由一个时间常数支配}——两个时间常数相差不到约两个数量级时，两个半圆会重
叠成一个压扁的弧，此时用「直径」读出的既不是 $R_{\mathrm{ct}1}$ 也不是 $R_{\mathrm{ct}2}$；
（3）\textbf{扩散不参与}，否则弧的低频一侧被 Warburg 尾「拉住」，外推得到的截距偏大。
要注意\textbf{压扁也可能来自 CPE 而非两个过程}，二者仅凭一张 Nyquist 图无法区分，必须看 Bode
相角图上有几个峰。

\paragraph{规则三：低频 $45$°斜线 $=$ 扩散控制} 成立需要在\textbf{两端}同时成立：
（1）\textbf{上界}——频率低到 Warburg 项主导，即 $\sigma\omega^{-1/2}$ 已远大于 $R_{\mathrm{ct}}$
的贡献；
（2）\textbf{下界}——频率还不够低，扩散层厚度 $\sqrt{2D/\omega}$ 仍\textbf{远小于}体系的最小
边界尺度（薄液层厚度、膜厚、多孔层厚度、对流边界层厚度）。把这两个条件写成可算的判据：半无限
扩散的有效下限约在
\begin{equation}
f_{\mathrm{low}}\gtrsim\frac{D}{2\pi\delta^{2}}
\label{eq:eis-d-fdiff}
\end{equation}
$f<f_{\mathrm{low}}$ 之后扩散层「看到」了边界，$45$°直线必然弯曲（弯向实轴或转为接近竖直的
电容性直线）。\tocheck{未对照原页面图像：式~\eqref{eq:eis-d-fdiff} 是本分片按半无限扩散的
特征扩散长度 $\delta\sim\sqrt{2D/\omega}$ 补出的操作性判据，源文本 §11.5.1 未给出该式；其常数
因子（$2\pi$ 与 $2$ 的取法）取决于所用特征长度的定义，须与 §11.5.1 的符号约定对照。}
数量级感受：$D=1\times10^{-5}$ cm$^{2}$·s$^{-1}$ 且薄液层厚 $\delta=100$ µm 时
$f_{\mathrm{low}}\approx0.03$ Hz——也就是\textbf{常规阻抗实验从 10 mHz 起扫
的那些低频点，早已越过了薄层体系的半无限下限}。在不搅拌的厚层溶液中（边界是本体溶液）这一条
通常自动满足；在薄膜电极、薄液层电解池与多孔催化剂层里，它是必须先算一遍的条件。

\subsubsection{等效电路的阻抗表达式（Randles 电路、Warburg 阻抗与 CPE）}

源文本把等效电路作为「基本思路」提出，并给出表 11-1 的元件清单与四类电极的 Nyquist 图，但
\textbf{没有写出这些电路的总阻抗表达式}。下面把最常用的几个表达式按学科标准形式补出，作为读图
（下一小节）的工具；凡源文本已给出形式的，都在公式后标出页码，凡由标准形式补出的，逐条说明。

\paragraph{Randles 电路与溶液电阻、电荷转移电阻、双电层电容} 把溶液电阻 $R_{\Omega}$（源文本符号，
即通常写作 $R_{\mathrm{s}}$ 者）与一个并联支路串联，并联支路由电荷转移电阻 $R_{\mathrm{ct}}$ 与
扩散阻抗 $Z_{\mathrm{W}}$ 串联后，再与双电层电容 $C_{\mathrm{dl}}$ 并联，即得最常用的
\textbf{Randles 电路}：

\begin{equation}
Z=R_{\Omega}+\left[\frac{1}{R_{\mathrm{ct}}+Z_{\mathrm{W}}}
+j\omega C_{\mathrm{dl}}\right]^{-1}
\label{eq:eis-d-randles}
\end{equation}

\tocheck{未对照原页面图像：Randles 电路的具体拓扑（$Z_{\mathrm{W}}$ 与 $R_{\mathrm{ct}}$ 串联、
二者整体与 $C_{\mathrm{dl}}$ 并联）与式~\eqref{eq:eis-d-randles} 是按源文本图 11-38 至图 11-41
的四类电极 Nyquist 图所隐含的结构、以及学科标准形式重建的；源文本 §11.5.1 未写出该式，故拓扑
本身未能与页面核对。}
当扩散可忽略（$Z_{\mathrm{W}}\to0$）时，式~\eqref{eq:eis-d-randles} 退化为
$R_{\Omega}$ 串联一个并联 $R_{\mathrm{ct}}C_{\mathrm{dl}}$ 网络，即规则二讨论的单个半圆；此时
\begin{equation}
Z'=R_{\Omega}+\frac{R_{\mathrm{ct}}}{1+(\omega R_{\mathrm{ct}}C_{\mathrm{dl}})^{2}},
\qquad
-Z''=\frac{\omega R_{\mathrm{ct}}^{2}C_{\mathrm{dl}}}{1+(\omega R_{\mathrm{ct}}C_{\mathrm{dl}})^{2}}
\label{eq:eis-d-rc}
\end{equation}
两式消去 $\omega$ 即得圆心在实轴上 $R_{\Omega}+R_{\mathrm{ct}}/2$ 处、半径为 $R_{\mathrm{ct}}/2$
的圆方程——这是「半圆直径 $=R_{\mathrm{ct}}$」与「半圆顶点频率
$\omega_{\max}=1/(R_{\mathrm{ct}}C_{\mathrm{dl}})$」两条读图规则的来源。\textbf{两条规则的几何
前提是同一个}：谱必须是一个完整的半圆；谱被压扁时式~\eqref{eq:eis-d-rc} 不再成立，圆的圆心
落到实轴下方，直径也不再等于 $R_{\mathrm{ct}}$。

\paragraph{Warburg 阻抗与扩散系数} 半无限线性扩散给出

\begin{equation}
Z_{\mathrm{W}}=\sigma\,\omega^{-1/2}(1-j),
\qquad
Z'_{\mathrm{W}}=-Z''_{\mathrm{W}}=\sigma\,\omega^{-1/2},
\qquad
\sigma=\frac{RT}{n^{2}F^{2}A\sqrt{2}}
\left(\frac{1}{c_{\mathrm{O}}^{*}\sqrt{D_{\mathrm{O}}}}
+\frac{1}{c_{\mathrm{R}}^{*}\sqrt{D_{\mathrm{R}}}}\right)
\label{eq:eis-d-warburg}
\end{equation}

实部与虚部恒相等，故相角恒为 $+45$°（按 $\varphi\equiv-\arg Z$），在 Nyquist 图上是一条斜率 $1$ 的直线
\src{相角 $45$°这一结论见书内 p.859 / PDF p.347}。由 $45$°直线在 $\omega^{-1/2}$ 轴上的斜率读出
$\sigma$ 后，若已知 $n$、$A$ 与浓度，即可估出扩散系数 $D$——这是阻抗谱少有的「能给出传质定量
参数」的场合。\tocheck{未对照原页面图像：$\sigma$ 的表达式（含系数 $RT/(n^{2}F^{2}A\sqrt{2})$
与括号内两项）以及 $Z_{\mathrm{W}}$ 的符号约定按学科标准形式重建，源文本 §11.5.1 只给出
「相角 $45$°」这一图形结论，未给出 $\sigma$ 的定义式，故未能与页面核对。}

\paragraph{常相位元件（CPE）} 实测半圆几乎从不以实轴为直径，而是一个\textbf{压扁的弧}。把
$C_{\mathrm{dl}}$ 换成\textbf{常相位元件}（constant phase element，CPE）即可描述这一现象：

\begin{equation}
Z_{\mathrm{CPE}}=\frac{1}{Q(j\omega)^{n}},\qquad 0<n\le1,
\qquad
\varphi_{\mathrm{CPE}}=\frac{n\pi}{2}
\label{eq:eis-d-cpe}
\end{equation}

$n=1$ 时 $Q$ 就是电容（单位为 F）；$n=0$ 时退化为电阻。\textbf{$n$ 偏离 $1$ 的物理含义是电极
表面不均匀}：粗糙表面、多晶面的不同晶面取向、吸附层覆盖不均、以及电流分布不均，都会使界面的
弛豫时间不是一个值，而是一段分布——分布越宽，$n$ 越小、弧越扁。\textbf{把 CPE 当理想电容去算
电容值，是在这一节里最容易发生的系统性错误}：$Q$ 的量纲是 $\mathrm{S.s^{n}}$，它不是电容；由
$Q$ 与 $n$ 反求有效电容必须做量纲换算
\begin{equation}
C_{\mathrm{eff}}=\left(Q\,R_{\mathrm{ct}}^{\,1-n}\right)^{1/n}
\label{eq:eis-d-cpeceff}
\end{equation}
而有效电容的取法（是否用 $R_{\Omega}$ 还是 $R_{\mathrm{ct}}$、或两者的并联）本身有不同约定，
因此\textbf{报告结果时必须同时给出 $n$ 与所用换算式的出处}，否则读者无法复算。
\tocheck{未对照原页面图像：式~\eqref{eq:eis-d-cpe} 中 $Q$ 的符号与 $0<n\le1$ 的区间、以及
式~\eqref{eq:eis-d-cpeceff} 的换算形式按学科标准形式重建，源文本 §11.5.1 未出现 CPE 或 $n$ 的
讨论（原书表 11-1 只有理想双电层电容 $C_{\mathrm{dl}}$），故该元件与原书的对照关系未能核对。}

\begin{principle}
\textbf{为什么等效电路能成立，以及它成立的代价。} 式~\eqref{eq:eis-d-randles} 一类的表达式之所以
可以描述一个真实的电极／溶液界面，靠的是三条前提：\textbf{线性}（扰动足够小，$i$–$E$ 关系可在
工作点附近线性化）、\textbf{因果}（响应不先于扰动）、\textbf{稳定}（测量期间体系不随时间漂移，
即处于稳态）。三条合起来等价于「系统可用线性时不变传递函数描述」，这正是阻抗定义
$Z(\omega)=E(\omega)/I(\omega)$ 的适用条件。Kramers--Kronig 关系是这三条前提的数学后果，因此
它可以反过来被用作\textbf{数据合格性的检验}：残差在全部频率上都小，才说明上面三条成立；否则
该组数据不应进入拟合\src{书内 p.859 / PDF p.347 给出「观察体系在稳态时对扰动的响应，是频率域
的测量」这一前提；Kramers--Kronig 检验为其数学后果，属本分片按学科标准形式补充}。
代价也要说清：\textbf{等效电路是等价描述，不是唯一描述}。同一组数据可以有多个电路满足它，
「线性元件」也只是在工作点附近成立——因此阻抗谱给出的是\textbf{同一稳态下}各步骤的时间尺度与
权重，它不能单独回答各步骤之间的先后关系。
\end{principle}

\subsubsection{测量装置与测量条件}

\begin{measure}
\textbf{（1）装置链路。} 波形发生器产生小幅正弦电势 $\to$ 恒电位仪施加到电化学系统 $\to$ 输出的
电流／电势信号经转换 $\to$ 锁相放大器或频谱分析仪 $\to$ 输出阻抗及其模量或相位角；改变正弦波
频率即得一系列不同频率下的阻抗、模量与相位角\src{书内 p.855—856 / PDF p.343—344}。链路上每一环
都会在某个频段引入自己的误差：恒电位仪的带宽与相位误差在高频端，锁相放大器的时间常数在低频端，
电缆与接线的电感与电容在高频端。\textbf{因此每次实验都应先用已知的 dummy cell（纯电阻或
$R$–$C$ 串联）验证整条链路}，特别是相位——相位误差不能靠多测平均消除。

\textbf{（2）扰动幅值。} 必须小到使响应线性，又大到使信噪比足够。判据不是记住一个数字，而是做
\textbf{幅度减半检验}：把交流幅值减半后重测，若全频段谱形不变，说明仍在线性区；若 $R_{\mathrm{ct}}$
随幅值系统性变小，说明已进入非线性区。燃料电池测量的交流振幅一般选择在 $5\sim15$ mV
\src{书内 p.862 / PDF p.350}；催化剂活性测量因电位控制更直接，通常取更小值。

\textbf{（3）频率范围。} 范围本身是结论的一部分。燃料电池测量中扫描频率范围一般在
1 mHz～100 kHz 之间\src{书内 p.862 / PDF p.350}。范围的选择须由
\textbf{待求的时间常数}反推：高频端要跨过 $1/(2\pi R_{\mathrm{ct}}C_{\mathrm{dl}})$ 至少两个
数量级（否则高频截距含 $R_{\mathrm{ct}}$ 成分），低频端要覆盖目标过程（扩散尾要看到 $45$°、
吸附弛豫要看到相角峰）。按照规则三下界的量级估算，低频端的点数应密集到足以分辨斜率，而不是按
对数均匀取点。

\textbf{（4）测量模式。} 阻抗测量一般有\textbf{恒电位模式}（扰动为交流电位，响应为测量电流）与
\textbf{恒电流模式}（扰动为交流电流，响应为测量电位）两种，从两种模式得到的结果基本上没有明显
不同，但各有优缺点\src{书内 p.862 / PDF p.350}。在恒电位模式中，一个小的电位都可能导致大的电流
变化，从而导致恒电位仪过载，因此尝试不同的交流电压以确保恒电位仪不过载是必要的；在催化剂活性
的测量中，控制电位能更好地控制电化学反应的发生，因此一般采用恒电位模式。在恒电流模式中，交流
和直流大小很容易被运算放大器确定，能有效防止电池和恒电位仪过载，因此能采用相对大的交流电，
从而得到更多的阻抗信息；对于电池阻抗测量，一般电流比较大，更适合用恒电流模式，恒电位模式容易
导致过载\src{书内 p.862 / PDF p.350}。

\textbf{（5）稳态必须先行建立。} 阻抗定义要求体系在测量期间是线性时不变系统，因此
\textbf{「稳态未建立就测」是本节一切误判的共同源头}。必须先在目标电位恒电位极化到电流漂移小于
约每分钟 $1\%$，再开始扫频；对慢过程（覆盖度还在演化、催化剂正在失活），应先在\textbf{一个
固定频率上监测 $\lvert Z\rvert$ 随时间的变化}，待其稳定后再扫全谱。\tocheck{未对照原页面图像：
「电流漂移小于约每分钟 $1\%$」与「先在单一频率监测 $\lvert Z\rvert$」两条操作判据按学科标准
形式补出，源文本 §11.5.1 只给出「观察体系在稳态时对扰动的响应」这一前提，未给出量化判据。}

\textbf{（6）记录条件。} 报告阻抗结果必须同时给出：交流幅值与模式、频率范围与每十倍频程的点数、
直流偏置电位、电极面积与几何、溶液组成与温度、参比电极种类、以及是否扣除与如何扣除
$R_{\Omega}$。缺任何一项，别处无法复现，也无法判断读出的数是否可比。
\src{本节（1）（2）（4）三条中凡标页码处为源文本 §11.5.2 对测量条件的直接陈述
（书内 p.862 / PDF p.350）；本分片只取与 §11.5.1 原理直接相关的条件，应用层面的讨论
（膜透过、催化剂、MEA、单池、电堆）不在本节范围内。}
\end{measure}

\subsubsection{从一组阻抗数据到等效电路：数据流与读图实例}

\begin{dataflow}
\textbf{第一层（构造）：} 由相量关系 $Z=Z'+jZ''$ 把每个频率的一对实数量组织成一个复数；
$|Z|^{2}=(Z')^{2}+(Z'')^{2}$，$\tan\varphi=-Z''/Z'$\src{式（11-123）至（11-125），
书内 p.858—859 / PDF p.346—347}。源文本在此处把 $Z_{\mathrm{Re}}$ 记作 $R$、把虚部印作
$X_C$（符号不严，见上文式~\eqref{eq:eis-d-zmod} 后的 \tocheck），那是式（11-121）的串联 $RC$
情形；本节其余处一律用一般记号 $Z'$、$Z''$，相角取测量约定 $\varphi\equiv-\arg Z$。

\textbf{第二层（定向）：} Nyquist 图上横轴 $Z'$ 向右增大、纵轴 $-Z''$ 向上增大，
\textbf{左侧频率高、右侧频率低}\src{书内 p.859 / PDF p.347}。因此沿曲线从左端走到右端 =
频率从高降到低；这条约定是后面所有「先看到谁」判断的基础。

\textbf{第三层（不变量检验）：} 纯电阻对应 $|Z|$–$f$ 水平线与 $\varphi=0$°；纯电容对应斜率
$-1$ 与 $\varphi=+90$°；Warburg 阻抗对应斜率 $-1/2$ 与 $\varphi=+45$°
\src{书内 p.859—860 / PDF p.347—348}。\textbf{先在 Bode 图上量出局部幂律指数，再回 Nyquist
图谈几何}——顺序颠倒就会把 Warburg 的直线段当圆弧读。

\textbf{完整数据流：} 逐频率数据 $\to$（检验）Kramers--Kronig 残差与正反扫重合性 $\to$
（变换）$\lg\lvert Z\rvert$ 与 $\varphi$ 对 $\lg f$ 作图、数相角峰个数定时间常数个数 $\to$
（分频段）高频平台定 $R_{\Omega}$、中频弧定 $R_{\mathrm{ct}}$ 与 $C_{\mathrm{dl}}$（或 $Q$ 与
$n$）、低频尾定扩散 $\to$（拟合）选定电路拓扑做非线性最小二乘 $\to$（验收）残差随机分布、
参数置信区间合理、跨电位／跨转速趋势自洽 $\to$（归属）把每个参数接到一个物理量与一条成立
条件上。\textbf{这条链上任何一步被跳过，后面的结论都不成立}；其中「数相角峰个数」是唯一能定
「电路里至少需要几个时间常数」的步骤，不能由拟合替代。
\end{dataflow}

\begin{figure}[tbp]
\centering
\begin{tikzpicture}
\tikzset{declare function={
  dpw(\z)  = 10^(\z);
  dwn(\z)  = 2*3.14159265358979*dpw(\z);
  dtc(\z)  = 100*0.5e-6;
  dpr(\z)  = dwn(\z)*dtc(\z);
  didx(\z) = 10+100/(1+(dpr(\z))^2);
  didy(\z) = 100*dpr(\z)/(1+(dpr(\z))^2);
  dq0      = (0.5e-6)^(0.85)*100^(0.85-1);
  dcpn     = 0.85;
  % 角度一律用「度」：pgfmath 的 sin/cos/tan 取度不取弧度（R3-C02 的实测探针
  %   sin(0.85*pi/2)=0.02328 是 sin(1.335°) 的值）。原先写 dcpb=n*pi/2 等于把
  %   θ 当成 1.335 rad 用，sin/cos 全部算错，弧被压成一条贴着横轴的线。
  dcpb     = dcpn*90;
  % u = Rct*Q*ωⁿ 才是无量纲组合；原写法 (ω*100*Q)^n 把整个乘积升到 n 次，
  %   量纲和数值都不对。下面的 a、d 分别是分母 1+u(cosθ+j sinθ) 的实部与虚部。
  dcpu(\z) = 100*dq0*(dwn(\z))^(dcpn);
  dcpa(\z) = 1+dcpu(\z)*cos(dcpb);
  dcpd(\z) = dcpu(\z)*sin(dcpb);
  dcpx(\z) = 10+100*dcpa(\z)/((dcpa(\z))^2+(dcpd(\z))^2);
  dcpy(\z) = 100*dcpd(\z)/((dcpa(\z))^2+(dcpd(\z))^2);
  % CPE 弧顶的解析值：对 u 求极值（u=1）得 Z'=Rs+Rct/2、-Z''=(Rct/2)tan(nπ/4)。
  %   标记点由这两个式子给出，不再手写孤立坐标。
  dcpxp    = 60;
  dcpyp    = 50*tan(dcpn*45);
  % Randles 加半无限 Warburg：Z = Rs + [1/(Rct+σ(1-j)/√ω) + jωC]^{-1}
  %   （式~\eqref{eq:eis-d-randles}）。原写法把它当成 Rct∥Warburg 的两元并联
  %   —— 丢失了 Cdl 支路 —— 并且把有量纲的 σ/√ω 直接与 1 相加（R3-C03 的根因）。
  %   下面改用归一化导纳 Ŷ = Y·Rct（无量纲）：
  %     t = σ/(Rct√ω)，S = (1+t)²+t²，Ŷ = (1+t)/S + j[t/S + ω C Rct]
  %   再取 Z - Rs = Rct/Ŷ。中间量都远离 pgfmath 的数值上下限，低频段也不会
  %   因为两个小量相减而丢精度。
  dkwz(\z) = 40/(100*sqrt(dwn(\z)));
  dkwS(\z) = (1+dkwz(\z))^2+(dkwz(\z))^2;
  dkwA(\z) = (1+dkwz(\z))/dkwS(\z);
  dkwB(\z) = dkwz(\z)/dkwS(\z)+dwn(\z)*0.5e-6*100;
  dkwN(\z) = (dkwA(\z))^2+(dkwB(\z))^2;
  dkwx(\z) = 10+100*dkwA(\z)/dkwN(\z);
  dkwy(\z) = 100*dkwB(\z)/dkwN(\z);
  % R2-C06：原先还声明过 dm / dmt / dmph / diph 四个量，已删除——
  %   ① 全书 grep 确认它们从未被任何 \addplot 调用（纯死代码）；
  %   ② 数学上也不成立：`diph` 给出的是 atan(ωRC)，按本节 φ≡−arg Z
  %      （tanφ=1/ωRC）恰是 φ 的余角；并且 pgfmath 的三角函数取「度」，
  %      要算 tan(76.5°) 就得写 tan(76.5)，写成 tan(dm*90) 之类是弧度思维。
  %   留着会被后来者当成可用函数，故删除而不是改写。本图不画相角曲线。
}}

\begin{axis}[
  name=eisda,
  width=70mm, height=62mm,
  unit vector ratio=1 1 1,
  xlabel={$Z'$ / $\Omega$}, ylabel={$-Z''$ / $\Omega$},
  % 轴伸到 62 Ω（刻度仍到 60），让两个弧顶标记落在框内而不是压在边框上。
  xmin=0, xmax=62, ymin=0, ymax=62,
  xtick={0,20,40,60}, ytick={0,20,40,60},
]
  \addplot[mainline,black,domain=0.3:4.9,samples=220,variable=\z]
    ({didx(\z)},{didy(\z)});
  \addplot[mainline,black,dashed,domain=0.3:4.9,samples=220,variable=\z]
    ({dcpx(\z)},{dcpy(\z)});
  \addplot[only marks,mark=*,mark size=1pt,black]
    coordinates {(60,50)};
  \addplot[only marks,mark=square*,mark size=1pt,black]
    coordinates {({dcpxp},{dcpyp})};
  \node[anchor=north west,font={\fontsize{7}{8.5}\selectfont}]
    at (axis cs:1.2,60.5) {(a)};
  % 曲线身份直接写在面板里（理想弧在 y>50 Ω 处、CPE 弧在 x>22 Ω 处都已是最上沿，
  %   这两处标注不会与任何曲线交叉）。图例因此取消，避免文字块压住弧顶。
  \node[anchor=west,font={\fontsize{7}{8.5}\selectfont}] at (axis cs:14,58) {实线 $n{=}1$};
  \node[anchor=west,font={\fontsize{7}{8.5}\selectfont}] at (axis cs:14,52) {虚线 $n{=}0.85$};
  % 说明块放进右下角的空白区（两条弧在 x>22 Ω、y<15 Ω 处都没有线），
  %   不再放到坐标区外，避免被裁掉。
  \node[anchor=south east,align=right,text width=28mm,font={\fontsize{7}{8.5}\selectfont}]
    at (axis cs:60.5,14.5)
    {理想弧顶点 $(60,\,50)$\\[1pt]
     CPE 弧顶点 $(60,\,39.4)$\\[1pt]
     $=(R_{\mathrm{ct}}/2)\tan(n\pi/4)$};
\end{axis}

\begin{axis}[
  name=eisdb,
  at={(eisda.right of east)}, anchor=left of west, xshift=6mm,
  width=74mm, height=60mm,
  xlabel={$Z'$ / $\Omega$}, ylabel={$-Z''$ / $\Omega$},
  xmin=0, xmax=300, ymin=0, ymax=260,
  xtick={0,50,100,150,200,250,300}, ytick={0,50,100,150,200,250},
]
  \addplot[mainline,black,densely dotted,domain=0.3:5.4,samples=300,variable=\z]
    ({didx(\z)},{didy(\z)});
  \addplot[mainline,black,domain=-1.7:5.4,samples=320,variable=\z]
    ({dkwx(\z)},{dkwy(\z)});
  % 六个读图点的坐标同样来自上面的函数（log10 f = -1.69897、-1、1.50515、
  %   3.50515、4.60206、4.90309，即 0.02 Hz、0.1 Hz、32 Hz、3.2 kHz、40 kHz、
  %   80 kHz），不是手写读数。
  \addplot[only marks,mark=o,mark size=1.1pt,black,mark options={fill=white}]
    coordinates {({dkwx(-1.69897)},{dkwy(-1.69897)})
                 ({dkwx(-1)},{dkwy(-1)})
                 ({dkwx(1.50515)},{dkwy(1.50515)})
                 ({dkwx(3.50515)},{dkwy(3.50515)})
                 ({dkwx(4.60206)},{dkwy(4.60206)})
                 ({dkwx(4.90309)},{dkwy(4.90309)})};
  % 45° 参考线只画一小段，放在扩散尾下方、说明块左侧的空白处（原先横跨全图，
  %   从三块说明文字上穿过去 —— R3-U01）。
  \addplot[mainline,black,dashed] coordinates {(140,10) (170,40)};
  \node[anchor=north east,font={\fontsize{7}{8.5}\selectfont}]
    at (axis cs:296,256) {(b)};
  % 面板 (b) 里最高的一条曲线只到 112.8 Ω，故 y>150 Ω 的整条带都是空白。
  \node[anchor=north west,align=left,text width=30mm,font={\fontsize{7}{8.5}\selectfont}]
    at (axis cs:8,250)
    {点线：$Z_{\mathrm{W}}{=}0$\\ 实线：$\sigma{=}40$\\ 虚线：$45$° 参考线};
  % 两块说明都落在空白区：上方（y>120 Ω 无曲线）与右下方（扩散尾下方）。
  \node[anchor=north west,align=left,text width=34mm,font={\fontsize{7}{8.5}\selectfont}]
    at (axis cs:8,182)
    {实线与点线在 $f\gtrsim 100$ Hz 几乎重合：只看高频段，\\ 分不开「$Z_{\mathrm{W}}$ 大」与「$R_{\mathrm{ct}}$ 大」};
  \node[anchor=south east,align=right,text width=22mm,font={\fontsize{7}{8.5}\selectfont}]
    at (axis cs:296,6)
    {高频截距 $\approx R_{\Omega}$\\ 只在弧闭合、\\ 无电感时成立};
\end{axis}
\end{tikzpicture}
\caption{三类等效电路的 Nyquist 图对照（模型计算）。取
$R_{\Omega}=10$ \unit{\ohm}、$R_{\mathrm{ct}}=100$ \unit{\ohm}、
$C_{\mathrm{dl}}=0.5$ µF（$\omega_{\max}=1/(R_{\mathrm{ct}}C_{\mathrm{dl}})$
对应 $f_{\max}\approx3.2$ kHz）；CPE 取 $n=0.85$、
$Q=C_{\mathrm{dl}}^{\,n}R_{\mathrm{ct}}^{\,n-1}\approx2.21\times10^{-6}$ S·s$^{0.85}$，
即低频极限电容与理想情形相同；Warburg 取 $\sigma=40$ \unit{\ohm.s^{-1/2}}。
三条曲线（实线、虚线、点线）分别由式~\eqref{eq:eis-d-randles}、
式~\eqref{eq:eis-d-cpe}、式~\eqref{eq:eis-d-rc} 按同一组参数算出：图中两个弧顶标记与
六个空心圆点都是这些函数的取值，不是手工写上的读数。
\textbf{(a)} 只画 $f=2$ Hz 至 79 kHz、$Z'\le62$ \unit{\ohm} 的窗口（两条弧的右半段落在
窗口外）：理想弧（$n=1$）的顶点在
$(R_{\Omega}+R_{\mathrm{ct}}/2,\;R_{\mathrm{ct}}/2)=(60,\,50)$，CPE 弧的顶点在
$(60,\,39.4)$——\textbf{横坐标相同，纵坐标只有 $(R_{\mathrm{ct}}/2)\tan(n\pi/4)$；
同样一个 $R_{\mathrm{ct}}$、同样一个低频极限电容，仅因 $n=0.85$ 就把弧的高度压到
$0.79$ 倍，而顶点频率并不改变}。两条曲线在 $f\to0$ 时都收敛到
$R_{\Omega}+R_{\mathrm{ct}}=110$ \unit{\ohm}（点线在 (b) 中可见 $10\to110$ \unit{\ohm}
的完整弧；CPE 弧的右半段落在 (a) 的窗口之外）。
\textbf{(b)} 频率 0.02 Hz 至 250 kHz：实线在 $f\lesssim 10$ Hz 之后已明显
偏离点线，朝向 $45$°参考线（虚线）——但\textbf{即使在 40 kHz，实线的截距仍是
10.6 \unit{\ohm} 而非 10.0 \unit{\ohm}、相角仍有 $+36.6$°（按 $\varphi\equiv-\arg Z$），即高频平台并未形成}，
这正是规则一第（1）条的量化含义；六个空心圆点依次标出后文读图实例所用的频率
（80 kHz、40 kHz、3.2 kHz、32 Hz、0.1 Hz、0.02 Hz），可以逐点核对。
注意 (b) 两个坐标轴的刻度尺度不同（横轴 300 \unit{\ohm} 与纵轴 260 \unit{\ohm}
画在等宽的图面上），故 $45$°参考线在图上并不呈 $45$°；判斜率请用数值而非图面角度。
局限：本图不含接触电阻、引线电感、有限扩散与吸附弛豫，真实数据在高频端与最低频端都会多出
本图没有的结构。
\modelmark\pedset{$R_{\Omega}$、$R_{\mathrm{ct}}$、$C_{\mathrm{dl}}$、$n$、$\sigma$ 为教学设定，
按式~\eqref{eq:eis-d-randles}、式~\eqref{eq:eis-d-warburg}、式~\eqref{eq:eis-d-cpe} 计算。}
\figsrc{结构取自 S2 §11.5.1 对四类电极 Nyquist 图与原书表 11-1 的论述
（书内 p.859—861 / PDF p.347—349）；具体数值为模型计算，非原书数据。}}
\label{fig:eis-d-model}
\end{figure}

\paragraph{读图实例：在一张模型 Nyquist 图上依次读出三类信息} 下面用图~\ref{fig:eis-d-model}
的数据走一遍完整读图流程。把\textbf{观察}与\textbf{推断}严格分开，因为这里的每一步推断都依赖
一个前提。

\begin{obs}
在图~\ref{fig:eis-d-model}(b) 的实线（含 Warburg）上（六个空心圆点即下列读数，
按频率从高到低排列）：
（1）从最高频端（80 kHz）开始，曲线左端落在 $Z'\approx 10.16$ \unit{\ohm}、
$-Z''\approx 3.97$ \unit{\ohm} 处，并未落在实轴上；
（2）频率降到 40 kHz 时 $Z'\approx 10.63$ \unit{\ohm}、$-Z''\approx 7.91$ \unit{\ohm}：
截距几乎没变，$-Z''$ 却翻了一倍；
（3）频率降到约 3.2 kHz 时，$-Z''$ 达到极大值约 50.1 \unit{\ohm}，此时
$Z'\approx 59.6$ \unit{\ohm}；
（4）继续降频，$-Z''$ 先减小到约 3.9 \unit{\ohm}（$f\approx 32$ Hz，此时
$Z'\approx 112.8$ \unit{\ohm}），再转而增大；
（5）在 $f\lesssim 10$ Hz 之后，$Z'$ 与 $-Z''$ 同步增大，在
0.1～0.02 Hz 段内 $Z'$ 从 160 \unit{\ohm} 增至 223 \unit{\ohm}、
$-Z''$ 从 50 \unit{\ohm} 增至 113 \unit{\ohm}，两点连线的斜率约为 $1.00$；
（6）把 (b) 中的点线（无扩散）与实线对比：两条线在 $f\gtrsim 100$ Hz 时肉眼不可分，
分岔从约 30 Hz 开始，到 0.02 Hz 时点线停在 $Z'=110$ \unit{\ohm}、
$-Z''\approx0$，而实线已走出 223 \unit{\ohm} 之外。
在 (a) 中对比理想弧（实线）与 CPE 弧（虚线）：两条弧的高频端都趋向同一点
（$Z'=R_{\Omega}$），低频端都趋向 $Z'=R_{\Omega}+R_{\mathrm{ct}}$（两端位置与 $n$ 无关），
但中段完全分开——理想弧顶点 $(60,\,50)$，CPE 弧顶点 $(60,\,39.4)$，后者只有前者的
$0.79$ 倍高，而顶点频率同为 $3.2$ kHz。
\end{obs}

\begin{infer}
（1）高频端不落在实轴上、且 $-Z''$ 仍有 4.0 \unit{\ohm}，说明\textbf{最高频处高频平台
尚未形成}，因此 (b) 左端那个 10.16 \unit{\ohm} \textbf{不能直接当作 $R_{\Omega}$}：
照读会把溶液电阻高估约 $1.6\%$；若改用 40 kHz 的读数 10.63 \unit{\ohm}，高估量升到
$6.3\%$。（该值是否可读，判据是 Bode 图上 $|Z|$ 与 $\varphi$ 是否出现平台。）
（2）极化电阻以弧的形式出现，说明该频段由电荷转移与双电层构成的\textbf{并联}支路支配；
在扩散尚未参与的高频段，这段弧的直径给出 $R_{\mathrm{ct}}$：理想弧的直径
$110-10=100$ \unit{\ohm}，与设定值一致。顶点条件
$\omega_{\max}=1/(R_{\mathrm{ct}}C_{\mathrm{dl}})$ 反过来给出
$C_{\mathrm{dl}}=1/(2\pi\times3183\times100)\approx0.5$ µF，也与设定值一致——
\textbf{但这两个数能读对，是因为本例恰好是完整半圆；换成 (a) 中的 CPE 弧，
用「直径」读 $R_{\mathrm{ct}}$ 仍得 100 \unit{\ohm}（两端截距与 $n$ 无关），
而用「弧顶高度 $=R_{\mathrm{ct}}/2$」读就只有 $2\times39.4=79$ \unit{\ohm}，
再代入 $C=1/(\omega_{\max}R_{\mathrm{ct}})$ 得 $0.63$ µF——与真值分别差 $21\%$ 与
$27\%$。压扁的弧不满足式~\eqref{eq:eis-d-rc} 的圆方程：顶点高度是
$(R_{\mathrm{ct}}/2)\tan(n\pi/4)$，不是 $R_{\mathrm{ct}}/2$。}
（3）$-Z''$ 先减到极小再增大，是本图上最容易被忽略的特征：极小值点的出现说明\textbf{有两个
不同时间尺度的过程串接在同一支路上}——高频弧（电荷转移）在约 30 Hz 处已衰减完毕，
而扩散项在更低频处才开始主导。因此「$-Z''$ 极小点」的位置是判断两个过程是否已分离的实用标志。
（4）低频段 $Z'$ 与 $-Z''$ 同步增大且连线斜率约 $1.00$，与式~\eqref{eq:eis-d-warburg}
的 $Z'_{\mathrm{W}}=-Z''_{\mathrm{W}}$ 一致，故该段可归为\textbf{半无限线性扩散（Warburg 阻抗）
控制}。斜率也因此是判断「扩散段是否已经独立」的实用判据：若把拟合区间一直取到弧上
（例如 3.2 kHz 与 0.02 Hz 两点连线，斜率只有约 $0.38$），正在衰减的电荷转移弧会把斜率
大幅拉低——\textbf{用斜率反算 $\sigma$ 必须先扣掉弧的贡献}。
（5）(b) 中点线与实线的分岔起点约 30 Hz，正是「扩散开始参与」的频率；而两条线在
$f\gtrsim 100$ Hz 不可分，说明\textbf{只看高频段的数据无法判断体系有没有扩散}。
这给出频率范围要求的一个定量形式：低频端必须至少比这个分岔频率低两个数量级，扩散尾才成形。
（6）(a) 中两条弧的高频端与低频端重合、中段分开，说明\textbf{弧的「两端」不携带
$n$ 的信息}——端点是 $R_{\Omega}$ 与 $R_{\Omega}+R_{\mathrm{ct}}$，与 $n$ 无关。因此
「用低频截距减高频截距得 $R_{\mathrm{ct}}$」这条做法即使在 CPE 情形下也成立
（本例中两条弧都给出 100 \unit{\ohm}）；\textbf{出错的是把中段几何（顶点高度、半高处宽度）
当作 $R_{\mathrm{ct}}/2$ 与 $1/(R_{\mathrm{ct}}C_{\mathrm{dl}})$ 去用}。
\end{infer}

\begin{pitfall}
\textbf{（1）拟合优度高不等于电路唯一。} 这是等效电路拟合的第一风险，也是本节所有其它风险的
放大器。\textbf{不同的电路可以对同一组数据给出几乎相同的拟合曲线}——只看图~\ref{fig:eis-d-model}(a)
的中段，同一段压扁的弧（跨度 $110-10=100$ \unit{\ohm}、顶点高 $39.4$ \unit{\ohm}）
至少可以用三种完全不同的电路拟合：①CPE（$n=0.85$）加 $R_{\mathrm{ct}}=100$ \unit{\ohm}；
②两个时间常数相差 $4$ 倍的 $(R\parallel C)$ 串联（各 $50$ \unit{\ohm}，即把这一段
100 \unit{\ohm} 拆成 $50+50$——峰值仍在 $3.2$ kHz、峰高 $40.0$ \unit{\ohm} 与
$39.4$ \unit{\ohm} 几乎相同，全频段与①的差别不超过约 $3$ \unit{\ohm}）；
③单个理想 $(R\parallel C)$ 串一个扩散项。三者的残差都可以做到很小，而它们给出的
「$R_{\mathrm{ct}}$」可以是 100 \unit{\ohm}（①）、两个各 50 \unit{\ohm}（②）、
或按弧顶高度误读出的 79 \unit{\ohm}——相差一倍以上。
\textbf{残差小只说明该电路在数值上能复现数据，不代表它对应真实的物理过程}。实践判据：换一组
完全不同的初值（或换一个拓扑相近的电路）重新拟合，若参数漂移到另一个解而残差几乎不变，就说明
该解不可辨识；此时只能报告「表观值」并把候选机理并列写出。

\textbf{（2）参数相关性使单个参数失去意义。} 这是风险（1）的数学根源。CPE 的 $Q$ 与 $n$ 高度
相关（$n$ 略变，$Q$ 可补偿性地变几个数量级）；$R_{\Omega}$ 与 $R_{\mathrm{ct}}$ 在同一数量级时
互相补偿；两个时间常数接近时，$R_{1}$、$C_{1}$、$R_{2}$、$C_{2}$ 四者其实只是两个可辨识的
组合量。因此\textbf{必须报告参数的置信区间与相关系数矩阵}——只报一组「最优值」而不报不确定度，
等于没有报告结果。报告 $Q$ 与 $n$ 时必须同时给出换算成电容所用的式子
（式~\eqref{eq:eis-d-cpeceff}），否则读者无法复算。

\textbf{（3）把 Warburg 当半圆。} 半无限扩散的 $45$°直线在只测到一段时很像大半径圆弧的一部分；
图~\ref{fig:eis-d-model}(b) 的实线在 10 Hz 以下那一段若单独放大，斜率 $1.00$ 的
「直线」完全可以被一个半径很大的圆拟合。直接按半圆读直径会把 $R_{\mathrm{ct}}$ 高估几倍。
\textbf{判别办法是换到 Bode 图}：扩散段的 $|Z|$ 应呈 $\omega^{-1/2}$、相角应稳定在约 $+45$°，
而圆弧段的相角会扫过一整段。

\textbf{（4）频率范围不够就下结论。} 图~\ref{fig:eis-d-model}(b) 已经给出量化例子：两条曲线
（有扩散／无扩散）在 $f\gtrsim 100$ Hz 不可分，在 40 kHz 处连高频平台都还没
形成。起始频率不够高，截距含 $R_{\mathrm{ct}}$，$R_{\Omega}$ 偏大；终止频率不够低，看不到扩散
尾或第二个弧，于是把「没看到」当成「不存在」，进而误判速控步。\textbf{频率范围不是实验参数，
它是结论的适用范围}，必须与结论一起报告。

\textbf{（5）把 CPE 当理想电容算电容值。} 这是最常见的一个算术错误，且\textbf{不报错、不告警}。
$Q$ 的量纲是 $\mathrm{S.s^{n}}$，$n=0.85$ 时 $Q=2.21\times10^{-6}$ 与同体系的
$C_{\mathrm{dl}}=0.5$ µF 相差 $4.4$ 倍，而且量纲不同——直接拿 $Q$ 当作电容会高估四倍以上；
式~\eqref{eq:eis-d-cpeceff} 的换算是必须做的。另一个同源错误是把
$C_{\mathrm{dl}}=1/(\omega_{\max}R_{\mathrm{ct}})$ 直接用在压扁的弧上——如图~\ref{fig:eis-d-model}(a)
所示，把弧顶高度当 $R_{\mathrm{ct}}/2$ 会读成 $R_{\mathrm{ct}}\approx79$ \unit{\ohm}，
代入后得 $0.63$ µF（真值 0.5 µF，高 $27\%$）；压得越扁，误差越大。
\textbf{只要 $n<0.95$，就不要用几何法读电容。}

\textbf{（6）未扣除溶液电阻。} $R_{\mathrm{ct}}$ 必须是在扣除 $R_{\Omega}$ 之后读出的弧直径；
用未校正的过电位去算动力学参数，误差随电流密度线性放大。更隐蔽的是：$R_{\Omega}$ 随电解质
浓度、电极间距与温度变化，\textbf{每次实验都应重新读，不能沿用上一批的值}。

\textbf{（7）把接触阻抗当电荷转移阻抗。} 电极与集流体、催化剂层与膜之间的接触电阻，以及参比
电极的鲁金毛细管位置不当引起的未补偿电阻，都会并入高频段，被拟合成一个大 $R_{\mathrm{ct}}$ 或
一个低 $n$ 的 CPE。这与（1）的危险叠加：接触电阻造成的低 $n$ 会被解释成「表面不均匀」这一
化学结论。\textbf{判别办法：改变电极压紧力或参比位置，若该「$R_{\mathrm{ct}}$」或「$n$」跟着
变，它就不是化学信息}。

\textbf{（8）稳态未建立就测。} 阻抗定义要求体系在测量期间是线性时不变系统。刚跳到目标电位、
表面覆盖度还在漂移、或催化剂正在失活时测得的谱，会在低频段出现「漂移型」假弧——而低频段恰好
是耗时最长、也最容易被漂移污染的一段。判据见「测量条件」第（5）条：先在单一频率上监测
$\lvert Z\rvert$ 的时间稳定性，再做全谱扫频；同一电位下正扫与反扫各测一次，两条谱重合才说明
稳态已建立。

\textbf{（9）扰动幅度过大。} 幅度过大会使响应非线性，Kramers--Kronig 检验不通过，并系统性
压低 $R_{\mathrm{ct}}$——因为大信号下电流–电势关系不再是线性的，等效电路这一整套线性工具
（包括式~\eqref{eq:eis-d-randles}）全部失效。用幅度减半检验把关。

\textbf{（10）把 $f_{\max}$ 当作 $\omega_{\max}$ 代入。} 式~\eqref{eq:eis-d-cdl} 中的
$\omega$ 是角频率。用 Hz 为单位直接代入会得到错 $2\pi$ 倍（约 $6.28$ 倍）的电容值。这是一个
纯粹的记账错误，但因为结果「看起来合理」，很少被当作可疑值复查。
\end{pitfall}

\begin{crosscheck}
阻抗谱给出的归属结论不应单独成立，至少与下列方法之一互查——每条互查针对的是\textbf{同一组
数据的一个具体参数}，而不是笼统地「再测一次」：

\textbf{（1）$R_{\mathrm{ct}}$ 对稳态极化／Tafel 斜率。} 在小过电位下
$R_{\mathrm{ct}}=RT/(nFi_0)$，因此由弧直径反推的 $i_0$ 应与 Tafel 外推值同数量级；差一个数量级
以上就说明弧里混了别的贡献（接触电阻、扩散、或第二个时间常数）。
\tocheck{未对照原页面图像：$R_{\mathrm{ct}}$ 与 $i_0$ 的换算关系式及电子数 $n$ 的约定按学科
标准形式给出，源文本 §11.5.1 未给出该式，须与页面核对。}

\textbf{（2）扩散项对旋转圆盘变转速。} 扩散控制的阻抗应在提高转速后缩短低频尾（边界层变薄），
而电荷转移控制的弧几乎不变。这条互查把 Warburg 与半圆分开，且不依赖任何模型假设——是本节最
有力的一条独立判据。RDE 的 Levich 斜率与阻抗 $\sigma$ 还可交叉校核同一个 $D$
（\S 11.4.2）。

\textbf{（3）$C_{\mathrm{dl}}$ 对循环伏安的电容电流。} 在无法拉第电流的电位区间做不同扫速的
循环伏安，由 $\mathrm{d}i/\mathrm{d}v$ 得到电容，应与阻抗高频弧给出的电容一致；两者一致才说明
高频弧确由双电层充电支配。

\textbf{（4）$R_{\Omega}$ 对电流中断法或正反馈补偿。} $R_{\Omega}$ 是阻抗谱最不可靠的一个数
（见「误判」第 6、7 条），必须用独立方法标定。三个值一致时，高频截距的归属才算落实。

\textbf{（5）时间常数个数对机理中的独立状态变量个数。} 源文本强调，从阻抗谱中时间常数的个数
及其数值大小可以获得各个步骤的动力学信息与电极表面状态变化的信息，还可观察电极过程中有无传质
过程的影响\src{书内 p.856 / PDF p.344}。把这条当作系统性的自检：\textbf{Bode 相角图上数出的
峰个数，应与所假定机理中独立状态变量的个数一致}；多了说明电路里塞进了不存在的元件，少了说明
两个过程的时间常数重叠、当前电路无法分辨它们。

\textbf{（6）与其它分片的一致性。} 本节是同一章内给出阻抗基本定义与 $Z$ 与元件关系的地方，
因此它与下一节（§11.5，电化学阻抗谱的应用与读图）的接口必须一致：\textbf{（i）}$R_{\Omega}$
的符号与含义（分片 c 沿用源文本符号 $R_{\mathrm{s}}$，指同一量）；\textbf{（ii）}Nyquist 图
纵轴取 $-Z''$ 的约定；\textbf{（iii）}相角取测量约定 $\varphi\equiv-\arg Z$（容性为正，纯电容
$+90$°、感性为负——分片 c 的读图实例即按此约定写「相角接近 $90$°」与「第四象限的负相角」）；
\textbf{（iv）}$\omega_{\max}$ 而非 $f_{\max}$ 用于
式~\eqref{eq:eis-d-cdl}。四处中任何一处不一致，读图结论就无法互相引用。

装配顺序已在 \texttt{chapters/ch11.tex} 中按文件头注释的处置办法（A）修正：本分片现排在
§11.4（分片 b）之后、§11.5 与 §11.6（分片 c）之前，故上文的「读图规则」在下文中确为
\textbf{上文}，§11.5.1 与 §11.5 的先后关系与实际排版一致。两分片之间原有一批重复段落
（阻抗定义与两种画法、等效元件表、三条读数规则、误判清单、测量条件），合并时按
\textbf{保留有源文本页码依据的一侧}处理：这些内容一律以本分片为准，分片 c 只保留
源文本 §11.5.2 特有的读图实例（甲醇氧化）与互补回查，不再重复。
\end{crosscheck}
}{}
\IfFileExists{chapters/parts/ch11-c.tex}{\subsection{电化学阻抗谱}

稳态极化曲线把所有步骤的过电位加总成一个数；阻抗谱则用一串不同频率的小幅正弦扰动，把总响应按\textbf{时间尺度}拆开——快步骤（溶液电阻、双电层充电）在高频段显现，慢步骤（电荷转移、吸附、扩散）在低频段显现\src{S2 §11.5，书内 p.866 / PDF p.354}。分段能力是它最大的价值，也是它最大的风险来源：把频率轴上某一段的形状直接翻译成一个元件，很容易得到「拟合优度很高、物理归属却错」的电路。

本分片源文本自书内 p.866（PDF p.354）起，只包含 §11.5 的后半（应用与读图）与 §11.6 全节；阻抗的基本定义、等效电路推导与 Nyquist／Bode 图的定义在源书 §11.5.1，不在本分片的源文本范围内——它们在前一节（§11.5 阻抗谱基本原理，依源文本写成，标有 \src{} 依据）。下文凡标 \tocheck{} 处，即按学科标准形式重建、未能与源文本对齐的部分。

\begin{probout}
 已知：电极过程的总电流至少含溶液电阻、双电层充电、电荷转移、吸附中间体与扩散五类贡献，稳态方法只能给出它们的总和。
 要问：能否把总响应分解成若干时间常数，分别读出溶液电阻 $R_{\mathrm{s}}$、双电层电容 $C_{\mathrm{dl}}$、电荷转移电阻 $R_{\mathrm{ct}}$ 与扩散阻抗，并据此判断哪一步是速率控制步骤？
 要给出：覆盖足够宽频率范围的一组阻抗数据、一个（或几个）能同时拟合全部频率的等效电路，以及每个参数的物理归属与成立条件——并明确「能拟合」不等于「唯一」。
\end{probout}

\subsubsection{读图实例：甲醇氧化的阻抗谱随电位的变化}

上一节（源书节号 §11.5.1，本项目中即紧接在前的一节）已给出阻抗的定义 $Z(\omega)=E(\omega)/I(\omega)$、Nyquist／Bode 两种画法、等效元件的来源与四类电极的特征谱，以及「高频截距」「半圆直径」「低频 45°斜线」三条读数规则各自成立的条件；本节不再重复，只处理源文本 §11.5.2 的真实数据。

源文本用同一组数据的两套图讨论甲醇在铂系催化剂上的氧化：Nyquist 形式的图 11-42 与以相角为主的图 11-43\src{S2 §11.5.2，书内 p.866 / PDF p.354}。下面把\textbf{观察}与\textbf{推断}严格分开，因为这里的每一步推断都依赖一个前提。

\begin{obs}
随直流电位升高，阻抗谱按四个区间改变形状：
（1）$0.10\sim0.20$ V：谱几乎是一条沿虚轴、相角接近 $90$° 的直线；
（2）$0.20\sim0.33$ V：相角迅速下降，但仍在 $0$° 以上，低频区近似电阻行为；谱线略偏离虚轴；
（3）$0.35$ V 起：谱延伸进第四象限，出现负相角（电感现象）；在 $0.35\sim0.48$ V 区间，电感行为位于低频区，且各条曲线随电位升高整体向高频方向移动；
（4）$0.48$ V 以上：谱突变，分裂为两段分离的曲线，高频段相角越出 $90$°。
但图中纵横坐标与相角的具体数值需对照原书图 11-42、图 11-43\tocheck{未对照原页面图像：源文本对 0.48 V 以上的描述两次出现「超过 90°」（高频段与低频段），疑为 OCR 重复或原文指低频段相角亦越出 0—90° 区间；图的纵轴定义（$-Z''$ 或相角／°）与频率范围无法在本分片源文本内确认}。
\end{obs}

\begin{infer}
（1）接近虚轴且相角接近 $90$° 的直线是理想极化电极（纯电容）特征：在这个电位区间几乎没有法拉第电流，电极／溶液界面只表现为双电层充电。
（2）相角下降而不越过 $0$°，说明出现了耗散性的法拉第通道；谱线开始偏离虚轴，与甲醇离解吸附的开始相对应。
（3）第四象限的负相角是电感行为，源文本把它归属于强吸附中间体 \ce{CO_{ad}} 的吸附弛豫。判据是\textbf{出现电位区间}与已知的 CO 吸附电位区间一致，而不是单看形状——仪器电感也会给第四象限，见下文误判清单第（1）条「第四象限的电感弧」。
（4）源文本采用的判断规则是：\textbf{低频区的过程是速率控制步骤}；因此随电位升高曲线向高频方向移动，就意味着该低频过程的速率在加快。据此，$0.35$ V 以前低频区为电阻行为、速控步是甲醇的离解吸附；$0.35\sim0.48$ V 低频区为电感行为、速控步是生成 \ce{CO_{ad}} 的步骤；$0.48$ V 以上低频区转为钝化行为、速控步是 \ce{CO_{ad} -> CO2}。谱在 $0.48$ V 的突变因此被解释为「CO 的生成与消耗两个步骤中，速控步发生了转移」。
（5）与上一条互补的判据是\textbf{斜率的数值}：$-Z''$ 对 $Z'$ 的斜率明显小于 $45$°（例如落在 $20$°～$30$°）时，通常意味着扩散层已「看到」边界、不再满足半无限扩散的前提，或有吸附弛豫叠加在扩散尾上——此时不能把该段直接当 Warburg 处理。这条判据只在频率范围够低、确实取到了扩散段之后才有意义。
\end{infer}

\begin{dataflow}
从一组阻抗数据到「速控步是谁」的完整链条中，\textbf{检验}（正反扫重合、Kramers--Kronig 残差）与\textbf{数时间常数}（Bode 相角图数峰）两步已写在上一节（源书 §11.5.1），此处不重复；真实数据特有的三步是：
（1）\textbf{定元件归属}——高频截距给 $R_{\mathrm{s}}$，最低频的极限行为决定是扩散（45°）、吸附（低频电容／电感弧）还是钝化（低频大电阻）。
（2）\textbf{跨电位比较}——只有同一套电路、同一套初值在全电位区间内都收敛，才能用「曲线向高频移动」这类趋势论证速率变化。
（3）\textbf{互查}——用独立实验（Tafel 斜率、变转速、原位谱学的覆盖度）验证归属，清单见下文；其中 Tafel 斜率与变转速两条属通用判据，已在上一节的互查清单里连同判据给出。
\end{dataflow}

\subsubsection{拟合的最大风险：参数相关性，以及「拟合好 $\neq$ 电路唯一」}

等效电路拟合的第一风险是\textbf{拟合优度高不等于电路唯一}：不同的电路可以对同一组数据给出几乎相同的拟合曲线。因此逐条列出误判清单仍有必要，尤其是下面几条在真实数据里最容易发生。判据、公式与量化例子见上一节的「误判」清单（十条），此处只列真实数据读图中另有表现的三条：

\paragraph{第四象限的电感弧}引线电感与强吸附中间体（如 \ce{CO_{ad}}）的吸附弛豫都给第四象限的弧。判别办法是改变接线方式或电极几何：\textbf{仪器电感应与频率无关、且能靠改接线消掉}；化学吸附弛豫则出现在确定的电位区间（如读图实例中的 $0.35\sim0.48$ V）。
\paragraph{吸附弛豫与双电层的重叠}吸附中间体的 $R_{\mathrm{ad}}$、$C_{\mathrm{ad}}$ 与 $C_{\mathrm{dl}}$ 的时间常数接近时，二者在 Nyquist 图上不可分辨，只能合并成一个压扁的弧（如读图实例中的甲醇体系）。此时单靠阻抗谱无法把覆盖度信息与双电层信息分开，必须引入原位谱学的覆盖度数据。
\paragraph{薄层与原位谱学联用时的第三套条件}若阻抗数据要与同一电位点的原位红外／拉曼对照，薄层电解池把传质条件改掉了，阻抗谱给出的扩散尾与厚层体系并不相同；两套数据的「同一电位」只有在覆盖度都达到稳态时才是同一个表面状态。

\subsubsection{测量条件与互补回查}

\begin{measure}
 \textbf{稳态与幅度}：体系在目标电位达到稳态后再扫频（判据见上一节测量条件第（5）条）；交流幅度取原书图 11-42／11-43 图注的 AC $10$ mV，并做第（2）条的幅度减半检验。
 \textbf{频率范围}：本分片读图实例所用原书图 11-42／11-43 的图注为 $20$ mHz 至 $5$ kHz；通用范围与点数的分配原则见上一节的测量条件第（3）条与规则三（源书 §11.5.1：低频端的点数要按「能否分辨出 45° 斜率」来定，而不是按对数均匀取点）。
 \textbf{电位区间}：至少覆盖目标反应从「无法拉第电流」到「扩散／钝化控制」的完整区间，才能用趋势论证速控步的转移（如读图实例中的 $0.10\sim0.48$ V）。
 \textbf{测量模式}：低阻抗、大电流体系（燃料电池单池、电池）宜用恒电流模式，避免恒电位仪过载；高阻抗体系（修饰电极、半导体电极）宜用恒电位模式\supp{恒电位／恒电流两种模式的适用判据属 §11.5.2 的讨论范围，不在本分片源文本内，此处只给判据}。
 \textbf{谱学联用前的采样条件}：若要与原位红外／拉曼的同一电位点对照，必须先让电位在该点稳定足够长时间（覆盖度达到稳态），否则两套数据的「同一电位」并不是同一个表面状态。
\end{measure}

\begin{crosscheck}
 阻抗谱的归属结论不应单独成立，至少与下列方法之一互查。上一节的互查清单（六条：$R_{\mathrm{ct}}$ 对 Tafel 斜率、扩散项对旋转圆盘变转速、$C_{\mathrm{dl}}$ 对循环伏安电容电流、$R_{\Omega}$ 对电流中断法、时间常数个数对独立状态变量个数、以及与其它分片的一致性）中，前两条已在上一节连同判据一起给出，此处不重复；真实数据读图另有表现的是三条：
 \textbf{（1）变温}：$R_{\mathrm{ct}}$ 的 Arrhenius 图给出表观活化能，可与同反应在稳态方法下的活化能比较。
 \textbf{（2）原位谱学}：第四象限的电感行为、低频吸附电容，应与原位红外／拉曼给出的 \ce{CO_{ad}} 覆盖度在同一电位区间出现——形状与分子信息对上了，归属才算落实。
 \textbf{（3）时间常数个数}：源文本强调，从阻抗谱中含有的时间常数个数及其数值大小可以推测电极过程的状态变量的情况，也可以观察电极过程中有无传质过程的影响，并由此得到溶液电阻、双电层电容及电极反应电阻等信息、推测反应机理\src{S2 §11.5，书内 p.866 / PDF p.354}（与上一节所引源书 p.856 处是同一条判据的两个出处）。把这条当作系统性的自检：时间常数个数应与你假定的机理中的独立状态变量个数一致，多了或少了都说明电路选错。
\end{crosscheck}

\subsection{电化学原位谱学技术}

常规电化学方法测的是界面上电流对电势或时间的响应，得到的是界面反应的综合信息，无法在分子水平上验证反应物、中间物种与产物，这给机理验证造成根本困难。原位谱学把谱学技术与电化学方法\textbf{同步}结合：在得到电流、电势等电化学信息的同时获得分子水平的信息，用于研究电极过程机理、电极表面特性、检测中间物种与产物、测定转移电子数、确定速率常数与扩散系数\src{S2 §11.6，书内 p.866—867 / PDF p.354—355}。本节按源文本的四个方面介绍：原位红外、原位拉曼、原位质谱与荧光单分子／单纳米粒子技术。

\begin{principle}
 原位（in situ）与离场（ex situ）的分界只有一条：\textbf{谱学测量是否在电化学调制进行的同时完成}。原位红外光谱电化学在反应体系被电化学调制与检测的同时监测红外信号变化；非原位红外则是在终止电化学反应之后对产物做红外研究\src{S2 §11.6.1，书内 p.867 / PDF p.355}。这一条分界决定了信息性质：原位方法能看到中间体与电位依赖，但必须解决「电解液对信号的干扰」与「传质受限」；离场方法信噪比与方法选择自由度大得多，但可能错过只在电位存在时才稳定的物种。
\end{principle}

\subsubsection{原位红外光谱电化学}

红外光谱电化学（infrared spectroscopy in electrochemistry）由红外光谱与电化学技术结合而成，可分为原位与非原位两类。原位红外既能检测电极表面吸附物种的身份与分子结构，也能探测电极表面及贴近表面薄液层中的反应物、产物与中间体，并为固／液界面提供双电层中物种的化学本质、成键方式及其随电位变化的信息。$1980$ 年 A. Bewick 等首次把电化学原位红外光谱成功用于电极／溶液界面研究，此后已建立的方法包括电化学调制红外光谱法（EMIRS）、差减归一化界面傅里叶变换红外光谱法（SNIFTIRS）、线性电位扫描反射光谱法（LPSIRS）、偏振调制红外光谱法（PMIRS）、傅里叶变换红外反射吸收光谱法（FTIRRAS）、时间分辨光谱电化学法（TRFTIRS）、二维红外光谱（2D-IR）与显微红外光谱电化学法（MFTIRS）等\src{S2 §11.6.1，书内 p.867 / PDF p.355}。

\begin{measure}
 \textbf{光路与窗口}：原位反射红外最常用外反射模式，红外辐射穿过窗口与溶液薄层，经研究电极表面反射后被检测。窗口材料必须对红外透明且不溶于所测溶液，如 \ce{CaF2}、\ce{Si}、\ce{ZnSe}。
 \textbf{薄层厚度是关键条件}：水及其他常用溶剂对红外辐射有强烈吸收，因此窗口与电极表面之间的液层必须很薄，约 $1\sim100$ nm。这一厚度同时决定了原位红外的时间响应与传质条件——薄层抑制了扩散，使本体浓度不再等于表面浓度，是原位红外定量结论的主要限制。
 \textbf{电极与池结构}：研究电极是封在玻璃或绝缘惰性材料中的圆盘，表面平整度要求高，固定在可调节电极与窗口距离的活塞末端，实验时充分接近窗口；通过调整反射镜角度选择最佳红外入射角，再用调制或差分技术取得最大检测信号\src{图 11-44、图 11-45，书内 p.868 / PDF p.356}。
\end{measure}

三种主要实验测试方式的差别在「减掉什么」上，理解这一点比记住缩写更重要\src{S2 §11.6.1，书内 p.868—870 / PDF p.356—358}：

\paragraph{EMIRS（电化学调制红外光谱）}
电势在「无研究物质产生」与「有电化学产物」两个位置之间以一定频率调制，用相敏检测器（PSD）解调，因而可以在测产物信号的同时排除溶剂与其他溶解物的吸收干扰，提高信噪比。测量的是电极电势由 $E_{1}$ 变为 $E_{2}$ 时表面相对反射率的改变量 $\Delta R$（相当于 $R(E_{2})-R(E_{1})$），最后以 $\Delta R/R$ 对波数作图。要点：\textbf{得到的是两个电位下的差谱}，因此单张谱不能直接给出绝对覆盖度。

\paragraph{SNIFTIRS（差减归一化界面傅里叶变换红外光谱）}
在傅里叶变换红外光谱基础上发展而来：单色仪被干涉仪取代，采集到的是强度随干涉仪物镜位置变化的干涉图；优点是可以多次采集并累加干涉图以提高信噪比。它测量两个不同电势下的谱图信号（电势分别取在所测物种的电化学过程发生与不发生处），在相应波长处作吸光度差减，再由计算机处理得到差谱。与 EMIRS 相比，它牺牲了调制解调的实时性，换来了可在宽波数范围同时取谱与更好的信噪比。

\paragraph{IRRAS（红外反射吸收光谱）}
依据红外反射光谱的表面选律：只有当吸附分子的偶极矩在电极表面法线方向上的分量不为零时，才能吸收平行于入射面的红外能量。装置上用光弹性调制器把入射光的偏振在 P 偏振与 S 偏振之间调制：P 光平行于入射面，可被电极表面与环境吸收；S 光垂直于入射面，不能与表面吸附物种相互作用，只包含环境对红外光的吸收。因此 P 光与 S 光的反射光谱之差就给出表面粒子的红外吸收信息。注意 IRRAS 在取谱时\textbf{电势保持不变}，得到的结果只代表该固定电势下的表面层信息，不同电位之间要各自成谱再比较。

\subsubsection{原位拉曼光谱}

\begin{principle}
 原位拉曼是在进行电化学过程的同时，用拉曼光谱研究电化学池中电极／溶液界面的性质：在测电位、电流获得电信号的同时，把一定波长的入射光照射在电极表面，再分析反射、透射或散射光束\src{S2 §11.6.2，书内 p.870 / PDF p.358}。
 散射的物理图像：光照射到分子或离子上，与分子发生碰撞并交换能量，光子的运动方向与能量随之改变。仅改变方向的弹性碰撞给出瑞利（Rayleigh）散射；方向与能量同时改变的非弹性碰撞给出拉曼（Raman）散射。散射光能量大于入射光能量的谱线（频率更高）称反斯托克斯线（紫伴线），小于入射光的称斯托克斯线（红伴线）\src{图 11-49，书内 p.871 / PDF p.359}。$1928$ 年 Raman 发现穿过透明介质的光频率发生改变，$1930$ 年获 Nobel 物理学奖。
 装置由拉曼光谱仪与原位电化学拉曼池两部分组成。光谱仪含激光源、收集系统（透镜组）、分光系统（光栅，或陷波滤光片结合光栅以滤除瑞利散射与杂散光）与检测系统（光电倍增管、半导体阵列检测器或多通道电荷耦合器件）。拉曼池具备工作电极、辅助电极、参比电极与通气装置，且必须配备光学窗口的密封体系，避免腐蚀性溶液与气体侵蚀仪器\src{S2 §11.6.2，书内 p.870—871 / PDF p.358—359}。
\end{principle}

\begin{measure}
 薄层条件是原位拉曼与原位红外共有的约束，但尺度不同：条件允许时应采用薄层溶液，电极与窗口间距取 $0.1\sim1$ mm，以尽量避免溶液信号的干扰；这一点对显微拉曼系统尤其重要，光学窗片或溶液层太厚会改变显微系统的光路，使表面拉曼信号的收集效率下降\src{S2 §11.6.2，书内 p.871 / PDF p.359}。
\end{measure}

\paragraph{应用一：CO 在贵金属表面的吸附}
CO 等小分子吸附在贵金属上时可视为偶极子，偶极子之间存在很强的相互作用（偶极耦合），导致 C--O 与 metal--CO 的拉曼频率蓝移；一般认为 metal--CO 的振动峰强于 C--O 峰。红外强度可代表金属上吸附 CO 的平均信号，而 SERS 研究发现了与红外类似的分压依赖行为，因此可以认为 SERS 信号能够代表纳米粒子上吸附 CO 分子的平均信号\src{S2 §11.6.2，书内 p.872 / PDF p.360}。这条结论是后面所有 SERS 定量讨论的前提，它的成立条件是「SERS 信号与分子数成正比」，而不是「SERS 峰强可比」。

\paragraph{应用二：酸性溶液中甲醇在 Pt 上的吸附与氧化}
在 $0.5$ mol/L \ce{CH3OH} + $0.5$ mol/L \ce{H2SO4} 溶液中扫甲醇氧化的循环伏安曲线并结合原位拉曼，可以得到甲醇在 Pt 上吸附氧化的完整图像\src{S2 §11.6.2，书内 p.872 / PDF p.360}：
\begin{obs}
 正扫至 $0.5$ V 之前，甲醇电氧化电流很小，可观察到 \qty{490}{cm^{-1}}、\qty{410}{cm^{-1}}、\qty{1060}{cm^{-1}} 及 \qty{1850}{cm^{-1}} 四个谱带；正扫过 $0.5$ V 后四个谱带明显减弱、氧化电流明显增大；扫至 $0.8$ V 时 CO 相关谱带几乎完全消失；扫至 $1.1$ V 时在 \qty{580}{cm^{-1}} 附近出现一个谱带。负扫时 \qty{580}{cm^{-1}} 谱带逐渐减弱，扫至 $0.8$ V 时消失，反扫的氧化电流出现峰值；负扫至 $0.5$ V 以下时 \qty{490}{cm^{-1}} 与 \qty{1060}{cm^{-1}} 谱带又逐渐增强。
\end{obs}
\begin{infer}
 四个谱带分别对应 Pt--CO$_{\mathrm{L}}$、Pt--CO$_{\mathrm{M}}$、C--O$_{\mathrm{L}}$ 与 C--O$_{\mathrm{M}}$ 的伸缩振动，说明该电位下电极表面 Pt 处于 CO「中毒」状态；谱带减弱与电流增大同步，说明电势扫至 $0.5$ V 时表面上电化学生成的 CO 开始被氧化；$0.8$ V 时 CO 谱带消失说明 CO 已完全氧化。$1.1$ V 出现的 \qty{580}{cm^{-1}} 谱带来源于 Pt 的多种氧化物种的 Pt--O 振动；负扫时它减弱并消失，说明氧化的 Pt 被还原、活性位增加，因此反扫氧化电流出现峰值。负扫至 $0.5$ V 以下时 Pt--CO 与 C--O 谱带重新增强，说明低电势下甲醇重新吸附于电极表面、脱氢后重新生成 CO；而 CO 在较低电势下无法被氧化成 \ce{CO2}，于是电极被 CO 毒化，阻碍甲醇进一步脱氢，氧化电流下降。
\end{infer}
这里\textbf{观察与推断的接口是谱带归属}：四个谱带（\qty{490}{cm^{-1}}、\qty{410}{cm^{-1}}、\qty{1060}{cm^{-1}}、\qty{1850}{cm^{-1}}）与 Pt--CO$_{\mathrm{L}}$、Pt--CO$_{\mathrm{M}}$、C--O$_{\mathrm{L}}$、C--O$_{\mathrm{M}}$ 伸缩振动的对应关系来自文献与标定实验，一旦归属错位，整段机理推断都会跟着错。因此这类工作的可信度取决于是否在同一体系里做了归属的独立验证（如同位素取代、参比化合物或理论计算）。

\begin{pitfall}
 \textbf{（1）拉曼信号弱是方法的固有短板。} 拉曼散射光子数极低——源文本给出的量级是荧光的 \num{1e-14}、红外的 \num{1e-6}，如此低强度的散射光用于研究表面或界面分子的结构与性质十分困难\src{S2 §11.6.2，书内 p.872—873 / PDF p.360—361}。常用对策是表面增强拉曼光谱（SERS）：把试样吸附于胶态金属（Cu、Ag、Au）或这些金属的粗糙表面上，再用普通拉曼方法测量，强度可达正常情况的 \num{1e5}～\num{1e6} 倍，因而可用于极低浓度的样品溶液，并能有效避免溶液中相同物种的信号干扰。
 \textbf{（2）SERS 的增强是条件依赖的，不是常数。} 增强因子取决于基底形貌、金属种类、激发波长与吸附构型，不同批次基底之间不可直接比较峰强。用 SERS 峰强做定量时，必须同批基底、原位标定；把 \num{1e5}～\num{1e6} 当作一个固定倍数去反算表面浓度是最常见的误用。
 \textbf{（3）薄层条件本身会改变被测体系。} 电极与窗口间距 $0.1\sim1$ mm 的薄层抑制了自然对流与本体扩散，使电极表面浓度低于本体；因此薄层谱学给出的电位—覆盖度关系不能直接外推到流动或厚层条件下的动力学。
 \textbf{（4）研究条件与实际工作条件差别太大。} 燃料电池基础研究一般在室温下用三电极体系，而组装成电堆后为提高反应速率常采用高温运行；考虑到 Nafion 膜的稳定性与活性以及三相水界面的要求，质子交换膜燃料电池的使用温度一般在 $80$ \degC 左右。温度不仅按 Arrhenius 关系改变反应速率，也会改变反应途径——甲醇、甲酸的直接途径与经 CO 的间接途径，氧还原的 2 电子与 4 电子过程，不同途径给出的电化学性能不同。若研究条件与实际情况相差太大，就很难由研究条件下的表征结果推断实际工作条件下的行为，因此有必要研制变温流动光谱电解池来模拟燃料电池的实际工作状态\src{S2 §11.6.2，书内 p.873 / PDF p.361}。
\end{pitfall}

\subsubsection{原位质谱}

质谱与电化学反应装置联用最早出现在 $20$ 世纪 $70$ 年代。$1971$ 年 Bruckenstein 和 Gadde 用气相质谱检测电化学反应产生的挥发性物质，率先建立电化学质谱技术（EMS）：把 Pt 工作电极制备在多孔玻璃膜上，采用膜进样方式，用多孔 Teflon 膜把电解液与质谱真空进样系统分隔开，只允许挥发性气体产物进入电离室；其特点是气体样品累积到一定量后一次性导入质谱作定性与定量分析，测量电极上气体产物的响应时间约 $20$ s。$1984$ 年 Wolter 和 Heitbaum 建立微分电化学质谱（DEMS），沿用多孔 Teflon 膜进样，区别在于用两级真空泵分别对电离室与质量分析系统分级抽真空，进样口利用压差把反应产生的气体快速吸入电离室，得到的质谱信号强度正比于电化学反应的法拉第电流；DEMS 实现了对挥发性产物的\textbf{连续}测量，总响应时间小于 $1$ s\src{S2 §11.6.3，书内 p.873—874 / PDF p.361—362}。

\begin{principle}
 DEMS 系统由电化学反应装置、膜接口与质谱仪三部分组成，膜接口是核心设计，实现液相与真空的连接：产物在膜两侧压差作用下从液相挥发进入真空系统，采样耗时极短，因而可在线检测。其定量基础是「离子电流正比于法拉第电流」，推导链条为\src{原书式 11-126—11-129，书内 p.874 / PDF p.362}：
 \begin{align}
  I_{\mathrm{F}}&=\frac{\mathrm{d}Q_{\mathrm{F}}}{\mathrm{d}t}=\frac{\mathrm{d}(zFn)}{\mathrm{d}t}=zF\frac{\mathrm{d}n}{\mathrm{d}t}
  \label{eq:dems-if}\\
  n_{\mathrm{i}}&=K_{1}K_{2}n
  \label{eq:dems-ni}\\
  Q_{\mathrm{i}}&=K_{3}n_{\mathrm{i}}
  \label{eq:dems-qi}\\
  I_{\mathrm{i}}&=\frac{\mathrm{d}Q_{\mathrm{i}}}{\mathrm{d}t}
   =K_{1}K_{2}K_{3}\frac{\mathrm{d}n}{\mathrm{d}t}
   =\frac{K_{1}K_{2}K_{3}}{zF}I_{\mathrm{F}}
  \label{eq:dems-ii}
 \end{align}
 令 $K=K_{1}K_{2}K_{3}/(zF)$，即得 $I_{\mathrm{i}}=KI_{\mathrm{F}}$；若法拉第电流效率不是 $100\%$，电流效率常数也应计入系数。式中 $K_{1}$ 为实际被采样的比例系数（产物在压差下进入真空系统的同时也会向溶液本体扩散），$K_{2}$ 概括「进入离子源」与「被电离进入质量分析器」两级比例。\tocheck{未对照原页面图像：源文本在 $K_{1}$ 与 $K_{2}$ 之间的定义句有脱字（OCR 作「可将这两项系数统一考虑，定义由此可得到」），$K_{2}$ 的严格定义需对照原书 p.874}
 由于影响 $K$ 值的因素众多（电化学实验环境、质谱仪参数、物种自身性质），不同物种在不同环境下的 $K$ 值可能都不相同，因此\textbf{特定物种的 $K$ 值必须在电化学条件下定量生成该物种、由质谱检测其离子电流来标定}，不能沿用文献值\src{S2 §11.6.3，书内 p.874 / PDF p.362}。
\end{principle}

对非挥发性产物的在线检测，$1986$ 年 Hambitzer 和 Heitbaum 把 N,\,N-二甲基苯胺电解氧化，并将电极表面的电解液连续导入质谱热喷雾离子室的前置进样毛细管中，实现了在线检测。随着质谱技术发展，电化学质谱联用结合了更多样品电离方式以分析不同形态与性质的产物，包括电喷雾电离、大气压化学电离、大气压光致电离、快原子轰击电离与电感耦合等离子体电离等。另一种可与更常用电极材料匹配的设计，是把 Teflon 膜接口放在距旋转圆柱电极很近处（约 $0.3$ mm）\src{S2 §11.6.3，书内 p.874—875 / PDF p.362—363}。

二次离子质谱（SIMS）是表征表面与薄膜的非原位 UHV 方法，可与电化学联用：用高能一次离子束（如 $15$ keV \ce{Cs+}）轰击表面，通过溅射侵蚀表面、由表面组成中产生二次离子并用质谱检测。利用一次离子束扫描可作二维表征，监测单离子强度与溅射时间的关系可得深度分布。SIMS 的测量极限比 XPS 或 AES 高得多（\num{1e-8}～\num{1e-4} 原子分数）。但 SIMS 并不是一种「真实表面」技术：二次离子的产生效率由一次离子束所产生的薄的离子注入层的三维性质决定，深度分布的假象可由此产生。把电极转移到 UHV 样品腔后，研究电极表面的其他可用方法是热脱附与激光脱附\src{S2 §11.6.3，书内 p.875 / PDF p.363}。

\begin{crosscheck}
 原位质谱与谱学方法的分工很清楚：红外／拉曼给\textbf{表面物种与成键}，质谱给\textbf{产物种类与生成速率}。二者对同一电位点必须给出一致的速率控制图像——若 DEMS 显示某产物在 $0.5$ V 起快速生成，而原位红外显示中间体覆盖度在同一电位开始下降，则两者共同支持「中间体在该电位被消耗」的机理。反过来，如果只报 DEMS 的产物分布而不报 $K$ 的标定方式，产物比例（常被用来判断 2 电子与 4 电子途径的比例）就只是相对信号，不能作为机理证据。
\end{crosscheck}

\subsubsection{荧光单分子与单纳米粒子催化}

纳米材料与催化领域的发展，需要能在微观层面解析纳米催化剂催化过程的新技术。荧光单分子／单纳米粒子技术基于荧光显微技术，利用现代单光子检测器的高灵敏度，实现单个纳米粒子表面荧光产物分子一个一个的观测，从而从微观反应动力学与热力学角度解析催化过程\src{S2 §11.6.4，书内 p.875—876 / PDF p.363—364}。

\begin{principle}
 该技术跟踪单个分子的荧光信号：荧光相关的前驱体分子在被催化发生氧化还原反应的过程中，在荧光态与非荧光态之间转化，通过跟踪荧光信号的强弱获取反应动力学信息。要拿到单分子与单纳米粒子的信号，需要同时满足三个条件\src{S2 §11.6.4，书内 p.877 / PDF p.365}：
 （1）\textbf{浓度极低}——催化剂浓度与反应物浓度都要极低，确保每次捕获到的是单次反应引起的荧光信号变化；
 （2）\textbf{界面全反射激发}——利用激光在界面处发生全反射时产生的渐逝波（倏逝波）激发界面处的荧光分子，避免背景溶液中的荧光分子被激发而抬高背景；
 （3）\textbf{超快单光子检测}——用现代单光子检测器捕捉信号，在获得良好信噪比之后，对大量荧光信号作统计分析，得到产物形成与脱附过程的信息。
 具体做法是把纳米粒子稀疏分散在石英载玻片或类似基底上，反应物流过表面时在纳米粒子表面活性位上生成可被激光激发荧光的产物分子。产物生成之前看不到该纳米粒子；催化反应一旦完成，生成的荧光分子立刻被激发发光，因而可以通过单分子荧光显微镜直接观察检测，并通过定位该荧光分子精确确定该纳米粒子的位置。用单光子检测器可实时记录单个产物分子\textbf{形成前的等待时间}与\textbf{形成后到脱附前在表面的停留时间}——这样就把常规方法中的「一个催化周期」拆成两个更细致的过程：产物分子的生成，以及产物的脱附与活性位的释放，从而分别考察单个纳米粒子表面的产物形成速率与脱附速率\src{S2 §11.6.4，书内 p.877 / PDF p.365}。
\end{principle}

最典型的系统研究出现在 $2008$ 年：陈鹏等以金纳米粒子为催化模型，首次用荧光单分子催化技术在单分子水平上研究了单个金纳米粒子催化羟胺与刃天青（resazurin）反应生成荧光产物试卤灵（resorufin）的过程，并首次提出适用于单分子单纳米粒子催化的动力学机理模型\src{S2 §11.6.4，图 11-51，书内 p.876—877 / PDF p.364—365}。该方法的独特价值在于能一个一个地研究看似一样的纳米粒子，鉴别粒子之间的活性差别（反应快慢、产物吸附强度或脱附速率），而这些信息是常规方法无法获得的；除了用统计方法得到大量粒子的平均信息，还能得出个别粒子的「个性」信息，有助于更精确地理解催化剂特性与其物理特性（粒径、组成、晶面类型）之间的相关性\src{S2 §11.6.4，书内 p.878 / PDF p.366}。

在超分辨成像方向，$2006$ 年 Hofkens 等用荧光显微镜实时观察催化反应过程中单个有机分子的化学转化，绘制出催化剂晶体的活性位点分布，并发现不同催化晶面具有催化不同化学反应的能力；$2009$ 年该组进一步把超分辨成像精度提高到可与电子显微镜匹敌；$2013$ 年 Majima 课题组用自主合成的探针对 \ce{TiO2} 光解水过程的活性位点作超分辨成像，发现可见光引起的还原反应发生在金纳米粒子附近几十纳米范围内，在微观层面证实了等离子引发的电子空穴分离机理；$2014$ 年方宁课题组用荧光超分辨成像揭示单个金—硫化镉纳米粒子表面的光催化性质，揭示了两种完全不同的、能量依赖的电荷分离机理；$2015$ 年 Hofkens 等对石油裂解过程产物分子作超分辨成像，发现分子筛催化石油裂解过程中产物主要在布朗斯特酸位点上形成，并揭示了不同粒子间的差异性。超分辨成像不仅能描绘静态活性位点，还能跟踪其位置变化，例如中国科学院长春应用化学研究所徐维林课题组 $2015$ 年对纳米金三角等催化剂纳米粒子作动态超分辨成像，发现金纳米粒子表面活性位点因重构而发生移动\src{S2 §11.6.4，书内 p.878 / PDF p.366}。

在动力学方向，$2008$ 年陈鹏等利用催化剂在荧光态与非荧光态之间转化的过程，实现对微观催化反应的动力学分析并揭示了三种不同的催化机理；$2014$ 年徐维林课题组首次结合量化计算与「用空心纳米笼困住反应分子」的方法，实现化学反应过程中各基元步骤的分子反应动力学分析；随后该组把荧光单分子催化技术拓展到电化学领域，揭示了电催化剂的衰减机理；同年该课题组通过对不同粒径纳米立方钯催化剂催化活性的解析，获得纳米粒子「边—面」活性差异，把单粒子催化技术拓展到亚单粒子水平。该技术虽刚发展起来，其应用正朝着更深更广的方向发展，已从催化动力学研究扩展到热力学研究\src{S2 §11.6.4，书内 p.879 / PDF p.367}。

\subsubsection{原位谱学的共同边界}

\begin{table}[htbp]
 \centering
 \caption{四类原位谱学方法的分工与边界。「关键条件」一列是使用该方法数据前必须确认的前提；「主要陷阱」一列给出最常见的误用。\src{据 S2 §11.6.1—§11.6.4 汇总，书内 p.867—879 / PDF p.355—367。}}
 \label{tab:ch11c-insitu-methods}
 \begin{tabularx}{\onecolwidth}{@{}p{20mm}p{29mm}p{30mm}X@{}}
  \toprule
  方法 & 拿到什么 & 关键条件 & 主要陷阱 \\
  \midrule
  原位红外（EMIRS／SNIFTIRS／IRRAS） & 吸附物种身份与成键、双电层物种随电位的变化 & 窗口材料透明且不溶；薄层 $1\sim100$ nm & 得到的是两个电位的差谱（EMIRS／SNIFTIRS），不能直接读绝对覆盖度；IRRAS 只代表固定电位 \\
  原位拉曼（含 SERS） & 振动指纹：金属—吸附物与分子内振动 & 薄层 $0.1\sim1$ mm；SERS 需粗糙／胶态贵金属基底 & 增强因子随基底与波长变化，峰强不可跨批比较；薄层改变传质 \\
  原位质谱（EMS／DEMS） & 挥发性产物的种类与生成速率；可作同位素标记 & 膜接口与压差；响应时间 $<1$ s（DEMS） & 系数 $K$ 必须原位标定；只测挥发性物种，非挥发性产物需另接电离方式 \\
  SIMS 与脱附法（非原位 UHV） & 表面组成、深度分布、脱附物种 & 需把电极转移进 UHV & SIMS 由注入层三维性质决定产额，会产生深度分布假象；不是真实表面技术 \\
  荧光单分子／单粒子 & 单粒子层面的生成与脱附动力学、活性位点空间分布 & 极低浓度；全反射倏逝波激发；单光子检测 & 分析对象是荧光产物分子，只适用于能生成荧光产物的反应体系 \\
  \bottomrule
 \end{tabularx}
\end{table}

\begin{pitfall}
 原位谱学的共同陷阱不是「分辨率不够」，而是\textbf{为了能测而改变了被测对象}：红外要求纳米级薄层、拉曼要求亚毫米级薄层、质谱要求真空接口、荧光要求极低浓度，每一项都在改变传质与表面状态。因此原位谱学给出的电位—覆盖度或电位—速率关系，只在同样的薄层与浓度条件下成立；把它们直接外推到厚层、流动或高电流工况（这正是燃料电池的真实工作状态），是这一类方法最容易被忽略的边界\src{S2 §11.6.2，书内 p.873 / PDF p.361}。
\end{pitfall}

\begin{crosscheck}
 原位谱学的结论应与三类独立信息互查：\textbf{（1）电化学响应}——谱带出现／消失的电位应与循环伏安上的电流峰或阶跃位置一致（如甲醇氧化中 $0.5$ V 的谱带减弱与电流增大同步）；\textbf{（2）非原位表征}——把同一电极取出做 XPS、SIMS 或脱附实验，检查表面物种与原位结论是否一致，并注意 UHV 转移过程中的表面变化；\textbf{（3）同位素标记}——用 \ce{^{13}C} 或 \ce{D} 取代确认谱带归属与产物来源，这是把「谱带—物种」对应关系从文献搬运变成自证的关键一步。
\end{crosscheck}

\srcnote{本分片覆盖 S2（下册）第 11 章 §11.5 后半（阻抗谱的应用与读图，书内 p.866 / PDF p.354）与 §11.6 全节（原位谱学技术及应用，书内 p.866—879 / PDF p.354—367）。§11.5.1 的阻抗基本定义、等效电路推导与 Nyquist／Bode 图定义不在本分片源文本范围内——它们写在本章另一分片（§11.5.1，依源文本写成、逐条标有 \src{} 依据），故本分片不再重复，只保留源文本 §11.5.2 及以后特有的内容：读图实例（甲醇氧化）、参数相关性与误判、测量条件与互补回查、以及 §11.6 全节；其中图 11-42／11-43 的坐标与相角取值、以及 DEMS 系数 $K_{1}$、$K_{2}$ 的定义按学科标准形式重建，存疑点已在正文用 \tocheck{} 就地标出。图 11-42、11-43、11-44～11-51 未对照原页面图像，本分片不再自绘阻抗模型图（原型计算的 Nyquist 模型图见上一节），未裁切原书谱图。}
}{}
\IfFileExists{chapters/parts/ch11-e.tex}{% ============================================================================
% ch11-e.tex —— §11.5.2 阻抗技术在电催化领域中的应用（前半）
%               现场/外场 EIS 的取舍、参比与工作电极制备、两种测量模式
% 源：sources/parts/ch11-b.txt 行 1095 起；书内 p.862—865 / PDF p.350—353
%
% 【节序说明 —— 装配者必读】
%   本文件在 chapters/ch11.tex 的 \input 顺序中排在最后（a→b→d→c→e），
%   但按源书节序，§11.5.2 应排在 §11.5.1（对应 part d）之后、
%   §11.6（对应 part c 后半）之前。即正确阅读顺序是 a→b→d→e→c，
%   而不是装配文件当前的 a→b→d→c→e。
%   本文件**不**修改 chapters/ch11.tex，节序调整由装配者统一处理。
% ============================================================================

\subsection{阻抗技术在电催化领域中的应用}

\src{书内 p.862 / PDF p.350}

前面讨论的都是同一个电极过程内部的等效元件。一旦把 EIS 用到一个真实的电催化
体系上，问题就分裂成两个层面：\emph{在哪里测}（现场还是外场）与\emph{怎么测}
（何种电极、何种扰动模式）。这两个层面的选择往往比后续的等效电路拟合更能
决定数据是否可用，因此本节先讨论它们，再进入典型谱图的判读。

\subsubsection{现场与外场阻抗谱}

\begin{principle}
\textbf{按测量对象划界的两种 EIS}。取决于研究目的，EIS 可以分为现场的
（in-situ）与外场的（ex-situ）。这一划分不是按仪器或频段，而是按\emph{被测
体系在测量时是否处于其真实工作状态}：现场 EIS 把单池或电堆整体作为被测
对象，在电池实际运行、带直流负载的条件下叠加交流扰动；外场 EIS 则把材料或
组件（催化剂、催化层、膜、GDE）取出，放进一个可控的电解池里单独表征。分界
线一旦划定，后面所有的取舍——装置成本、可获得的参数、数据的可解释性——
都由它派生出来。
\end{principle}

\begin{measure}
\textbf{现场 EIS：装置、代价与人为因素}。现场 EIS 经常用于单池或电堆的分析。
它的测量装置通常是复杂、昂贵且耗时的。更棘手的是运行期间——尤其在高直流
负载时——由于电缆和连线等的干扰，可能引入人为因素，使谱的高频段不再反映
电池本身。这里的因果关系值得说清楚：高直流负载意味着大电流，大电流在电缆
与连线上产生可观的电压降与感抗，而感抗项随频率上升而增大，于是高频端出现的
那个“电感尾巴”究竟来自电池还是来自测量回路，无法从谱形本身分辨。这也是
现场 EIS 最典型的失效方式。
\end{measure}

\begin{measure}
\textbf{外场 EIS：三电极体系与电解质选择}。正因为上述代价，外场 EIS 一般是
很多燃料电池研究者的首选。外场 EIS 一般采用传统的三电极体系在硫酸电解质中
进行，参比电极一般为饱和甘汞电极（SCE）、标准氢电极（NHE）、\ce{Ag/AgCl}
或 \ce{Hg/HgO} 电极，有时也用动态氢电极（DHE）作参比电极；工作电极一般通过
将薄层催化剂铸在干净的玻碳电极或热解石墨电极上，或者将电极直接浸泡在含有
欲测催化剂的溶液中，有时，气体扩散电极（GDE）也可以用作工作电极。
\src{书内 p.863 / PDF p.351}
\end{measure}

\begin{probout}
\textbf{现场与外场 EIS 的取舍判据}——不是优劣排序，而是“什么目的配什么代价”：

\begin{itemize}
  \item \textbf{要问“这个电池/电堆现在出了什么问题”}$\rightarrow$ 选现场 EIS。
        只有现场测量才能保留膜、气体扩散电极、MEA、单池各元件在真实工作
        条件下的贡献，也才可能定位问题元件。\emph{代价}：装置复杂、昂贵、
        耗时；高直流负载下电缆与连线引入人为因素，高频段不可信。
  \item \textbf{要问“这个催化剂/材料本身好不好”}$\rightarrow$ 选外场 EIS。
        三电极体系把工作电极电位测准，排除了对电极与欧姆降的纠缠，数据
        可解释性最高。\emph{代价}：离开真实工况，电解质、温度、气体分压、
        水管理都与单池不同，测得的活性排序未必能外推到电堆。
  \item \textbf{两者都要}时，正确用法是外场 EIS 建立机理与参数基准，现场 EIS
        检验该基准在工况下是否仍然成立，而不是把两者的数值直接相减当作
        “工况劣化量”。
\end{itemize}
\end{probout}

\subsubsection{电极制备与测量模式}

\begin{dataflow}
\textbf{从电极制备到一条可信的阻抗谱}。参数的可信度在电流进入恒电位仪之前
就已经被决定了，链条如下：

\begin{enumerate}
  \item \textbf{参比电极选型}（SCE / NHE / \ce{Ag/AgCl} / \ce{Hg/HgO} / DHE）
        $\rightarrow$ 决定电位标度与液接电位误差；
  \item \textbf{工作电极制备}（薄层催化剂铸于干净玻碳或热解石墨，或电极直接
        浸泡于含催化剂的溶液，或用 GDE）$\rightarrow$ 决定催化剂的真实负载量、
        分散状态与接触电阻；
  \item \textbf{扰动模式选择}（恒电位 / 恒电流）与交流振幅、频率范围
        $\rightarrow$ 决定谱是否在仪器的线性区内被测得；
  \item \textbf{等效电路拟合} $\rightarrow$ 输出 $R_\mathrm{s}$、
        $C_\mathrm{dl}$、$R_\mathrm{ct}$ 等参数。
\end{enumerate}

这条链条是\emph{串联}的：第 1、2 环节引入的误差无法在第 4 环节通过拟合消除，
只会在拟合结果里伪装成某个元件参数。
\end{dataflow}

\begin{pitfall}
\textbf{为什么参比与工作电极制备决定数据可信度}。三点具体的因果：

\begin{itemize}
  \item \textbf{参比电极决定“电位”这个自变量的真值。} 不同参比（SCE、NHE、
        \ce{Ag/AgCl}、\ce{Hg/HgO}、DHE）各自的电极电位不同，报告数据时必须
        写明参比与电解质；只写“在 0.5 V 下测得”而不写参比，等于没有给出
        电位。DHE 之所以称为“动态”氢电极，正是因为它的电位随体系中氢分压
        与酸度而移动，用它可以减小与实际工况的偏离，但也就失去了固定标度的
        便利。
  \item \textbf{工作电极制备决定“有多少催化剂在工作”。} 薄层催化剂铸在玻碳
        或热解石墨上，要求基底干净：基底上的残留物既贡献额外的界面电容，
        又可能自身有电化学活性。而“将电极直接浸泡在含有欲测催化剂的溶液中”
        这一做法，催化剂是悬浮或溶解状态，其负载量与附着状态无法定量，得到
        的阻抗只能作相对比较。GDE 则引入了额外的气体扩散层与三相界面，谱中
        会多出传质相关的特征。
  \item \textbf{两种制备方式的差异会在谱上表现为同一个 $R_\mathrm{ct}$ 的不同
        数值}，而谱形本身无法告诉你究竟是本征活性不同，还是接触与负载不同。
        这正是必须在论文中完整描述电极制备细节的原因。
\end{itemize}
\end{pitfall}

\begin{measure}
\textbf{两种测量模式及其各自的代价}。阻抗测量一般有两种模式，一种是恒电位
模式，另一种是恒电流模式。从这两种模式中得到的结果基本上没有明显不同。然而，
取决于测量对象，这两种模式各有优缺点。\src{书内 p.862 / PDF p.350}

\emph{恒电位模式}：扰动是交流电位，响应是通过研究体系的测量电流。由于一个
小的电位都可能导致大的电流变化，从而导致恒电位仪的过载，因此，尝试不同的
交流电压确保恒电位仪不会过载是非常必要的。在燃料电池的测量中，交流振幅一般
选择在 5$\sim$15 mV，扫描频率范围一般在 1 mHz$\sim$100 kHz 之间。在催化剂
活性的测量中，控制电位能更好地控制电化学反应的发生，因此一般采用恒电位模式。

\emph{恒电流模式}：扰动是交流电流，响应是通过研究体系的测量电位。由于交流
和直流大小很容易被运算放大器所确定，能有效防止电池和恒电位仪过载，因此能
采用相对大的交流电，从而得到更多的阻抗信息。对于电池阻抗测量，一般电流比较
大，因此更适合用恒电流模式测量，恒电位模式容易导致过载。
\src{书内 p.862 / PDF p.350}
\end{measure}

\begin{center}
\small
% 注意：本项目 calc 包会改写 \dimexpr，任何 `p{\dimexpr...\onecolwidth\relax}`
% 形式的尺寸运算都会报 "Illegal unit of measure"（calc Error: `.' invalid）。
% 故首列改用固定 26 mm（版心 165 mm），其余由 tabularx 按 \onecolwidth 分配，
% 全程不做尺寸运算。\onecolwidth = 0.80\textwidth ≈ 132 mm。
\begin{tabularx}{\onecolwidth}{@{}>{\raggedright\arraybackslash}p{26mm}
  >{\raggedright\arraybackslash}X >{\raggedright\arraybackslash}X@{}}
\toprule
 & 恒电位模式 & 恒电流模式 \\
\midrule
扰动量 & 交流电位 & 交流电流 \\
响应量 & 测量电流 & 测量电位 \\
主要风险 & 小电位变化即可引起大电流，恒电位仪过载 &
            交流与直流幅值不易独立设定 \\
适用对象 & 催化剂活性测量（控制电位即控制反应） &
            电池与电堆阻抗（电流大，恒电位易过载） \\
常用参数 & 燃料电池中振幅 5$\sim$15 mV；频率 1 mHz$\sim$100 kHz &
            可采用相对大的交流电，阻抗信息更多 \\
\bottomrule
\end{tabularx}
\end{center}

\begin{pitfall}
\textbf{四类会把谱读错的误判}。

\begin{enumerate}
  \item \textbf{把接触电阻当电荷传递电阻。} 催化剂层与玻碳或热解石墨基底之间
        的接触电阻、以及电缆与接头引入的欧姆项，都会在 Nyquist 图上表现为
        高频端的实轴截距或一个额外的小半圆。若不先做空白基底测量作对照，
        很容易把这个本不属于电极反应的电阻读成 $R_\mathrm{ct}$，从而低估
        催化活性。判据是：$R_\mathrm{ct}$ 随电位（过电位）应有系统性变化，
        接触电阻不会。
  \item \textbf{恒电位仪过载导致谱失真。} 恒电位模式下，小的电位扰动即可引起
        大电流，一旦超出恒电位仪的电流量程，仪器不再维持设定的交流电位，谱形
        被削平或出现伪感抗，而软件仍会照常输出数据点。因此“尝试不同的交流
        电压确保恒电位仪不会过载是非常必要的”——这句话不是操作建议，而是
        数据有效性的前提。
  \item \textbf{外场 EIS 的电解质与实际工况不符。} 外场 EIS 常在硫酸电解质中
        用三电极体系进行，而真实燃料电池是固体聚合物电解质、含液态水与气体
        扩散层。电解质种类、离子强度、pH 与传质环境都不同，因此在硫酸中测得
        的 $R_\mathrm{ct}$ 排序不能直接当作电堆中的性能排序。
  \item \textbf{现场 EIS 的电缆干扰。} 高直流负载下，大电流经电缆与连线产生的
        干扰会把测量回路自身的阻抗叠加到谱上，通常表现为高频区的伪感抗或
        不可解释的相角。判据是：更换电缆走线或缩短线长后高频段发生变化，
        则该段特征是人为因素而非电池特征。
\end{enumerate}
\end{pitfall}

\begin{crosscheck}
\textbf{读图实例：原书图 11-41 的 Nyquist 图}。该图为溶液电阻不可忽略时电化学
极化与浓差极化同时存在电极的等效电路及 Nyquist 图。
\src{书内 p.862 / PDF p.350}

先说几何，再说归属。对于电化学极化和浓差极化同时存在的电极，该电极过程由
电荷传递过程和扩散过程共同控制。
\end{crosscheck}

\begin{obs}
谱形由两段构成：

\begin{itemize}
  \item 高频区：一个半圆，其半径为 $R_\mathrm{ct}/2$；
  \item 低频区：一条 45° 的直线。
\end{itemize}
\end{obs}

\begin{infer}
由此可读出控制步骤的归属：

\begin{itemize}
  \item 半圆表示电荷传递过程为控制步骤——半圆的直径给出 $R_\mathrm{ct}$；
  \item 45° 直线表示由电极反应的反应物或产物的扩散控制；
  \item 因此该电极过程由电荷传递过程和扩散过程\emph{共同}控制。
\end{itemize}
\end{infer}

\begin{crosscheck}
\textbf{观察与推断的界线必须守住}。图谱本身只给出“半圆 + 45° 直线”这一几何
事实；半圆 $\rightarrow$ 电荷传递控制、直线 $\rightarrow$ 扩散控制，是建立在
等效电路模型之上的解释。若换了模型（例如存在两个时间常数、半圆被压扁、或
直线偏离 45°），同一段几何形状可以有别的解释。
\end{crosscheck}

\begin{pitfall}
读原书图 11-41 时最容易犯的两个错：一是把高频端半圆在实轴上的\emph{截距}
当成 $R_\mathrm{ct}$——只有当溶液电阻可忽略时截距才近似等于
$R_\mathrm{ct}$，而此图的前提恰恰是“溶液电阻不可忽略”，所以必须先扣除溶液
电阻 $R_\mathrm{s}$，再取半圆直径；二是把低频 45° 直线的起点当作扩散开始
出现的位置精确定值，实际上直线段与半圆的过渡是渐变的，起点的判定依赖拟合
区间的人为选择。
\end{pitfall}

此处已对照原页面图像（书内 p.862 / PDF p.350）：原书图~11-41 确为高频区一个
半径 $R_\mathrm{ct}/2$ 的半圆加低频区一条 $45^\circ$ 的直线，图中另以竖直虚线标出
“电荷传递控制”与“扩散控制”的分界；等效电路元件印作 $R_\Omega$、$R_\mathrm{ct}$、
$C_\mathrm{dl}$ 与方框化的 $\tilde Z_\mathrm{W}$，半圆一侧的实轴标出 $R_\Omega$、
$R_\Omega+R_\mathrm{ct}/2$、$R_\Omega+R_\mathrm{ct}$。与本节的谱形叙述一致。

此处已对照原页面图像（书内 p.861--862 / PDF p.349--350）：图~11-40（p.861）与
图~11-41（p.862）的图题\textbf{印刷完全相同}，均为“溶液电阻不可忽略电化学极化
电极的等效电路及 Nyquist 图”，属原书自身图题重复；正文分别以“其等效电路及
Nyquist 如图 11-40 所示”“……如图 11-41 所示”引用两图，本节按正文分工处理无误。

此处已对照原页面图像（书内 p.862 / PDF p.350）：原书正文印“交流振幅一般选择在
5$\sim$15 mV，扫描频率范围一般在 1 mHz$\sim$100 kHz 之间”，数值上下限与本处一致
（原书用全角波浪号“～”，此处用 $\sim$，属排版差异）。

此处已对照原页面图像（书内 p.861--862 / PDF p.349--350）：原书印“参比电极一般为
饱和甘汞电极（SCE）、标准氢电极（NHE）、Ag/AgCl 或 Hg/HgO 电极，有时也用动态氢
电极（DHE）作参比电极”，\ce{Hg/HgO} 的写法与“动态氢电极（DHE）”的中文名、缩写
均无误。

\src{书内 p.862--863 / PDF p.350--351}
}{}
